% !TEX root = ../print.tex
% Chapter 5, first part: the price of the plane. Draft sections 5.1-5.3, lines 743-866 of
% draft/spatial-hemigroup-pyramids.md; 06-helmholtz.tex continues the chapter with the draft's
% sections 5.4-5.6. The "Checks for section 5" (lines 1029-1044) are records, not transcribed.
%
% Moves (adr/0004, decision 3). Lemma 5.3 (lem:lattice-infinitely-divisible) moved OUT to chapter 3
% and is cited behind. Two pieces moved IN from chapter 4: the paragraph after Proposition 4.7 (draft
% line 646), which says which masks Proposition 4.7 covers, is the node prop:creation-grid-masks
% after Proposition 5.6; the grid case of Theorem 4.10(4) (draft lines 711 and 721) is the node
% prop:grid-frame-step-not-selected after Proposition 5.4.
%
% Nodes: def:grid-family (new label, the draft's paragraph after Convention 5.1), def:exact-senses,
% def:exact-grid-pyramid, prop:grid-pyramid-not-isotropic [A], prop:grid-frame-step-not-selected,
% prop:exact-masks (regraded [T] in wave 3g), prop:resolvent-masks-exact, def:level-records (new label),
% lem:records-creation (new label), prop:creation-grid-masks (new label),
% prop:binomial-records (new label, the draft's open binomial records, closed 2026-11-12),
% prop:line-and-arrays [A], prop:binomial-knot-relation (new label, Remark 5.8(1)).
% Review round 2, batch C (2026-11-23): prop:expanding-grid-pyramid-not-isotropic [A], not ready
% (decision D1), and prop:no-isotropic-refinable-finite-mask [T].
% Review round 3, batch C (2026-11-23): prop:grid-pyramid-singular-kernels [A], after the remark on
% the line that follows prop:expanding-grid-pyramid-not-isotropic, not ready for its clause (3);
% prop:band-limited-refinable-mask [A], after prop:no-isotropic-refinable-finite-mask, not ready
% (decision D2).
%
% What stays prose, and why (one line each):
% - the chapter's lead and the section leads: they say what the sections prove;
% - the comment after Convention 5.1 (the senses answer different questions): explanation;
% - the comment after the grid family (the sampled Gaussian): the reason for the restriction;
% - the comment after Definition 5.2 (relation to the isotropic axioms): it names where each
%   relation is proved;
% - the paragraph before Proposition 5.4 (the jumps of a law): the picture behind the proof;
% - one remark on the line per section (after Propositions 5.4 and 5.5), which absorbs Remark
%   5.8(3) and the draft's sentence on [V-line v0.1, Thm. 5.2] inside the proof of 5.4(1);
% - the comments after Propositions 5.4, 5.6, 5.7 and the two moved-in nodes: each claim in them
%   is a reference to a node;
% - Remark 5.8(2) and (4): each sentence cites a node; Remark 5.8(1)'s mathematics is the node
%   prop:binomial-knot-relation, its closing question is marked "not addressed here".
%
% The blind reviews of this file (records/pyramids/REVIEW-blind-05-price.md and
% REVIEW-blind2-05-price.md) are answered at the anchors; their "Answers" sections say how.
%
% The split (adr/0007, 2026-11-16). This file ends the theory paper. Its last section,
% sec:discrete-helmholtz, moved here verbatim from the start of 06-helmholtz.tex (the draft's
% section 5.4 up to Proposition 5.10(2)): def:grid-family, prop:resolvent-masks-exact and
% prop:creation-grid-masks state results for those kernels. The chapter lead's last paragraph
% still describes the draft's sections 5.4-5.6, of which only that first part stays in this
% chapter: the writing session rewords it.

\chapter{The price of the plane}
\label{ch:price}

% draft: the lead of section 5, lines 745-749.
Image pyramids live in the plane. A grid pyramid applies one \emph{mask}, a symmetric probability
law $w$ on the lattice whose weights are the smoothing step of one level. The steps it applies from
octave to octave, $w$ dilated to each coarser grid, form its \emph{ladder}
(Definition~\ref{def:exact-grid-pyramid}). Call the pyramid exact when a continuous covariant
family runs through that ladder. On a grid in the plane, exactness in this sense and rotation
invariance exclude each other. This chapter proves it. It then measures, for the two exact
pyramids that the next chapter computes, how far from isotropic each is and what each pays in
aliasing.

The reason is the carrier of the \Levy{} measure. An exact pyramid's mask is infinitely divisible.
An infinitely divisible law carried by a lattice has its \Levy{} measure on the lattice. The
covariant kernel inherits a \Levy{} measure on countably many points. A nonzero rotation-invariant
\Levy{} measure in $d \ge 2$, finite away from the origin, cannot sit on countably many points. On
the line the orthogonal group is $\Orth(1) = \{\pm 1\}$, and the last step is empty.

Section~\ref{sec:price-senses} fixes what \emph{exact} means, since on a
grid the word has three senses. Section~\ref{sec:price-obstruction} defines the exact grid pyramid
and proves the obstruction, with the frame \Matern{} pyramid as an exact family that meets it.
Section~\ref{sec:price-masks} says which masks have an exact pyramid, among them the resolvent
masks of the computations, and what the pyramid algorithm returns. The rest of the chapter gives
those resolvents as grid families: the discrete Helmholtz family, an exact cascade on one grid,
and two stencils besides the five-point one. It then treats the other reading of \emph{exact}, a
family whose subsample is again a member, and shows that the discrete Helmholtz family fails it at
every index in the plane. Last, it measures the price for the frame \Matern{} pyramid and the
five-point resolvent pyramid. The price is a function of the ratio $\rho = h/(ct)$ of the grid
spacing $h$ to the range $ct$ of the smoothing at scale $t$, with $c$ the range constant of the
\Matern{} pyramid (Theorem~\ref{thm:matern-pyramid}).

Throughout, $\Gamma \subset \RR^d$ is a full-rank lattice: $h\ZZ^d$, or the hexagonal lattice in
the plane. For $c > 0$, $D_c\mu$ is the image of a measure $\mu$ under $x \mapsto cx$, as in
Definition~\ref{def:family-dimension-d}, so that $\widehat{D_c\mu}(\omega) = \widehat{\mu}(c\omega)$.
For $R \in \Orth(d)$, $D_R\mu$ is its image under $x \mapsto Rx$.

\section{Three senses of exact}
\label{sec:price-senses}

A negative statement about exactness is only as strong as the sense it uses, so the senses come
first. Sense (ii) needs one more word, the grid family, which comes first.

% draft: the paragraph after Convention 5.1 (section 5.1), line 757, its first two sentences.
% CHANGED: made a definition node of its own (new label def:grid-family), since sense (ii) of
% def:exact-senses is phrased in terms of it; the example that follows it in the draft is the
% prose after the node.
% CHANGED: "defined by grid operations" made a closed list (blind review, flag 9): the stencils
% are the two the paper treats, and "the resolvent of a stencil and its convolution powers" is
% read with every index gamma; this narrows the draft's open phrase.
% CHANGED: the bare Delta of the resolvent equation is named, Delta_h or Delta_9, since Delta is
% the continuum Laplacian in the neighbouring chapters (second blind review, flag 6). The closing
% example (the sampled Laplace family) moved to the prose after the node (flag 11).
\begin{definition}[grid family]
\label{def:grid-family}
\lean{Pyramids.IsStencilResolvent, Pyramids.stencilFamily, Pyramids.maskPowers}\leanok
\uses{prop:discrete-helmholtz, prop:nine-point-family}
A \emph{grid family} on a lattice is one of the following families of kernels on that lattice.
\begin{enumerate}
\item For the $(2d+1)$-point Laplacian $\Delta_h$ on $h\ZZ^d$
(Proposition~\ref{prop:discrete-helmholtz}): the convolution powers $\kappa_t^{*\gamma}$, $t > 0$,
$\gamma \ge 1$ an integer, of its resolvents $\kappa_t$, the summable solutions of
$(1 - t^2\Delta_h)\kappa_t = \delta_0$. For the nine-point Laplacian $\Delta_9$ on $h\ZZ^2$
(Proposition~\ref{prop:nine-point-family}): the same, with $\Delta_9$ in place of $\Delta_h$.
\item For a mask $w$ on the lattice: its convolution powers $w^{*n}$, $n \ge 1$.
\end{enumerate}
% CHANGED (fidelity review, 2026-11-15, R11): the status line says that the closed list is a
% convention the declarations do not type.
\statusT\quad\emph{(A definition. \texttt{stencilFamily L} takes an arbitrary operator $L$ and
\texttt{maskPowers} an arbitrary array; that the operators are $\Delta_h$ and $\Delta_9$ only, and that
a family of samples of a continuum family is no grid family, is a convention of the text, read by
the statements, not typed.)}
\end{definition}

A family defined as the samples of continuum kernels is a grid family only where it equals one of
these. The sampled Laplace family on $h\ZZ$ is one, since it equals the family of the three-point
Laplacian (Proposition~\ref{prop:recursion-helmholtz}(2)). The restriction matters for sense (ii)
below. The samples of an isotropic Gaussian on $h\ZZ^2$, subsampled by $q$, are the samples of the
same Gaussian on $qh\ZZ^2$. Read literally, they would be exact in sense (ii), and they are not a
grid family.

% draft: Convention 5.1 (section 5.1), line 755.
% CHANGED: the convention's two sentences of rationale ("The sense is deliberately weak ..." and
% "They are different questions: ...") moved to the prose after the node, which now holds only
% the three definitions and the convention on sense (iii); "this section" became "this chapter";
% long sentences split at their semicolons, with no change of content.
% CHANGED: "C = sum_m lambda(qm) < 1" now says when the inequality is strict (blind review,
% flag 23): the draft asserted it unconditionally.
% CHANGED: in (ii) the ratio q is bound (an integer q >= 2), and "of some range and some index"
% is said for item (1) of def:grid-family only; for item (2) the coarse member is a power of the
% dilated mask (second blind review, flag 5).
\begin{definition}[three senses of \emph{exact}]
\label{def:exact-senses}
\lean{Pyramids.IsExactLawArray, Pyramids.IsExactSampleArray, Pyramids.latticePt, Pyramids.IsExactSubsampleIn, Pyramids.IsExactLatticePyramid}\leanok
% CHANGED (the split, adr/0007, 2026-11-16): prop:helmholtz-subsampling leaves the uses, and its
% parenthetical citation leaves item (2). It named where sense (ii) is applied, which is the
% computation paper; nothing in the definition rests on it. Meaning unchanged.
% CHANGED (the theory paper's blueprint round, 2026-11-16, review of Section 6, flags COUNT and
% ENVELOPE): "three senses in this paper" became "three senses for a pyramid or a grid
% computation", since "exact cascade" is a fourth use, now named in the prose after the node. The
% usage notes left the environment for that prose: the floating-point reading and aliasing from (i),
% the value of C and the impulse-response reading from (ii), and the statements that use (iii).
% prop:sampled-laplace-exact leaves the uses with the aliasing sentence. The three definitions are
% unchanged.
\uses{def:grid-family, def:exact-grid-pyramid}
The word \emph{exact} has three senses for a pyramid or a grid computation, each fixed by what it
is attached to.
\begin{enumerate}
% CHANGED (skeleton review 05, item 1): "in exact arithmetic" names the reading the typed predicate
% has; floating point is not modelled, and "up to rounding" is kept for a run of the algorithm.
\item A \emph{grid computation} is exact when the arrays it produces, in exact arithmetic, equal
% CHANGED (2026-11-23, review round 3, s4-5.R11): wording.
% CHANGED (2026-11-23, review round 3, s4-5.R11 review): wording.
the kernels they stand for. These are a lattice law itself, or a continuum kernel through its renormalized samples on the current grid, as each use of the sense names.
\item Let $q \ge 2$ be an integer. The \emph{subsample} of a member $\lambda$ of a grid family on
a lattice $\Gamma$ is its restriction to the coarser grid $q\Gamma$: on $h\ZZ^d$, the array
$m \mapsto \lambda(qhm)$ of every $q$-th value along each axis. It is exact when, renormalized to
mass one, it is a member of the same family on $q\Gamma$. For item (1) of
Definition~\ref{def:grid-family} that is a kernel $\kappa_\tau^{*\gamma}$ of the same stencil on
$q\Gamma$, of some range $\tau > 0$ and some index $\gamma \ge 1$. For item (2) it is a power
$(D_qw)^{*n}$, $n \ge 1$, of the same array on $q\Gamma$. The renormalizing constant is the mass
the subsample keeps, $C = \lambda(q\Gamma)$, which is required to be positive.
\item A \emph{grid pyramid} is exact when a continuous, discretely covariant hemigroup carries its
ladder (Definition~\ref{def:exact-grid-pyramid}).
\end{enumerate}
\statusT
\end{definition}

The three senses answer different questions. Sense (i) is a property of an implementation. A
floating-point run equals the arrays of an exact computation up to rounding. The difference between
the renormalized samples and the continuum kernels themselves is aliasing
(Proposition~\ref{prop:sampled-laplace-exact}(3)). It comes from sampling, and an exact computation
has it too. Proposition~\ref{prop:line-and-arrays}(2) is a statement in sense (i).

Sense (ii) is a property of a family. The constant $C$ is below $1$ whenever $\lambda$ charges a
point off the coarse grid. The subsample is the impulse response of smoothing by $\lambda$ on the
fine grid and then keeping every $q$-th value, so the averaging over neighbouring fine values is in
$\lambda$ itself. The sense is deliberately weak, so that a negative statement in it is as strong as
possible. On the line, the subsample of a sampled Laplace kernel is the kernel of the same range on
the coarser grid (Proposition~\ref{prop:sampled-laplace-exact}(2)). In the plane,
Proposition~\ref{prop:helmholtz-subsampling}(1) is a statement in sense (ii), with the range named.

Sense (iii) is a property of a mask and a ratio. It is the sense of every statement that calls a
pyramid exact in dimension $d \ge 2$, and of \emph{exact} in every sentence that carries
Proposition~\ref{prop:grid-pyramid-not-isotropic}.

The phrase \emph{exact cascade} (Propositions~\ref{prop:discrete-helmholtz}(2)
and~\ref{prop:nine-point-family}(2)) names none of the three. It is used for a family on one grid
whose members and increments satisfy the hemigroup relation exactly.

\section{The exact grid pyramid, and why it is not isotropic}
\label{sec:price-obstruction}

Exactness requires a continuous, discretely covariant family through the ladder of one mask. The
definition comes first, then the obstruction. It rests on
Lemma~\ref{lem:lattice-infinitely-divisible}, which puts the \Levy{} measure of an infinitely
divisible law carried by a lattice on that lattice.

% draft: Definition 5.2 (section 5.2), lines 763-770.
% CHANGED: the standing assumption on Gamma (a full-rank lattice in R^d) is written into the
% definition, where the draft has it in the section's opening paragraph.
% CHANGED: "positive in the sense of the pyramid literature" (an unanchored attribution, checks
% for section 5, item 11) replaced by what "positive" excludes, signed weights (blind review,
% flag 20).
% CHANGED: the ladder's steps get their own symbol w_n := D_{q^n}w (blind review, flag 1); the
% draft wrote them mu_{q^n,q^{n+1}}, the symbol of the carrying family, and (c)-(d) now prove
% that the family meets them.
% CHANGED: "isotropic if some kernel is rotation invariant" now reads "some kernel mu_t with
% t > 0": mu_0 = delta_0 is rotation invariant, and Proposition 5.4(2) is stated for t > 0
% (the q0629 referee's finding 7 on the same point).
% CHANGED: the rest of the draft's paragraph after (a)-(d) (the relation to the isotropic axioms,
% the line, the knot relation) moved to the prose after the node and to the remark on the line.
% CHANGED (the theory paper's blueprint round, 2026-11-16, review of Section 6, flag NOTATION on
% "positive"): the sentence "The mask is positive: its weights are nonnegative, as opposed to a
% mask with signed weights." is dropped. Every probability law has nonnegative weights, so it
% excluded nothing; "positive" leaves the titles of prop:grid-pyramid-not-isotropic too.
\begin{definition}[exact grid pyramid in $\RR^d$]
\label{def:exact-grid-pyramid}
\lean{Pyramids.IsLatticeMask, Pyramids.dilLadder, Pyramids.CarriesLatticeLadder, Pyramids.IsExactLatticePyramid, Pyramids.knotKernel, Pyramids.IsSORotationInvariant, Pyramids.IsIsotropicCarrier}\leanok
Let $d \ge 1$, let $\Gamma \subset \RR^d$ be a full-rank lattice, let $q > 1$, and let $w$ be a
symmetric probability law on $\Gamma$ with $w \ne \delta_0$, the \emph{mask}. The
\emph{grid pyramid} of $(w, q)$ on $\Gamma$ is its \emph{ladder}, the sequence of octave steps
$w_n := D_{q^n}w$, $n \in \ZZ$, the same array applied on the grid $q^n\Gamma$ at every level. It is
\textbf{exact} if some family $(\mu_{s,t})_{0 \le s \le t}$ of probability laws on $\RR^d$ carries
it, in the sense that
\begin{enumerate}
\item[(a)] $\mu_{s,t} * \mu_{t,u} = \mu_{s,u}$ for $s \le t \le u$, and $\mu_{t,t} = \delta_0$
(hemigroup);
\item[(b)] $(s,t) \mapsto \mu_{s,t}$ is weakly continuous on $\{0 \le s \le t\}$ (continuity in
scale);
\item[(c)] $\mu_{qs,qt} = D_q\mu_{s,t}$ for $0 \le s \le t$ (discrete covariance);
% CHANGED (2026-11-23, review round 3, s6-7.S16): wording.
\item[(d)] $\mu_{1,q} = w$.
\end{enumerate}
% CHANGED (2026-11-23, review round 3, s6-7.S16 review): wording.
A carrying family need not be symmetric and need not have (ND). The \emph{kernels} are $\mu_t := \mu_{0,t}$. The \emph{knots} are the scales $q^n$, the levels of
the pyramid, and the $\mu_{q^n}$ are the \emph{knot kernels}, all dilates of $\mu_1$ by (c). The
pyramid is \emph{isotropic} if some kernel $\mu_t$ with $t > 0$ is rotation invariant, that is
% CHANGED (2026-11-23, review round 3, s6-7.R9): wording.
% CHANGED (2026-11-23, review round 3, s6-7.R9 review): wording.
$D_R\mu_t = \mu_t$ for every $R \in \mathrm{SO}(d)$. $\mathrm{SO}(1)$ is trivial, so at $d = 1$ every pyramid
is isotropic in this sense.
% CHANGED (fourth blind pass, 2026-11-15, the session's decision): exactness is defined for every
% symmetric probability law w != delta_0 on R^d, as IsExactLatticePyramid reads no lattice and
% prop:exact-masks(1) is stated in that setting; the grid pyramid on a lattice is unchanged. The
% definition before ended at the sentence above.
% CHANGED (the theory paper's blueprint round, 2026-11-16, review of Section 6, flag NOTATION "The
% same words are used for every symmetric"): the lattice-free extension has its own name, the exact
% ladder of (w, q), and "grid pyramid" is kept for a mask carried by a lattice. The old wording:
% "The same words are used for every symmetric probability law w != delta_0 on R^d, carried by a
% lattice or not: its ladder is the sequence D_{q^n}w, n in Z, and the grid pyramid of (w, q) is
% exact if some family of probability laws on R^d has (a)-(d)." The notion is unchanged.
For every symmetric probability law $w \ne \delta_0$ on $\RR^d$, carried by a lattice or not, the
\emph{ladder} of $(w, q)$ is the sequence $D_{q^n}w$, $n \in \ZZ$. It is \textbf{exact}, the
\emph{exact ladder of $(w, q)$}, if some family of probability laws on $\RR^d$ has (a)--(d).
Kernels, knots, knot kernels and isotropy are read as above. For a mask on $\Gamma$ the ladder is
the grid pyramid, and it is exact exactly when the grid pyramid is. The words \emph{grid pyramid}
are kept for a mask carried by a lattice.
% CHANGED (fidelity review, 2026-11-15, R15): the status line says how (b) is typed.
\statusT\quad\emph{(A definition. The weak continuity (b) is typed as the continuity of
$(s,t) \mapsto \widehat\mu_{s,t}(\omega)$ on $\{0 \le s \le t\}$ for each $\omega$, which is
equivalent for probability laws by \Levy{}'s continuity theorem.)}
\end{definition}

By (c) and (d), $\mu_{q^n, q^{n+1}} = D_{q^n}\mu_{1,q} = w_n$: a family that meets one octave step
meets the whole ladder, as in Definition~\ref{def:grid-pyramid-line}. Conditions (a)--(c) are what
the axioms of \prog{iso v0.1}{Def.~3.2, p.~11} say about the laws of the operators $\Phi_{s,t}$,
with two changes. Rotation invariance is dropped: it is (I3) there, and (A3$_d$) in the operator
form of its Def.~3.1. And covariance is required at the one ratio $q$. Exactness requires a continuous
family through the ladder, not for the knot relation $D_q\mu_1 = \mu_1 * w$ alone. That choice is
what lets Lemma~\ref{lem:lattice-infinitely-divisible} act, and
Proposition~\ref{prop:binomial-knot-relation} shows why the weaker notion does not suffice.

% draft: the paragraph after Lemma 5.3, line 776.
Conditions (a)--(c) allow isotropic families: the isotropic \Matern{} pyramid of
Theorem~\ref{thm:matern-pyramid} satisfies them. What excludes isotropy is (d), the one condition
about the grid, together with the continuity (b). Read the \Levy{} measure of a law as its
distribution of jumps. Continuity makes every step of the family infinitely divisible, so the mask
has a \Levy{} measure, and by Lemma~\ref{lem:lattice-infinitely-divisible} a mask on the lattice
jumps only by lattice vectors. The knot kernel applies the mask at every finer scale, by
(\ref{eq:knot-product}) below. Its jumps are the lattice vectors shrunk by powers of $q$, a denser
set but still a countable one. A rotation-invariant law jumps equally in every direction, and in
$d \ge 2$ a sphere has uncountably many points. On the line the sphere is two points, and nothing
is excluded. Without (b) the argument has nothing to start from: the binomial mask satisfies the
knot relation with a limit that is not infinitely divisible
(Proposition~\ref{prop:binomial-knot-relation}).

The knot kernel of a carrying family is fixed by the mask. Only the mask is assumed symmetric,
not the family, so the lemma below is not a case of Proposition~\ref{prop:ladder-knot-kernel}(1).

% CHANGED (review response r1, batch 2b, 2026-11-23, finding B7): new node. The product formula,
% the symmetry of the knot kernel and its data, with the telescoping of D_q nu_1 = nu_w + nu_1 to
% nu^*, were written out in prop:grid-pyramid-not-isotropic(1) and in (a) implies (b) of
% prop:exact-masks(1); both cite it now. Its declaration assembles the ones those two Lean proofs
% call (CarriesLatticeLadder.product, .symm_one and .knot_data); linkage reads no dotted name, so the
% node is bound to the assembly.
\begin{lemma}[the knot kernel of a carrying family]
\label{lem:knot-data}
\lean{Pyramids.knot_kernel_data}\leanok
\uses{def:exact-grid-pyramid, lem:steps-infinitely-divisible, lem:symmetric-levy-khintchine,
prop:ladder-knot-kernel, thm:cited-A40}
Let $d \ge 1$ and $q > 1$, let $w$ be a symmetric probability law on $\RR^d$, and let a family
$(\mu_{s,t})_{0 \le s \le t}$ of probability laws on $\RR^d$ have (a)--(d) of
Definition~\ref{def:exact-grid-pyramid}. Then
\[
  \widehat{\mu_1}(\omega) = \lim_{M \to \infty}\prod_{m=1}^M\widehat{w}\bigl(q^{-m}\omega\bigr),
  \qquad \omega \in \RR^d,
\]
and $\mu_1$ is symmetric. If $w$ has the \Levy{} data $(A_w, \nu_w)$, then $\mu_1$ has the data
$\bigl(A_w/(q^2 - 1), \nu^*\bigr)$, with $\nu^* := \sum_{m \ge 1}D_{q^{-m}}\nu_w$.
\statusA\quad\ledger{A41}\quad\ledger{A40}\quad\ledger{A39}\quad\emph{(\textbf{Assignment.}
\ledger{A41} carries the infinite divisibility of $\mu_1$, through
Lemma~\ref{lem:steps-infinitely-divisible}. \ledger{A40} carries that its transform has no zero,
which the product formula reads. \ledger{A39} carries, through
Lemma~\ref{lem:symmetric-levy-khintchine}, the data of $\mu_1$ and their uniqueness. Held as [T],
with the proof below: the product formula, the reality of its limit, the equation of the data, and
its one solution $\nu^*$. \textbf{Hypotheses} (the node's, each with what \texttt{knot\_kernel\_data}
requires): $d \ge 1$: same; $q > 1$ and (a)--(d): same, \texttt{CarriesLatticeLadder}, with (b)
typed as the continuity of the transforms; $w$ a symmetric probability law: same; the data of $w$:
same, \texttt{IsSymData}, quantified in the last part. The limit is typed as the convergence of the
partial products to $\widehat{\mu_1}(\omega)$ at each $\omega$.)}
\end{lemma}

\begin{proof}
\leanok
\uses{lem:steps-infinitely-divisible, lem:symmetric-levy-khintchine, prop:ladder-knot-kernel,
thm:cited-A40}
\emph{The product.} The law $\mu_1 = \mu_{0,1}$ is infinitely divisible, by
Lemma~\ref{lem:steps-infinitely-divisible}. Its (I1) is (a), convolution being commutative, and its
(I2) follows from (b), as $x \mapsto e^{i\langle\omega, x\rangle}$ is bounded and continuous. By
(a), (c) and (d), $D_q\mu_1 = \mu_{0,q} = \mu_{0,1} * \mu_{1,q} = \mu_1 * w$, so
$\widehat{\mu_1}(q\omega) = \widehat{w}(\omega)\widehat{\mu_1}(\omega)$, and by induction
\[
\widehat{\mu_1}(\omega) = \widehat{\mu_1}(q^{-M}\omega)\prod_{m=1}^M\widehat{w}(q^{-m}\omega),
\qquad M \ge 1 .
\]
The transform of $\mu_1$ has no zero (\ledger{A40}). As $M \to \infty$,
$\widehat{\mu_1}(q^{-M}\omega) \to \widehat{\mu_1}(0) = 1$ by continuity of the transform, so the
partial products $\widehat{\mu_1}(\omega)/\widehat{\mu_1}(q^{-M}\omega)$ converge to
$\widehat{\mu_1}(\omega)$ at every $\omega$.

\emph{Symmetry.} The transform of the symmetric law $w$ is real, so $\widehat{\mu_1}$ is a limit of
products of real numbers, and it is real. The reflection of $\mu_1$ has the conjugate transform,
which is the same, so $\mu_1$ is symmetric.

\emph{The data.} Let $w$ have the data $(A_w, \nu_w)$. Being symmetric and infinitely divisible,
$\mu_1$ has data $(A_1, \nu_1)$, by Lemma~\ref{lem:symmetric-levy-khintchine}. The law $D_q\mu_1$
has the data $(q^2A_1, D_q\nu_1)$, by the substitution $y = qx$ in the exponent. The law
$\mu_1 * w$ has the data $(A_1 + A_w, \nu_1 + \nu_w)$, since exponents add. They are one law, so
uniqueness gives $q^2A_1 = A_1 + A_w$, that is $A_1 = A_w/(q^2 - 1)$, and
$D_q\nu_1 = \nu_w + \nu_1$. Applying $D_{q^{-1}}$ $M$ times,
\[
\nu_1 = \sum_{m=1}^M D_{q^{-m}}\nu_w + D_{q^{-M}}\nu_1 .
\]
For a Borel set $B \subseteq \{|x| > \varepsilon\}$, $\varepsilon > 0$, the last term is
$\nu_1(q^MB) \le \nu_1(\{|x| > q^M\varepsilon\})$, which tends to $0$, a \Levy{} measure being
finite off every neighbourhood of $0$. So $\nu_1$ and $\nu^*$ agree off every ball about $0$, and
neither charges $0$: $\nu_1 = \nu^*$, the argument of Proposition~\ref{prop:ladder-knot-kernel}(1).

Every step is machine-checked, and the cited facts are \ledger{A40}, with \ledger{A41} and
\ledger{A39} through the lemmas used.
\end{proof}

% draft: Proposition 5.4 (section 5.2), lines 778-802.
% CHANGED: renamed lambda_a to lambda^L_a, the sampled Laplace mask (BLUEPRINT-INPUTS), with the
% dictionary to Proposition 2.23's lambda_t (a = a_t) written in (blind review, flag 5).
% CHANGED: Burt and Adelson's name gets its citation, the page the draft gives at line 471
% (blind review, flag 21).
% CHANGED: clause (3)'s gloss ("the grid form of the frame Matern family of Proposition 6.6(1), at
% index 1 along every axis, the one index at which the subsample of that family is exact,
% Proposition 5.13(4)") moved to the prose after the proof; the knot kernel "the product over the
% axes of (2.11)" is written out, with the constant 2a/(1-a)^2 of Proposition 2.21(4).
% CHANGED: Burt and Adelson's name for the knot kernel moved from clause (1) to the prose after
% the proof, a relation to prior theory kept out of the statement (second blind review, flag 10).
% CHANGED (the theory paper's blueprint round, 2026-11-16, review of Section 6, flags SCOPE and
% NOTATION): the title names the lattice, since def:exact-grid-pyramid's last paragraph extends the
% words to laws on no lattice, under which the Matern pyramid is exact and isotropic; "positive"
% leaves the title and clause (3), as it excluded nothing. In (2) "More precisely, if" became "If
% moreover", the finite-group sentence adding a hypothesis. The old title: "no exact positive grid
% pyramid is isotropic in $d \ge 2$"; the old first sentence of (3): "Exact positive grid pyramids
% exist in every dimension." No claim changes.
\begin{proposition}[no exact grid pyramid on a lattice is isotropic in $d \ge 2$]
\label{prop:grid-pyramid-not-isotropic}
\lean{Pyramids.grid_pyramid_not_isotropic_union, Pyramids.grid_pyramid_not_isotropic_1, Pyramids.grid_pyramid_not_isotropic_2, Pyramids.grid_pyramid_not_isotropic_3}
\leanok
\uses{def:exact-grid-pyramid, lem:steps-infinitely-divisible, lem:symmetric-levy-khintchine,
lem:knot-data, lem:lattice-infinitely-divisible, prop:sampled-laplace-exact, thm:cited-A39,
thm:cited-A40, thm:cited-A41, thm:cited-A42}
Let $(w, q)$ be an exact grid pyramid on the lattice $\Gamma \subset \RR^d$, carried by the family
$(\mu_{s,t})$, and put $\Gamma_q := \bigcup_{k \in \ZZ}q^k\Gamma$, a countable set, which for
integer $q$ and $\Gamma = h\ZZ^d$ is $h\ZZ[1/q]^d$.
\begin{enumerate}
\item \emph{(The covariant kernel.)} The mask $w$ is infinitely divisible: it is the compound
Poisson law with a finite \Levy{} measure $\nu_w \ne 0$ carried by $\Gamma \setminus \{0\}$. The
knot kernel is determined by the mask,
\begin{equation}\label{eq:knot-product}
\mu_1 = \lim_{M \to \infty} D_{q^{-1}}w * D_{q^{-2}}w * \dots * D_{q^{-M}}w, \qquad
\widehat{\mu_1}(\omega) = \prod_{m \ge 1}\widehat{w}\bigl(q^{-m}\omega\bigr).
\end{equation}
Every kernel $\mu_t$, $t > 0$, is infinitely divisible with no Gaussian part and a nonzero \Levy{}
measure $\nu_t$ carried by $\Gamma_q$, and at the knots
$\nu_{q^n} = \sum_{m \ge 1} D_{q^{n-m}}\nu_w$.
\item \emph{(No isotropy.)} If $d \ge 2$, no kernel $\mu_t$, $t > 0$, is rotation invariant. If
moreover $\operatorname{supp} w$ spans $\RR^d$, then for every $t > 0$ the group
$\{R \in \Orth(d) : D_R\mu_t = \mu_t\}$ is finite.
% CHANGED (2026-11-23, review round 3, s6-7 3.3): wording.
% CHANGED (2026-11-23, review round 3, s6-7 3.3 review): wording.
\item \emph{(Existence: the frame \Matern{} pyramid on the grid.)} Exact grid pyramids on a lattice exist in every dimension. Clause (3) does not assume the opening hypotheses. For $0 < a < 1$ let $\lambda^{\mathrm L}_a(hn) := \tfrac{1-a}{1+a}a^{|n|}$, $n \in \ZZ$,
be the sampled Laplace mask on $h\ZZ$: the kernel $\lambda_t$ of
Proposition~\ref{prop:sampled-laplace-exact} at $a = a_t = e^{-h/t}$, written in its ratio $a$. For
every $q > 1$ the grid pyramid of $((\lambda^{\mathrm L}_a)^{\otimes d}, q)$ on $h\ZZ^d$ is exact.
Its knot kernel is
\[
\widehat{\mu_1}(\omega) = \prod_{i=1}^d\prod_{m \ge 1}\Bigl[1 + \frac{2a}{(1-a)^2}
\bigl(1 - \cos(q^{-m}h\omega_i)\bigr)\Bigr]^{-1},
\]
and its \Levy{} measure
\begin{equation}\label{eq:frame-levy}
\nu_1 = \sum_{i=1}^{d}\sum_{m \ge 1}\sum_{n \ne 0}\frac{a^{|n|}}{|n|}\,\delta_{nhq^{-m}e_i}
\end{equation}
is carried by the coordinate axes. The symmetry group of every knot kernel is the hyperoctahedral
group of signed permutations of the axes.
\end{enumerate}
% CHANGED (wave 3g, 2026-11-15): every clause is machine-checked, so statement and proof are marked
% proved; the status line says what each entry carries on the checked route. A40 is now spent for
% the product formula, A42 for the symmetries only (the dilates of the data are a change of
% variables), A39(iii) for exactness through prop:ladder-knot-kernel(2), and the ball integral is no
% longer needed.
% CHANGED (2026-11-15, a strengthening): the Lean of the set identity for Gamma_q drops 2 <= q and
% h > 0, which its proof does not read; the print keeps the node's setting.
% CHANGED (fidelity review, 2026-11-15, R26/R50/R69): the status line lists the declarations' hypotheses beside the node's (LINKAGE rule 6).
% CHANGED (fidelity review, 2026-11-15, R29): the status line says how the product formula and the identity for Gamma_q are typed.
\statusA\quad\ledger{A41}\quad\ledger{A40}\quad\ledger{A39}\quad\ledger{A42}\quad\emph{(\textbf{Assignment.}
\ledger{A41} is Khintchine's theorem [\sscite{sato1999levy} Thm.~9.3, p.~47]. It carries the
infinite divisibility of every law of the family in clause (1), through
Lemma~\ref{lem:steps-infinitely-divisible}. \ledger{A40} (Lemma~7.5,
p.~32) carries that the transform of the knot kernel has no zero, which the product formula reads,
through Lemma~\ref{lem:knot-data}.
That the transform is real is the reality of the factors of the product.
\ledger{A39} is Thm.~8.1, pp.~37--38. Its existence (i), p.~37, gives the triplets; its uniqueness
(ii), p.~38, gives the data of the mask and of the knot kernel, through
Lemmas~\ref{lem:symmetric-levy-khintchine} and~\ref{lem:knot-data}, the additivity
of triplets in clause (1), and the passage of a symmetry of a kernel to its \Levy{} measure in
clauses (2) and (3); its existence (iii), p.~38, gives the family of clause (3), through
Proposition~\ref{prop:ladder-knot-kernel}(2). \ledger{A42} (Prop.~11.10, p.~65) gives the triplet of
a linear image. It is printed for a \Levy{} process and read at time 1. It is spent for the
orthogonal maps of clauses (2) and (3). The lattice carrier of the mask is
Lemma~\ref{lem:lattice-infinitely-divisible}, which spends its own entries (A42, A48). Proved here:
the weak limit in the product formula, the carrier, the atom count of clause (2), and all
of clause (3) from the \Levy{} measure of $\lambda^{\mathrm L}_a$ of
Proposition~\ref{prop:grid-pyramids-infinitely-divisible}(4). The product formula and the data
of the knot kernel are Lemma~\ref{lem:knot-data}'s. Every clause is machine-checked. The
set identity for $\Gamma_q$ is checked for every natural $q$ and real $h$, its hypotheses $q \ge 2$
and $h > 0$ dropped on 2026-11-15. \textbf{Not claimed}: anything about the laws between the knots
beyond their carrier; a symmetry group when the support of the mask does not span.
The limit in (\ref{eq:knot-product}) is checked as the convergence of the transforms of the partial
convolutions and of the partial products, not as weak convergence of laws; the two are equivalent by
\Levy{}'s continuity theorem. The identity for $\Gamma_q$ divides $h$ by $q^k$ with $q$ a natural
number, a junk value at $q = 0$, where the identity still holds.
\textbf{Hypotheses} (the node's, each with what its declarations require): $d \ge 1$: absent from
\texttt{grid\_pyramid\_not\_isotropic\_1}, same in \texttt{grid\_pyramid\_not\_isotropic\_3};
$\Gamma$ a full-rank lattice: same, typed as a \texttt{Submodule} over $\ZZ$ with
\texttt{DiscreteTopology} and \texttt{IsZLattice}; $q > 1$: same, inside
\texttt{CarriesLatticeLadder}; $w$ a symmetric probability law on $\Gamma$, $w \ne \delta_0$: same
(\texttt{IsLatticeMask}); $(\mu_{s,t})$ with (a)--(d): same (\texttt{CarriesLatticeLadder}), (b)
typed as the continuity of the transforms; for the identity for $\Gamma_q$, $q$ an integer, $q > 1$,
and $\Gamma = h\ZZ^d$, $h > 0$: weaker, \texttt{grid\_pyramid\_not\_isotropic\_union} requires only
$q \in \NN$ and $h \in \RR$. (2): $d \ge 2$: same; $\operatorname{supp} w$ spans $\RR^d$: same, as
the span of the atoms of $w$. (3): $d \ge 1$, $h > 0$, $0 < a < 1$, $q > 1$: same.)}
\end{proposition}

% CHANGED (wave 3g, 2026-11-15)
% The proof follows the machine-checked route. (1): the data of the knot kernel come from the
% symmetry of mu_1 and the equation D_q nu_1 = nu_w + nu_1, whose one solution is nu* (the argument
% of prop:ladder-knot-kernel(1)), in place of the ball integral I(x), the Riemann-Lebesgue lemma and
% Tonelli; the dilates of the data are a change of variables in the exponent. (2): every point of the
% sphere through an atom is an atom, so the sphere would lie in the countable Gamma_q; the atoms of
% the mask lie in the span of the atoms of nu_w because the transform of the mask is 1 on its
% orthogonal complement, in place of the compound Poisson sum. (3): the data of the mask are read off
% the folded Levy measure of lambda^L_a (prop:grid-pyramids-infinitely-divisible(4)), and exactness
% is prop:ladder-knot-kernel(2) with the logarithmic moment, in place of the explicit coupling on the
% line and its product; the converse of the symmetry group is read on transforms. Statement
% unchanged.
% CHANGED (fourth blind pass, 2026-11-15): the data are written in the convention of
% lem:family-structure throughout, the mask's triplet (2A_w, nu_w, 0) named for the lattice lemma;
% the converse inclusion of the description of Gamma_q; the case y = x of the rotation; the Levy
% measure of the sampled Laplace mask computed here, where it was read from another node's proof;
% (a)-(c) from the family of prop:ladder-knot-kernel(2), (b) by Levy's continuity theorem.
% CHANGED (the theory paper, 2026-11-16): the closing sentence on checking and cited facts shortened; no step changed.
% CHANGED (the theory paper's blueprint round, 2026-11-16, review of Section 6, flags PROSE,
% POINTER and NOTATION): the infinite divisibility of every mu_{s,t} is cited from
% lem:steps-infinitely-divisible, where the null array and Khintchine's theorem were written out
% here; borrowed facts are cited through \ledger alone, the direct Sato page cites (Thm. 8.1,
% (8.1), (11.10), Def. 9.2, Thm. 9.3, Lemma 7.5, Prop. 11.10) dropped; the triplet's Gaussian
% matrix is A^S and the data are (A, nu) of lem:family-structure, where they were (A, nu, beta)
% and (A', nu); the finite sets of (2) are Y_i, as A_i would read as a data matrix; the
% permutation of (3) is pi, as sigma is the walk step of prop:resolvent-masks-exact; the status
% line's A41 is carried through the new lemma. No step changed.
% CHANGED (review response r1, batch 2b, 2026-11-23, findings B2 and B7): the data of a symmetric
% infinitely divisible law are lem:symmetric-levy-khintchine, and the product formula, the symmetry
% and the data of the knot kernel are lem:knot-data, both cited here as the Lean proof calls
% exists_symData and CarriesLatticeLadder.product and .knot_data. They were written out here. No
% step changed.
\begin{proof}
\leanok
\uses{lem:steps-infinitely-divisible, lem:symmetric-levy-khintchine, lem:knot-data,
lem:lattice-infinitely-divisible, prop:ladder-knot-kernel, prop:grid-pyramids-infinitely-divisible,
thm:cited-A39, thm:cited-A41, thm:cited-A42}
Throughout, an infinitely divisible law has the generating triplet $(A^{\mathrm S}, \nu, \gamma)$ of
\ledger{A39}(i), unique by \ledger{A39}(ii). Triplets add under convolution, since the transform of
a convolution is the product of the transforms and the exponents of the formula of \ledger{A39}(i)
add. A symmetric infinitely divisible law has unique \Levy{} data $(A, \nu)$ and the triplet
$(2A, \nu, 0)$, by Lemma~\ref{lem:symmetric-levy-khintchine}. The data of a dilate $D_c\mu$, $c \ne 0$, are $(c^2A, D_c\nu)$, by the
substitution $y = cx$ in the exponent. The mask is a probability law other than $\delta_0$, so
$d \ge 1$: on $\RR^0$ every law is $\delta_0$.

(1) \emph{Every $\mu_{s,t}$ is infinitely divisible}, by
Lemma~\ref{lem:steps-infinitely-divisible}. Its (I1) is (a), convolution being commutative, and its
(I2) follows from (b), as $x \mapsto e^{i\langle\omega, x\rangle}$ is bounded and continuous.

\emph{The mask.} $w = \mu_{1,q}$ is infinitely divisible, symmetric and carried by $\Gamma$. It has
data $(A_w, \nu_w)$, so the triplet $(2A_w, \nu_w, 0)$, and by
Lemma~\ref{lem:lattice-infinitely-divisible}, applied to that triplet, $A_w = 0$ and the \Levy{}
measure $\nu_w$ is finite and carried by $\Gamma \setminus \{0\}$. So
$\widehat{w} = e^{-\psi_w}$ with $\psi_w(\omega) := \int(1 - \cos\langle\omega, x\rangle)\,
\nu_w(dx)$, which is the transform of the compound Poisson law
$e^{-\nu_w(\RR^d)}\sum_{k \ge 0}\nu_w^{*k}/k!$. So $w$ is that law. Finally $\nu_w \ne 0$, since
$w \ne \delta_0$.

\emph{The product formula.} By Lemma~\ref{lem:knot-data}, the partial products
$\prod_{m=1}^M\widehat{w}(q^{-m}\omega)$ converge to $\widehat{\mu_1}(\omega)$ at every $\omega$.
They are the transforms of the partial convolutions $D_{q^{-1}}w * \dots * D_{q^{-M}}w$, which
therefore converge weakly to $\mu_1$ by L\'evy's continuity theorem. This is
(\ref{eq:knot-product}).

\emph{\Levy{} data of the knot kernel.} By the same lemma, $\mu_1$ is symmetric with the data
$(A_w/(q^2 - 1), \nu^*) = (0, \nu^*)$, as $A_w = 0$, where $\nu^* := \sum_{m \ge 1}D_{q^{-m}}\nu_w$.
Write $\nu_1 := \nu^*$. This measure is carried by $\bigcup_{m \ge 1}q^{-m}\Gamma$, and $\mu_1$ has
no Gaussian part.

\emph{All kernels.} By (c), $\mu_{q^n} = D_{q^n}\mu_1$, whose data are $(0, D_{q^n}\nu_1)$, and
$D_{q^n}\nu_1 = \sum_{m \ge 1}D_{q^{n-m}}\nu_w$, the formula at the knots. Now let $t > 0$ and
choose $n$ with $q^n \le t < q^{n+1}$. By (a), $\mu_{q^{n+1}} = \mu_t * \mu_{t,q^{n+1}}$ and
$\mu_t = \mu_{q^n} * \mu_{q^n,t}$, all four factors infinitely divisible. Triplets add. The Gaussian
parts of $\mu_t$ and $\mu_{t,q^{n+1}}$ are positive semidefinite and sum to $0$, so both vanish.
The \Levy{} measures give $D_{q^n}\nu_1 \le \nu_t \le D_{q^{n+1}}\nu_1$. The upper bound puts
$\nu_t$ on $\bigcup_{m \ge 1}q^{n+1-m}\Gamma \subseteq \Gamma_q$. The lower bound makes it nonzero,
as $D_{q^n}\nu_1 \ge D_{q^{n-1}}\nu_w \ne 0$. The set $\Gamma_q$ is countable, a countable union of
images of the lattice, which is countable as a discrete subset of $\RR^d$. For $\Gamma = h\ZZ^d$ and
integer $q$, a point $q^khm$, $m \in \ZZ^d$, with $k \ge 0$ is $h(q^km) \in h\ZZ[1/q]^d$, and with
$k = -j < 0$ it is $h\,(m/q^j) \in h\ZZ[1/q]^d$. Conversely, a point of $h\ZZ[1/q]^d$ is $h\,m/q^j$
for some $m \in \ZZ^d$ and $j \ge 0$, a common denominator of its coordinates, and this is
$q^{-j}(hm) \in q^{-j}\Gamma$. This is the description of $\Gamma_q$ in the statement.

(2) Let $t > 0$ and $R \in \Orth(d)$ with $D_R\mu_t = \mu_t$. By \ledger{A42} with $U = R$,
the \Levy{} measure of $D_R\mu_t$ is $R_*\nu_t$, as $R$ maps only $0$ to $0$, so uniqueness gives
$R_*\nu_t = \nu_t$, and $\nu_t(\{Rx\}) = \nu_t(\{x\})$ for every $x$. By (1), $\nu_t \ne 0$ is
carried by the countable set $\Gamma_q$, so it has an atom $x$, and $x \ne 0$. Suppose $\mu_t$ is
invariant under $\mathrm{SO}(d)$ and $d \ge 2$. For every $y$ with $|y| = |x|$ some
$R \in \mathrm{SO}(d)$ has $Rx = y$. If $y = x$, take $R = I$. Otherwise the reflection in the
hyperplane orthogonal to $x - y$ maps $x$ to $y$, as $|x| = |y|$, and it is followed by the
reflection in a hyperplane containing $y$, which exists as $d \ge 2$, to make the determinant $+1$. So every point of the sphere $\{|y| = |x|\}$ is an atom
of $\nu_t$, and an atom of $\nu_t$ lies in $\Gamma_q$. But the sphere is uncountable: the points
$|x|\,(e_1 + se_2)/|e_1 + se_2|$, $s \in \RR$, lie on it and are distinct, the ratio of their second
coordinate to their first being $s$. This contradiction proves the first sentence.

For the finite group, let $V$ be the span of the atoms of $\nu_w$. The measure $\nu_w$ is carried
by the countable set $\Gamma$, and its points outside $V$ carry no mass, so $\nu_w(\RR^d \setminus
V) = 0$. For $\eta$ orthogonal to $V$, $\langle\eta, x\rangle = 0$ for $\nu_w$-almost every $x$, so
$\widehat{w}(\eta) = e^{-\psi_w(\eta)} = 1$. Then $\cos\langle\eta, x\rangle = 1$ for $w$-almost
every $x$, and the same at $\eta/k$ for every integer $k \ge 1$ forces $\langle\eta, x\rangle = 0$
for $w$-almost every $x$. So $w$ is carried by the orthogonal complement of $V^\perp$, which is $V$,
and every atom of $w$ lies in $V$. If the support of $w$, its set of atoms, spans $\RR^d$, so does
$V$, and some atoms $b_1, \dots, b_d$ of $\nu_w$ form a basis. With $n$ as in (1),
$\nu_t \ge D_{q^n}\nu_1 \ge D_{q^{n-1}}\nu_w$, so the points $x_i := q^{n-1}b_i$ are atoms of
$\nu_t$, and they form a basis. For each $i$ let $Y_i$ be the set of points $y$ with
$|y| = |x_i|$ and $\nu_t(\{y\}) \ge \nu_t(\{x_i\})$. It is finite, because
$\nu_t(\{|y| > |x_i|/2\}) < \infty$ and each point of $Y_i$ carries at least the mass
$\nu_t(\{x_i\}) > 0$. A symmetry $R$ is orthogonal and has $\nu_t(\{Rx_i\}) = \nu_t(\{x_i\})$, so
$Rx_i \in Y_i$. A linear map is determined by its values on a basis, so the group has at most
$\prod_i|Y_i|$ elements.

(3) \emph{The mask on the line.} The proof of
Proposition~\ref{prop:grid-pyramids-infinitely-divisible}(4) sums the transform of
$\lambda^{\mathrm L}_a$ as two geometric series and expands $\log|1 - ae^{ih\omega}|^2$, which gives
\[
\widehat{\lambda^{\mathrm L}_a}(\omega) = \frac{(1-a)^2}{1 - 2a\cos h\omega + a^2}
= \frac{1}{1 + \frac{2a}{(1-a)^2}\bigl(1 - \cos h\omega\bigr)}
= \exp\Bigl(-\int(1 - \cos\omega v)\,\varpi(dv)\Bigr),
\]
with the finite folded \Levy{} measure $\varpi := \sum_{k \ge 1}\tfrac{2a^k}{k}\delta_{kh}$. Its
symmetric form $\varpi_s := \tfrac12\bigl(\varpi + \varpi(-\,\cdot)\bigr) =
\sum_{n \ne 0}\tfrac{a^{|n|}}{|n|}\delta_{nh}$ has the same integral of $1 - \cos\omega v$.

\emph{The mask in $\RR^d$.} Let $w := (\lambda^{\mathrm L}_a)^{\otimes d}$ and
$\nu_w := \sum_i(\iota_i)_*\varpi_s$, with $\iota_i(v) := ve_i$. The transform of a product law is
the product of the transforms of its factors, so
\[
\widehat{w}(\omega) = \prod_{i=1}^d\widehat{\lambda^{\mathrm L}_a}(\omega_i)
= \exp\Bigl(-\sum_{i=1}^d\int(1 - \cos\omega_iv)\,\varpi_s(dv)\Bigr)
= \exp\Bigl(-\int(1 - \cos\langle\omega, x\rangle)\,\nu_w(dx)\Bigr),
\]
and $\nu_w$ is finite, symmetric, without atom at $0$. So $w$ has the data $(0, \nu_w)$. It is a
mask on $h\ZZ^d$: a probability law; symmetric, since its transform is real; carried by
$h\ZZ^d$, since each coordinate is carried by $h\ZZ$; and not $\delta_0$, since $\nu_w \ne 0$ and
$\delta_0$ has the data $(0, 0)$.

\emph{Exactness.} Since $\log_+ x \le x$ for $x \ge 0$,
\[
\int\log_+|x|\,\nu_w(dx) = d\sum_{k \ge 1}\frac{2a^k}{k}\log_+(kh)
\le 2dh\sum_{k \ge 1}a^k < \infty .
\]
So Proposition~\ref{prop:ladder-knot-kernel}(2) gives a symmetric $q$-covariant family with
$S_qt = qt$ and $\mu_{1,q} = w$. Its cascade (I1) is (a), convolution being commutative; its
continuity (I2) of the transforms gives (b) by \Levy's continuity theorem; and its covariance at
$q$ with $S_qt = qt$ is (c). So the grid pyramid of $(w, q)$ is exact, for every $q > 1$.

\emph{Every family carrying it.} By (1), the knot kernel is (\ref{eq:knot-product}), and its
factors are
$\widehat{w}(q^{-m}\omega) = \prod_i\bigl[1 + 2a(1-a)^{-2}\bigl(1 - \cos(q^{-m}h\omega_i)\bigr)
\bigr]^{-1}$, those displayed in the statement. By Lemma~\ref{lem:knot-data}, as $w$ is symmetric
with data $(0, \nu_w)$,
\[
\nu_1 = \sum_{m \ge 1}D_{q^{-m}}\nu_w
= \sum_{i=1}^d\sum_{m \ge 1}\sum_{n \ne 0}\frac{a^{|n|}}{|n|}\,\delta_{nhq^{-m}e_i},
\]
which is (\ref{eq:frame-levy}), carried by the coordinate axes.

\emph{The symmetry group.} Let $R$ be a signed permutation, $Re_j = s_je_{\pi(j)}$ with
$s_j = \pm1$. Then $(R^{\tp}\omega)_j = s_j\omega_{\pi(j)}$. Each factor
$f_m(u) := \bigl[1 + 2a(1-a)^{-2}\bigl(1 - \cos(q^{-m}hu)\bigr)\bigr]^{-1}$ is even, so the
partial product $\prod_{m \le M}\prod_jf_m(\omega_j)$ takes the same value at $R^{\tp}\omega$ as at
$\omega$, and in the limit $\widehat{\mu_1}(R^{\tp}\omega) = \widehat{\mu_1}(\omega)$, that is
$D_R\mu_1 = \mu_1$. As $D_{q^n}$ commutes with $R$, the same holds for every knot kernel
$\mu_{q^n} = D_{q^n}\mu_1$. Conversely let $R \in \Orth(d)$ have $D_R\mu_{q^n} = \mu_{q^n}$. Then
$D_R\mu_1 = \mu_1$, and $R_*\nu_1 = \nu_1$ as in (2). The point $(h/q)e_i$ is an atom of $\nu_1$
(the term $m = 1$, $n = 1$), so $R(h/q)e_i$ is an atom of $\nu_1$, hence a point $ce_j$ of the
axes, with $|c| = h/q$ because $R$ is orthogonal. So $Re_i = \pm e_{\pi(i)}$ for some
$\pi(i)$. The map $\pi$ is injective: if $\pi(i) = \pi(k)$, then $Re_i = \pm Re_k$, so
$e_i = \pm e_k$ and $i = k$. So $\pi$ is a permutation, and $R$ is the signed permutation
$e_i \mapsto \pm e_{\pi(i)}$.

% CHANGED (the theory paper, 2026-11-16): the closing sentence names the cited facts by \ledger{}
% instead of raw entry numbers, and (2) gives Sato's page for Prop. 11.10. Meaning unchanged.
Every step is machine-checked, and the cited facts are \ledger{A41}, \ledger{A40}, \ledger{A39},
\ledger{A42} and \ledger{A48}.
\end{proof}

Run downward in scale, the product (\ref{eq:knot-product}) is the equivalent weighting function of
Burt and Adelson [\sscite{burt1983laplacian}, p.~533]. The frame pyramid of clause (3) is the
grid form of the frame \Matern{} family of
Proposition~\ref{prop:separable-frame}(1), at index $1$ along every axis. That is the one index at
which the subsample of that family is exact in sense (ii)
(Proposition~\ref{prop:helmholtz-subsampling}(4)).

% CHANGED (review round 2, batch C, 2026-11-23): new node, the author's decision D1 of
% records/pyramids/PLAN-review-response-theory-round2.md (open question (3) of the theory paper's
% section 7.2, similarities x -> qRx): a proposition with its printed proof, not ready, as
% prop:subsampling-discrepancy is. Its statement is typed in the skeleton against a matrix form of
% def:exact-grid-pyramid (CarriesMatrixLadder, a new definition; CarriesLatticeLadder is unchanged).
% CHANGED (2026-11-23, review round 3 batch C, item 4, s6-7.S2 and s1-3 3.1): a strengthening. The
% matrix M is invertible, where it was expanding ("every eigenvalue of M, real or complex, has
% modulus greater than 1"); the label keeps its old name. The proof argues from one atom of the
% Levy measure of a kernel, nu_t >= M^{n-1}_* nu_w, which needs M invertible only; the Gaussian
% part, the series for nu_1, the decay of M^{-n} and the countable carrier Gamma_M are no longer
% read. The skeleton declaration is renamed matrix_grid_pyramid_not_isotropic, over
% CarriesMatrixLadder with M invertible.
\begin{proposition}[no exact grid pyramid with a matrix dilation is isotropic in $d \ge 2$]
\label{prop:expanding-grid-pyramid-not-isotropic}
\lean{Pyramids.matrix_grid_pyramid_not_isotropic}
\notready
\uses{def:exact-grid-pyramid, prop:grid-pyramid-not-isotropic, lem:steps-infinitely-divisible,
lem:lattice-infinitely-divisible, thm:cited-A39, thm:cited-A41, thm:cited-A42}
Let $d \ge 2$, let $\Gamma \subset \RR^d$ be a full-rank lattice, let $w$ be a symmetric
probability law on $\Gamma$ with $w \ne \delta_0$, and let $M$ be an invertible real $d \times d$
matrix. Let $q > 1$, and let a family $(\mu_{s,t})_{0 \le s \le t}$ of probability laws on
$\RR^d$ have (a), (b) and (d) of Definition~\ref{def:exact-grid-pyramid} and, in place of (c),
\begin{enumerate}
\item[(c$_M$)] $\mu_{qs,qt} = M_*\mu_{s,t}$ for $0 \le s \le t$, where $M_*\mu$ is the image of
$\mu$ under $x \mapsto Mx$.
\end{enumerate}
% CHANGED (review round 2, batch C, blind register review, 2026-11-23): the remarks on what is not
% assumed (M Gamma in Gamma, M a multiple of an orthogonal matrix) and the quincunx example move to
% the prose before the node; "exact" is said through the conditions, which define it for D_q only.
% CHANGED (2026-11-23, review round 3, s6-7 3.1): wording.
Then no kernel $\mu_t$, $t > 0$, is rotation invariant. So no family with (a), (b), (c$_M$) and (d) makes the grid pyramid of $(w, M)$, the ladder $M^n_*w$ applied on the grids $M^n\Gamma$, isotropic in the sense of Definition~\ref{def:exact-grid-pyramid}.
% CHANGED (2026-11-23, review round 3 batch C, item 4): the assignment follows the atom argument,
% and the price of the formalization is recomputed from it (it was 600 to 1000 lines for the
% expanding form, whose proof also needed the decay of M^{-n} and the tail of nu_1).
\statusA\quad\ledger{A41}\quad\ledger{A39}\quad\ledger{A42}\quad\emph{(\textbf{Assignment.}
\ledger{A41} carries the infinite divisibility of every law of the family, through
Lemma~\ref{lem:steps-infinitely-divisible}. \ledger{A39} gives the triplets, their additivity and
their uniqueness, and \ledger{A42} the triplet of a linear image, here under the powers of $M$ and
the rotations. The lattice carrier of the mask is Lemma~\ref{lem:lattice-infinitely-divisible},
which spends its own entries (A42, A48). Proved here, in print only: everything else. Nothing in the
proof reads that $M$ is expanding, a multiple of an orthogonal matrix or that $M\Gamma \subseteq
\Gamma$. \textbf{Left}, and why the node is not marked proved: the proof is not formalized, by the
author's decision (D1 of the second review round's plan); its statement is typed in the skeleton.
Priced at about 250 to 400 lines: the steps $\mu_{q^{n-1},q^n} = M^{n-1}_*w$ for every integer $n$
from (c$_M$) and (d) (about 50); the \Levy{} measure of $M^{n-1}_*w$ through \ledger{A42} and the
bound $\nu_t \ge M^{n-1}_*\nu_w$ from the additivity of triplets over three factors (about 100 to
150); an atom of $\nu_w$ off $0$ (Lemma~\ref{lem:lattice-infinitely-divisible}); and the rotation
step, from that of Proposition~\ref{prop:grid-pyramid-not-isotropic}(2), with the finiteness of
$\nu_t$ off a ball in place of the countable carrier (about 100 to 150). \textbf{Hypotheses} (the
node's, each with what \texttt{matrix\_grid\_pyramid\_not\_isotropic} requires): $d \ge 2$: same;
$\Gamma$ a full-rank lattice: same, as in Proposition~\ref{prop:grid-pyramid-not-isotropic}; $w$ a
symmetric probability law on $\Gamma$, $w \ne \delta_0$: same (\texttt{IsLatticeMask}); $M$
invertible: same (\texttt{IsUnit} in \texttt{CarriesMatrixLadder}); $q > 1$, (a), (b), (c$_M$)
and (d): same (\texttt{CarriesMatrixLadder}), (b) typed as the continuity of the transforms;
rotation invariance: same, under $\mathrm{SO}(d)$ (\texttt{IsSORotationInvariant}).)}
\end{proposition}

% CHANGED (2026-11-23, review round 3 batch C, item 4): the proof argues from one atom of nu_t. It
% replaces the old route, which ran the proof of prop:grid-pyramid-not-isotropic(1) with M_* for
% D_q: the knot relation M_* mu_1 = mu_1 * w, the Gaussian part M A_1 M^T = A_1 forced to 0 and the
% series nu_1 = sum_m M^{-m}_* nu_w through the decay of M^{-n} (Gelfand's formula), the bounds
% between the knots, and the countable carrier Gamma_M of nu_t.
\begin{proof}
\uses{prop:grid-pyramid-not-isotropic, lem:steps-infinitely-divisible,
lem:lattice-infinitely-divisible, thm:cited-A39, thm:cited-A41, thm:cited-A42}
Triplets are those of \ledger{A39}. They add under convolution, and the triplet of a law is unique.

\emph{The mask.} Every $\mu_{s,t}$ is infinitely divisible, by
Lemma~\ref{lem:steps-infinitely-divisible}, as in Proposition~\ref{prop:grid-pyramid-not-isotropic}(1).
The mask $w = \mu_{1,q}$ is carried by $\Gamma$, so by Lemma~\ref{lem:lattice-infinitely-divisible}
its triplet has Gaussian part $0$ and a finite \Levy{} measure $\nu_w$ carried by
$\Gamma \setminus \{0\}$, and $\nu_w \ne 0$ as $w \ne \delta_0$. A nonzero measure carried by a
countable set charges one of its points, so $\nu_w$ has an atom $b \ne 0$.

\emph{The steps.} By (c$_M$) at $(s/q, t/q)$, $\mu_{s,t} = M_*\mu_{s/q,t/q}$, and as $M$ is
invertible, $\mu_{s/q,t/q} = M^{-1}_*\mu_{s,t}$. By induction,
$\mu_{q^ks,q^kt} = M^k_*\mu_{s,t}$ for every integer $k$, and with (d),
$\mu_{q^{n-1},q^n} = M^{n-1}_*w$ for every integer $n$.

\emph{An atom of $\nu_t$.} Let $t > 0$, and choose an integer $n$ with $q^n \le t$. By (a),
\[
\mu_t = \mu_{q^{n-1}} * \mu_{q^{n-1},q^n} * \mu_{q^n,t} ,
\]
three infinitely divisible laws, so the \Levy{} measure $\nu_t$ of $\mu_t$ is the sum of theirs. By
\ledger{A42} with $U = M^{n-1}$, the \Levy{} measure of $\mu_{q^{n-1},q^n} = M^{n-1}_*w$ is
$M^{n-1}_*\nu_w$, as $M^{n-1}$ maps only $0$ to $0$. So $\nu_t \ge M^{n-1}_*\nu_w$, and the point
$x := M^{n-1}b \ne 0$ has $\nu_t(\{x\}) \ge \nu_w(\{b\}) > 0$.

\emph{No isotropy.} Suppose $\mu_t$ is invariant under $\mathrm{SO}(d)$. For
$R \in \mathrm{SO}(d)$, \ledger{A42} with $U = R$ gives the \Levy{} measure $R_*\nu_t$ to
$D_R\mu_t = \mu_t$, so $R_*\nu_t = \nu_t$ by uniqueness, and $\nu_t(\{Rx\}) = \nu_t(\{x\})$. As
$d \ge 2$, every point $y$ of the sphere $\{|y| = |x|\}$ is $Rx$ for some $R \in \mathrm{SO}(d)$, and
the sphere is infinite, both as in the proof of Proposition~\ref{prop:grid-pyramid-not-isotropic}(2).
So for every $N$ the sphere holds $N$ distinct points of mass $\nu_t(\{x\}) > 0$ each, and
$\nu_t(\{|y| \ge |x|\}) \ge N\nu_t(\{x\})$. But a \Levy{} measure is finite off every
neighbourhood of $0$, and $|x| > 0$. This contradiction proves the first sentence; the second is the
first, read through the definition of isotropy.

% CHANGED (review round 2, batch C, blind register review, 2026-11-23): the sentences on what the
% proof reads of M, said before the node, are dropped here.
The cited facts are
\ledger{A41}, \ledger{A39} and \ledger{A42}, with \ledger{A48} through
Lemma~\ref{lem:lattice-infinitely-divisible}.
\end{proof}

% draft: Remark 5.8(3), line 864, the remark on the line of Definition 5.2's comment (line 770)
% and the sentence on [V-line v0.1, Thm. 5.2] in the proof of Proposition 5.4(1) (line 788),
% consolidated into the section's one remark on the line (blind review, flag 18).
% CHANGED: "the admissible ones" now reads "exact in the sense of def:grid-pyramid-line", the word
% that definition uses (second blind review, flag 1).
\emph{On the line.} In $d = 1$ the exact grid pyramids are those that are exact in the sense of
Definition~\ref{def:grid-pyramid-line}, on operators (Proposition~\ref{prop:line-and-arrays}(1)). There
Proposition~\ref{prop:grid-pyramid-not-isotropic}(2) is empty, since $\Orth(1) = \{\pm 1\}$, and
clause (1) is the product formula of Proposition~\ref{prop:grid-pyramids-infinitely-divisible}(2)
on $\Gamma = h\ZZ$. The argument for infinite divisibility in clause (1) is that of
\prog{V-line v0.1}{Thm.~5.2, p.~26}, which does not use the dimension.

% CHANGED (2026-11-23, review round 3 batch C, item 13, s4-5.R2): new prose and a new node.
For an integer ratio the kernels of an exact grid pyramid have no density at all, in every
dimension. The mask is carried by the grid, so its transform is $1$ on the dual grid, and the knot
relation then holds the transform of the knot kernel at one value along a sequence of frequencies
that tends to infinity.

% CHANGED (2026-11-23, review round 3 batch C, item 13, s4-5.R2 and plan item 1): new node.
% Clauses (1) and (2) are machine-checked; clause (3) is proved in print only, so the node is
% not ready, and its statement is typed in the skeleton. The singularity is the zero-one argument
% of Jessen and Wintner's law of pure types, written out for this series (the verifier's route, on
% the countable group hZ[1/q]^d), so that no cited fact enters beyond the ledger's.
\begin{proposition}[the kernels of an exact grid pyramid are singular]
\label{prop:grid-pyramid-singular-kernels}
\lean{Pyramids.grid_pyramid_singular_kernels_1, Pyramids.grid_pyramid_singular_kernels_2,
Pyramids.grid_pyramid_singular_kernels_3}
\notready
\uses{def:exact-grid-pyramid, lem:steps-infinitely-divisible, lem:knot-data,
prop:grid-pyramid-not-isotropic, thm:cited-A39, thm:cited-A40, thm:cited-A41}
Let $q \ge 2$ be an integer and $h > 0$, and let $(w, q)$ be an exact grid pyramid on $h\ZZ^d$,
carried by the family $(\mu_{s,t})$.
\begin{enumerate}
\item \emph{(One value along a ray.)} For every coordinate vector $e_i$,
\[
  \widehat{\mu_1}\bigl(2\pi q^ke_i/h\bigr) = \widehat{\mu_1}\bigl(2\pi e_i/h\bigr) \ne 0,
  \qquad k \ge 0 .
\]
\item \emph{(No density.)} No kernel $\mu_t$, $t > 0$, is absolutely continuous with respect to
Lebesgue measure.
\item \emph{(Singular.)} Every kernel $\mu_t$, $t > 0$, is singular with respect to Lebesgue
measure: it is carried by a Lebesgue-null set.
\end{enumerate}
\statusA\quad\ledger{A41}\quad\ledger{A40}\quad\ledger{A39}\quad\emph{(\textbf{Assignment.}
\ledger{A41} carries the infinite divisibility of $\mu_1$ and of the steps of the family, through
Lemma~\ref{lem:steps-infinitely-divisible}. \ledger{A40} carries that the transform of $\mu_1$ has
no zero, the inequality of (1) and the lower bound of (2). \ledger{A39} carries, in (3), the
triplets of the steps between two knots and their additivity; the data of the mask are those of
Proposition~\ref{prop:grid-pyramid-not-isotropic}(1), which spends its own entries (A42, A48).
Proved here: (1) and (2), machine-checked, the Riemann--Lebesgue lemma being Mathlib's; (3) in
print only, with two textbook facts of probability that are not ledger entries, the tightness of a
weakly convergent sequence of laws and the $\pi$--$\lambda$ theorem. \textbf{Not claimed}: that the
kernels have no atom, which holds as their \Levy{} measures have infinite mass
[\sscite{sato1999levy}, Thm.~27.4, p.~175] and which no node reads; a ratio that is not an integer,
for which the frequencies $q^k\xi$ leave the dual grid. \textbf{Left}, and why the node is not
marked proved: clause (3) is not formalized. Priced at about 500 to 800 lines: the independent
variables on a product space and the convergence in probability of their partial sums (about 200),
the zero-one step for an event that agrees almost surely with an event of every tail $\sigma$-algebra
(about 100 to 200), the null set $N + h\ZZ[1/q]^d$ (about 50), and the compound Poisson carrier of
a step between two knots (about 150 to 250). \textbf{Hypotheses} (the node's, each with what its
declarations require): $q \ge 2$ an integer: same, a natural number \texttt{q} with $1 < q$ inside
\texttt{CarriesLatticeLadder}; $h > 0$: same; $(w, q)$ an exact grid pyramid on $h\ZZ^d$ carried by
$(\mu_{s,t})$: same, \texttt{IsLatticeMask} on \texttt{gridSet d h} and
\texttt{CarriesLatticeLadder}, (b) typed as the continuity of the transforms; $d \ge 1$, implicit in
$w \ne \delta_0$: same, derived. (1): $e_i$: same, \texttt{ebasis i}, the frequency written
$(q^k \cdot 2\pi/h)\,e_i$. (2): absolute continuity: same, \texttt{AbsolutelyContinuous} with respect to
\texttt{volume}. (3): singularity: same, \texttt{MutuallySingular} with \texttt{volume}, in the
skeleton.)}
\end{proposition}

\begin{proof}
\uses{lem:steps-infinitely-divisible, lem:knot-data, prop:grid-pyramid-not-isotropic,
thm:cited-A39, thm:cited-A40, thm:cited-A41}
Write $\xi := 2\pi e_i/h$ and $G := h\ZZ[1/q]^d$, which is
$\bigcup_{k \ge 0}q^{-k}h\ZZ^d$ (Proposition~\ref{prop:grid-pyramid-not-isotropic}(1)).

(1) By (a), (c) and (d), $D_q\mu_1 = \mu_{0,q} = \mu_1 * w$, so
$\widehat{\mu_1}(q\omega) = \widehat{w}(\omega)\,\widehat{\mu_1}(\omega)$, as in the proof of
Lemma~\ref{lem:knot-data}. For $k \ge 0$ and $n \in \ZZ^d$,
$\langle q^k\xi, hn\rangle = 2\pi q^kn_i$ is an integer multiple of $2\pi$, $q$ being an integer.
So $e^{i\langle q^k\xi, x\rangle} = 1$ at every point $x$ of $h\ZZ^d$, which carries $w$, and
$\widehat w(q^k\xi) = 1$. Hence $\widehat{\mu_1}(q^{k+1}\xi) = \widehat{\mu_1}(q^k\xi)$, and by
induction $\widehat{\mu_1}(q^k\xi) = \widehat{\mu_1}(\xi)$ for every $k \ge 0$. The law
$\mu_1 = \mu_{0,1}$ is infinitely divisible, by Lemma~\ref{lem:steps-infinitely-divisible}, so its
transform has no zero (\ledger{A40}).

(2) Let $t > 0$, and choose an integer $n \ge 0$ with $t \le q^n$. By (a),
$\mu_{q^n} = \mu_t * \mu_{t,q^n}$, so $\widehat{\mu_{q^n}} = \widehat{\mu_{t,q^n}}\,\widehat{\mu_t}$,
and as the transform of a probability law has modulus at most $1$,
$|\widehat{\mu_t}(\omega)| \ge |\widehat{\mu_{q^n}}(\omega)|$ for every $\omega$. By (c),
$\mu_{q^n} = D_{q^n}\mu_1$, so at $\omega_k := q^k\xi$ (1) gives
$\widehat{\mu_{q^n}}(\omega_k) = \widehat{\mu_1}(q^{n+k}\xi) = \widehat{\mu_1}(\xi)$. So
$|\widehat{\mu_t}(\omega_k)| \ge |\widehat{\mu_1}(\xi)| > 0$ for every $k$, while
$|\omega_k| = 2\pi q^k/h \to \infty$. If $\mu_t$ had a density, its transform would tend to $0$ at
infinity, by the Riemann--Lebesgue lemma. So $\mu_t$ is not absolutely continuous.

(3) \emph{A random series.} Let $X_1, X_2, \dots$ be independent random vectors, $X_m$ with law
$D_{q^{-m}}w$, and $T_M := X_1 + \dots + X_M$, $T_0 := 0$, whose law is
$\rho_M := D_{q^{-1}}w * \dots * D_{q^{-M}}w$. By (\ref{eq:knot-product}), $\rho_M \to \mu_1$
weakly. A weakly convergent sequence of probability laws on $\RR^d$ is tight: for $\eta > 0$ there
is $R$ with $\rho_j(\{|x| > R\}) < \eta$ for every $j$. For $M < M'$ the difference
$T_{M'} - T_M = X_{M+1} + \dots + X_{M'}$ has the law $D_{q^{-M}}\rho_{M'-M}$, so for
$\varepsilon > 0$,
\[
  P\bigl(|T_{M'} - T_M| > \varepsilon\bigr) = \rho_{M'-M}\bigl(\{|x| > q^M\varepsilon\}\bigr) < \eta
  \qquad\text{once } q^M\varepsilon \ge R .
\]
So $(T_M)$ is a Cauchy sequence in probability. It converges in probability to a random vector
$S$, and almost surely along a subsequence, and the law of $S$ is the weak limit of the $\rho_M$,
which is $\mu_1$.

\emph{A zero-one law.} Each $X_m$ takes its values in $q^{-m}h\ZZ^d \subseteq G$. The set $G$ is a
group, the union of the increasing groups $q^{-k}h\ZZ^d$ (increasing as $q$ is an integer), so
$T_M \in G$ for every $M$. Let $A \subseteq \RR^d$ be a Borel set with $A + g = A$ for every
$g \in G$, and $E := \{S \in A\}$. Then $E = \{S - T_M \in A\}$ for every $M$. The vector
$S - T_M$ is the limit in probability of $T_{M'} - T_M$ as $M' \to \infty$, so it agrees almost
surely with the limit along a subsequence, a function of $X_{M+1}, X_{M+2}, \dots$. So $E$ agrees
almost surely with an event $E_M$ of $\mathcal{F}_M := \sigma(X_{M+1}, X_{M+2}, \dots)$, for every
$M \ge 0$. The event $E_M$ is independent of $\sigma(X_1, \dots, X_M)$, and so is $E$, which
differs from it by a null set. The union of the $\sigma(X_1, \dots, X_M)$, $M \ge 1$, is a
$\pi$-system that generates $\mathcal{F}_0$, so by the $\pi$--$\lambda$ theorem $E$ is independent
of $\mathcal{F}_0$, which contains $E_0$. Then $P(E) = P(E \cap E_0) = P(E)\,P(E_0) = P(E)^2$, and
$\mu_1(A) = P(E)$ is $0$ or $1$. This is Kolmogorov's zero-one argument as Jessen and Wintner use
it for their law of pure types [\sscite{jessen1935distribution}, Thm.~35, p.~86],
[\sscite{sato1999levy}, Prop.~27.18, p.~182].

\emph{The knot kernel is singular.} By (2), $\mu_1$ is not absolutely continuous, so some Borel set
$N$ of Lebesgue measure $0$ has $\mu_1(N) > 0$. The set $A := N + G = \bigcup_{g \in G}(N + g)$ is a
countable union of null Borel sets, so it is a null Borel set, and $A + g = A$ for $g \in G$. As
$\mu_1(A) \ge \mu_1(N) > 0$, the zero-one law gives $\mu_1(A) = 1$: $\mu_1$ is carried by the null
set $A$.

\emph{Every kernel.} Let $t > 0$, and choose $n \in \ZZ$ with $q^n \le t < q^{n+1}$. By (a),
$\mu_t = \mu_{q^n} * \mu_{q^n,t}$, and $\mu_{q^n} = D_{q^n}\mu_1$ is carried by the null set
$q^nA$. The step $\mu_{q^n,t}$ is carried by a countable set. Indeed, by (a), (c) and (d),
$D_{q^n}w = \mu_{q^n,q^{n+1}} = \mu_{q^n,t} * \mu_{t,q^{n+1}}$, three infinitely divisible laws
(Lemma~\ref{lem:steps-infinitely-divisible}), whose triplets of \ledger{A39} add. By
Proposition~\ref{prop:grid-pyramid-not-isotropic}(1) the mask has the data $(0, \nu_w)$ with
$\nu_w$ finite and carried by $h\ZZ^d \setminus \{0\}$, so the symmetric law $D_{q^n}w$ has the
triplet $(0, D_{q^n}\nu_w, 0)$. Two positive semidefinite Gaussian parts with sum $0$ vanish, so
$\mu_{q^n,t}$ has a triplet $(0, \nu', \gamma')$ with $\nu' \le D_{q^n}\nu_w$, a finite measure
carried by $q^nh\ZZ^d$. Its transform is therefore
$\exp\bigl(i\langle b, \omega\rangle + \int(e^{i\langle\omega, x\rangle} - 1)\,\nu'(dx)\bigr)$ with
$b := \gamma' - \int_{\{|x| \le 1\}}x\,\nu'(dx)$, the transform of
$\delta_b * e^{-\nu'(\RR^d)}\sum_{k \ge 0}\nu'^{*k}/k!$. So $\mu_{q^n,t}$ is this law, carried by
the countable set $C := b + q^nh\ZZ^d$. The set $q^nA + C$ is a countable union of translates of
$q^nA$, so it is null, and
\[
  \mu_t(q^nA + C) = \int\mu_{q^n}(q^nA + C - y)\,\mu_{q^n,t}(dy) = 1,
\]
since $q^nA + C - y \supseteq q^nA$ for $y \in C$. So $\mu_t$ is singular.

Clauses (1) and (2) are machine-checked, and their cited facts are \ledger{A40} and \ledger{A41};
clause (3) is proved here in print only, with \ledger{A39} and the entries of
Proposition~\ref{prop:grid-pyramid-not-isotropic}(1).
\end{proof}

The proof of (1) and (2) uses of the grid only a point $\xi \ne 0$ of its dual, a point with
$\langle\xi, \gamma\rangle \in 2\pi\ZZ$ for every $\gamma$ of the lattice. Every full-rank lattice
has one, so the proposition holds on every full-rank lattice, with such a point in place of
$2\pi e_i/h$ in (1) and $\Gamma_q$ in place of $h\ZZ[1/q]^d$ in the proof of (3). So the pyramid principle, which reads continuous densities, never applies to a
family that carries an exact grid pyramid of integer ratio.

The selection of the \Matern{} pyramid in the plane rejects the frame pyramid's octave step on
the continuum, while each of its axis passes has the implicit-Euler principle of
Definition~\ref{def:implicit-euler-principle} (Theorem~\ref{thm:selection-stratum-two}(4)). On the
grid the same holds, in every dimension $d \ge 2$, and the reason can be read off the inverse of
the mask.

% draft: Theorem 4.10(4) (section 4.3), line 711, the grid case: "The octave step of the frame
% Matern pyramid ... fails it in d >= 2 ... on the grid, while each of its axis passes has it";
% its proof, line 721: "The grid is Proposition 5.11(4)." Moved here by adr/0004, decision 3.
% CHANGED: the grid principle is the N-principle of def:neighbourhood-principle (06-helmholtz.tex),
% read on hZ^d; clause (1) is stated for every finite symmetric neighbourhood and every d >= 2,
% where prop:grid-implicit-euler-principle-frame has N_8 on Z^2; clause (2) reads "each of its
% axis passes has it" as the principle on any neighbourhood containing the two axis neighbours of
% that axis. The Z^2, N_8 case is cited, not claimed again.
% CHANGED: the principle is that of def:implicit-euler-principle on a neighbourhood of hZ^d, which
% is defined there for a probability on hZ^d; the rescaling to Z^d through
% def:neighbourhood-principle is dropped (second blind review, flag 13). No change of content.
\begin{proposition}[the frame step on the grid is not selected]
\label{prop:grid-frame-step-not-selected}
\lean{Pyramids.grid_frame_step_not_selected_1, Pyramids.grid_frame_step_not_selected_2}
\leanok
% CHANGED (wave 3h, 2026-11-15): the uses follow the checked proof; the inverse of the mask comes
% from the three-point recursion of prop:grid-implicit-euler-principle-frame, not from the
% transform of prop:sampled-laplace-exact, which is not read.
% CHANGED (the split, adr/0007, 2026-11-16): prop:grid-implicit-euler-principle-frame and
% prop:grid-implicit-euler-principle leave the uses here and in the proof. Since the fourth blind
% pass the proof checks the three-point recursion itself; the one step it read from the computation
% paper, the associativity of summable convolution against a bounded function, is stated with its
% reason. The checked proof reads the same three helper facts, now in a file of this paper
% (Formalization/Pyramids/Line/KconvBdd.lean). Statement unchanged.
\uses{def:implicit-euler-principle, prop:grid-pyramid-not-isotropic}
Let $d \ge 2$, $h > 0$ and $0 < a < 1$, and let $w := (\lambda^{\mathrm L}_a)^{\otimes d}$ be the
frame mask on $h\ZZ^d$ of Proposition~\ref{prop:grid-pyramid-not-isotropic}(3), the octave step of
the frame \Matern{} pyramid on the grid. Let $\mathcal{N}$ range over the finite symmetric sets of
nonzero points of $h\ZZ^d$, and read the principle on $\mathcal{N}$ as in
Definition~\ref{def:implicit-euler-principle}.
\begin{enumerate}
\item For every such $\mathcal{N}$, the frame mask $w$ does not have the principle on $\mathcal{N}$.
\item For each $i$, the axis pass along $e_i$, the law $\lambda^{\mathrm L}_a$ carried by the axis
$h\ZZ e_i$, has the principle on $\mathcal{N}$ for every such $\mathcal{N}$ that contains $\pm he_i$.
\end{enumerate}
% CHANGED (wave 3h, 2026-11-15): machine-checked, so statement and proof are marked proved; the
% inverse of the mask is read off its three-point recursion, and no transform is summed.
% CHANGED (2026-11-15, a strengthening): the Lean of clause (1) holds for every set N of indices,
% and that of clause (2) for every d >= 1 and every set N containing the two indices +-e_i; the
% proofs do not read that N is finite, symmetric or free of 0, nor d >= 2 in clause (2). The print
% keeps the node's setting.
% CHANGED (fidelity review, 2026-11-15, R26/R50/R69): the status line lists the declarations' hypotheses beside the node's (LINKAGE rule 6).
\statusT\quad\emph{(Cites nothing. At $d = 2$ and $\mathcal{N} = \mathcal{N}_8$, clause (1) is
Proposition~\ref{prop:grid-implicit-euler-principle-frame}, and at $d = 1$ the computation of
clause (2) is that node's ``line''. This node extends both to every $d \ge 2$ and every
neighbourhood with the same inverse. Machine-checked, with no axiom; the proof reads neither the
finiteness nor the symmetry of $\mathcal{N}$, nor that $0 \notin \mathcal{N}$.
\textbf{Hypotheses} (the node's, each with what its declarations require): $d \ge 2$: same in
\texttt{grid\_frame\_step\_not\_selected\_1}, absent from \texttt{grid\_frame\_step\_not\_selected\_2};
$h > 0$: absent, the arrays and $\mathcal{N}$ are read in the indices $\ZZ^d$ of $h\ZZ^d$;
$0 < a < 1$: same; $\mathcal{N}$ finite, symmetric, $0 \notin \mathcal{N}$: weaker, any set of
indices. (2): $\pm he_i \in \mathcal{N}$: same, as $\pm e_i$ in indices.)}
\end{proposition}

% CHANGED (wave 3h, 2026-11-15)
% The proof follows the machine-checked route. The inverse on the line is read off the three-point
% recursion of the geometric values (as in prop:grid-implicit-euler-principle-frame), in place of
% the transform of lambda^L_a; the inverse in d dimensions is checked as a finite sum over the box
% {-1,0,1}^d, factored coordinate by coordinate; f(0) < 0 is read off the signs of its terms.
% Statement unchanged.
% CHANGED (fourth blind pass, 2026-11-15): the three-point recursion is checked here, where it was
% read from the proof of prop:grid-implicit-euler-principle-frame; the commutativity of summable
% convolution is said where the inverse is used on the other side.
\begin{proof}
\leanok
We read kernels on $h\ZZ^d$ in their indices $n \in \ZZ^d$, so that $he_i$ is the index $e_i$ and
$h$ plays no further role; $\lambda := \lambda^{\mathrm L}_a$ is a kernel on $\ZZ$. Put
$\varsigma := a/(1-a)^2 > 0$.

\emph{The inverse on the line.} The values of $\lambda$ satisfy the three-point recursion
\[
(1 + 2\varsigma)\lambda(m) - \varsigma\bigl(\lambda(m+1) + \lambda(m-1)\bigr) = \delta_0(m),
\qquad m \in \ZZ .
\]
Indeed, let $c := \tfrac{1-a}{1+a}$, and note $\varsigma(1-a)^2 = a$. For $m \ne 0$ the left side
is $c\,a^{|m|-1}\bigl((1 + 2\varsigma)a - \varsigma(1 + a^2)\bigr) = c\,a^{|m|-1}\bigl(a -
\varsigma(1-a)^2\bigr) = 0$. For $m = 0$ it is $c\,(1 + 2\varsigma(1-a)) = c\,(1 + 2a/(1-a)) = 1$.
So the kernel $m_1 := (1 + 2\varsigma)\delta_0 - \varsigma(\delta_1 + \delta_{-1})$, carried by
$\{-1, 0, 1\}$, has $m_1 * \lambda = \delta_0$.

\emph{The inverse of the frame mask.} Let $M := m_1^{\otimes d}$, $M(n) := \prod_jm_1(n_j)$, a
kernel carried by the box $B := \{-1, 0, 1\}^d$. For a kernel $g$ on $\ZZ$ and $n \in \ZZ^d$, the
convolution $(M * g^{\otimes d})(n)$ is a finite sum over $B$, and expanding a product of $d$ finite
sums gives
\[
(M * g^{\otimes d})(n) = \sum_{y \in B}\prod_jm_1(y_j)\,g(n_j - y_j)
= \prod_j\,\sum_{y_j \in \{-1, 0, 1\}}m_1(y_j)\,g(n_j - y_j) .
\]
At $g = \lambda$ each factor is $\delta_0(n_j)$ by the recursion, so $M * w = \delta_0$. The mask
$w$ is summable: a finite set of indices lies in the box of its coordinates, over which the sum of
$w$ is a product of $d$ finite sums of $\lambda$, each at most $1$. Convolution of two summable
kernels is associative against a bounded function, since the double sum converges absolutely,
by the product of the two masses and the bound.

(1) Let $u := -1$ on the set $D$ of the four indices $\pm e_1 \pm e_2$ and $u := 0$ elsewhere, and
let $f := M * u$, bounded since $M$ is finitely supported and $u$ is bounded. Then
$w * f = (w * M) * u = u$, as $w * M = M * w = \delta_0$, convolution of summable kernels being
commutative. Since $u \le 0 = u(0)$, the origin is a maximum of $u$ over the whole
grid, so it is an $\mathcal{N}$-maximum for every $\mathcal{N}$. Next,
$f(0) = \sum_yM(y)\,u(-y)$. A term is nonzero only if $-y \in D$, that is $y \in D$, and there
$y_1, y_2 \in \{\pm1\}$ and $y_j = 0$ for $j \ge 3$, so
\[
M(y) = m_1(\pm1)\,m_1(\pm1)\prod_{j \ge 3}m_1(0) = \varsigma^2(1 + 2\varsigma)^{d-2} > 0 .
\]
So every term of $f(0)$ is $\le 0$, the four nonzero ones are $-\varsigma^2(1 + 2\varsigma)^{d-2}$,
and
\[
f(0) = -4\varsigma^2(1 + 2\varsigma)^{d-2} < 0 = u(0) .
\]
The principle would require $u(0) \le f(0)$ at an $\mathcal{N}$-maximum, so it fails. At $d = 2$,
$f(0) = -4\varsigma^2$.

(2) Let $\kappa$ be the law $\lambda$ carried by the axis $\ZZ e_i$, and
$m^{(i)} := (1 + 2\varsigma)\delta_0 - \varsigma(\delta_{e_i} + \delta_{-e_i})$. Then
$(m^{(i)} * \kappa)(n) = (1 + 2\varsigma)\kappa(n) - \varsigma\bigl(\kappa(n + e_i) +
\kappa(n - e_i)\bigr)$. If $n$ lies on the axis, so do $n \pm e_i$, and this is the recursion at
$n_i$, equal to $\delta_0(n_i) = \delta_0(n)$. Otherwise $n$ has a nonzero coordinate $n_j$,
$j \ne i$, which $n \pm e_i$ share, so the three values of $\kappa$ vanish, and so does
$\delta_0(n)$. So $m^{(i)} * \kappa = \delta_0$, and $\kappa$ is summable, its values being those
of $\lambda$. For a bounded $f$ and $u := \kappa * f$, associativity gives
$f = (m^{(i)} * \kappa) * f = m^{(i)} * u$, that is
\[
f(x) = u(x) + \varsigma\bigl(2u(x) - u(x + e_i) - u(x - e_i)\bigr) .
\]
At an $\mathcal{N}$-maximum $x$ of $u$ with $\pm e_i \in \mathcal{N}$ the bracket is nonnegative,
so $u(x) \le f(x)$. At an $\mathcal{N}$-minimum it is nonpositive, so $f(x) \le u(x)$.

Every step is machine-checked, with no cited fact.
\end{proof}

% CHANGED: the comparison with prop:grid-implicit-euler-principle, a statement on Z^2, is made at
% d = 2 only (second blind review, flag 7).
In every $d \ge 2$ the inverse of the frame mask charges the diagonal neighbours $\pm he_1 \pm he_2$
with the positive weight $\varsigma^2(1 + 2\varsigma)^{d-2}$. At $d = 2$ this is the language of
Proposition~\ref{prop:grid-implicit-euler-principle}: the inverse of a mask with the principle has
nonpositive weights off the origin, and the inverse of the frame mask does not. The failure does
not depend on the neighbourhood, because the counterexample's origin is a global maximum.

\section{Which masks are exact, and what the algorithm computes}
\label{sec:price-masks}

Proposition~\ref{prop:grid-pyramid-not-isotropic}(1) gives a necessary condition: the mask is
infinitely divisible. With a logarithmic moment it is also sufficient, in every dimension. That
settles the masks of the computations, the frame mask and the resolvents of the five-point and
nine-point stencils, and it says what the pyramid algorithm returns.

% draft: Proposition 5.5 (section 5.3), lines 808-825.
% CHANGED (the theory paper's blueprint round, 2026-11-16): the law w need not be carried by a
% lattice in (1), so the title and (a) say "exact ladder", the name def:exact-grid-pyramid now gives
% the lattice-free notion. The old title: "which masks have an exact grid pyramid"; the old (a):
% "The grid pyramid of $(w, q)$ is exact." No claim changes.
\begin{proposition}[which laws have an exact ladder]
\label{prop:exact-masks}
\lean{Pyramids.exact_masks_1, Pyramids.exact_masks_1_poisson, Pyramids.exact_masks_2}
\leanok
% CHANGED (wave 3g, 2026-11-15): the uses follow the checked proof, which reaches every cited fact
% through the nodes it uses.
% CHANGED (fourth blind pass, 2026-11-15): thm:cited-A39 and thm:cited-A41 join the uses, which
% the proof now cites itself.
% CHANGED (the theory paper's blueprint round, 2026-11-16): lem:steps-infinitely-divisible joins
% the uses; the proof cites it for the infinite divisibility of the laws of a carrying family.
\uses{def:exact-grid-pyramid, lem:steps-infinitely-divisible, lem:symmetric-levy-khintchine,
lem:knot-data, lem:lattice-infinitely-divisible, lem:synthesis,
prop:grid-pyramid-not-isotropic, prop:ladder-knot-kernel, prop:no-bounded-step, thm:cited-A39,
thm:cited-A41}
% CHANGED (2026-11-15, a strengthening, the author's decision)
% The setting is split clause by clause: the criterion holds for every symmetric probability law
% w other than delta_0 on R^d; the lattice is assumed for the last sentence of (1) and for (2). The
% old wording: "Let $w \ne \delta_0$ be a symmetric probability law on $\Gamma$ and $q > 1$." with
% (1)'s last sentence "An infinitely divisible $w$ is the compound Poisson law ..." and (2) "If
% these hold, ...".
Let $w \ne \delta_0$ be a symmetric probability law on $\RR^d$ and $q > 1$.
\begin{enumerate}
\item \emph{(Criterion.)} The following are equivalent.
\begin{enumerate}
\item[(a)] The ladder of $(w, q)$ is exact.
\item[(b)] $w$ is infinitely divisible and $\int\log_+|x|\,\nu_w(dx) < \infty$.
\item[(c)] $w$ is infinitely divisible and $\int\log_+|x|\,w(dx) < \infty$.
\end{enumerate}
% CHANGED (fourth blind pass, 2026-11-15): the lattice is bound where it is assumed. The old
% wording: "If $w$ is carried by $\Gamma$, an infinitely divisible $w$ ..." and (2) "Let $w$ be
% carried by $\Gamma$."
If $w$ is carried by a full-rank lattice $\Gamma$, an infinitely divisible $w$ is the compound
Poisson law with a finite symmetric \Levy{} measure $\nu_w$ carried by $\Gamma \setminus \{0\}$, and
its support is unbounded.
\item \emph{(Ladder and family.)} Let $w$ be carried by a full-rank lattice $\Gamma$. If (a)--(c)
hold, the knot
kernel $\mu_1$ is that of
(\ref{eq:knot-product}), with no Gaussian part and \Levy{} measure $\nu^* := \sum_{m \ge 1}
D_{q^{-m}}\nu_w$. The families carrying the ladder are not unique. One of them is the
\emph{uniform coupling}, with \Levy{} measures $\nu_t := \nu^* + (\log_qt)\,\nu_w$ for
$t \in [1, q]$, $\nu_{q^nt} := D_{q^n}\nu_t$, and no Gaussian part.
\end{enumerate}
% CHANGED (wave 3g, 2026-11-15): every clause is machine-checked, so statement and proof are marked
% proved. The checked proof reads every cited fact through a node it uses (A46 through
% prop:no-bounded-step, A49 and A50 through prop:ladder-knot-kernel(2), A39 through
% lem:family-structure, lem:synthesis and the nodes above), so the node is [T] by the grading
% convention, as prop:grid-pyramids-infinitely-divisible is on the line.
% CHANGED (fourth blind pass, 2026-11-15): regraded [A]. The proof now writes out the infinite
% divisibility of the laws of a carrying family and the data of a symmetric infinitely divisible
% law, which it had read from inside the proofs of lem:family-structure and
% prop:grid-pyramid-not-isotropic, and so cites A41 and A39 itself. The status was "\statusT ...
% (No ledger entry: every cited fact is carried by the nodes used. Lemma family-structure gives the
% infinite divisibility of every law of a carrying family (Khintchine's theorem, A41) and the
% symmetric data of a symmetric infinitely divisible law (A39); ...".
% CHANGED (fidelity review, 2026-11-15, R26/R50/R69): the status line lists the declarations' hypotheses beside the node's (LINKAGE rule 6).
\statusA\quad\ledger{A41}\quad\ledger{A39}\quad\emph{(\textbf{Assignment.} \ledger{A41},
Khintchine's theorem, carries the infinite divisibility of every law of a carrying family, through
Lemma~\ref{lem:steps-infinitely-divisible}; it is read in (a) implies (b).
\ledger{A39} carries, through Lemma~\ref{lem:symmetric-levy-khintchine}, that a symmetric
infinitely divisible law has symmetric data and that the data of a law are unique, read for the
mask in (1) and for the \Levy{} measure throughout. Not carried by either, and reached through the
nodes used: the symmetry and the data of the knot kernel, Lemma~\ref{lem:knot-data} (A40); the lattice carrier, Lemma~\ref{lem:lattice-infinitely-divisible} (A42, A48); the unbounded
support, Proposition~\ref{prop:no-bounded-step} (A46); the knot kernel of (2),
Proposition~\ref{prop:grid-pyramid-not-isotropic}(1); the equivalence of the two moments and the
existence of a family from the moment, Proposition~\ref{prop:ladder-knot-kernel}(2) (A49, A50); and
the couplings, which Lemma~\ref{lem:synthesis} builds from their data (A39(iii)). Held as [T], with
the proof below: the moment, the
assembly, that a symmetric point mass is $\delta_0$, and the \Levy{} measures of the uniform
coupling at every scale. Every clause is machine-checked. The criterion (1) is stated for every symmetric probability law
$w \ne \delta_0$ on $\RR^d$, since its proof does not read that $w$ is carried by $\Gamma$ (the
author's decision of 2026-11-15, a strengthening). The lattice enters the last sentence of (1) and,
through $A_w = 0$, the absence of a Gaussian part in (2), which assume it. \textbf{Not claimed}: the behaviour of the partial products when the moment fails.
\textbf{Hypotheses} (the node's, each with what its declarations require): $d \ge 1$: absent, no
declaration requires it; $w$ a probability law on $\RR^d$, symmetric, $w \ne \delta_0$: same; $q > 1$:
same in \texttt{exact\_masks\_1} and \texttt{exact\_masks\_2}, absent from
\texttt{exact\_masks\_1\_poisson}, which does not read $q$. (1), last sentence: $w$ carried by a
full-rank lattice $\Gamma$: same (\texttt{IsLatticeMask}, $\Gamma$ a \texttt{Submodule} over
$\ZZ$ with \texttt{DiscreteTopology} and \texttt{IsZLattice}); $w$ infinitely divisible: same.
(2): $w$ carried by $\Gamma$: same, as in (1); (a)--(c): same, required as (a), which gives (b) and (c)
by (1); $\nu_w$ the \Levy{} measure of $w$: same, bound by the hypothesis
\texttt{HasLevyMeasure}~$w$~$\nu_w$, which (b) supplies.)}
\end{proposition}

% CHANGED (wave 3g, 2026-11-15)
% The proof follows the machine-checked route. (a) implies (b) reads the data of the knot kernel
% from prop:grid-pyramid-not-isotropic(1) and the moment from nu* by the lower bound of the bracket
% alone; (b) implies (a), the non-uniqueness and the uniform coupling are the families of
% prop:ladder-knot-kernel(2), built by lem:synthesis, in place of the construction written out here;
% the second family is the one with the weight (log_q t)^2; the Levy measures of the uniform
% coupling are read off the data that lem:synthesis gives at every scale. The unbounded support is
% prop:no-bounded-step. Statement unchanged.
% CHANGED (fourth blind pass, 2026-11-15): the data of a symmetric infinitely divisible law and the
% infinite divisibility of the laws of a carrying family are proved here from A39 and A41 (now in
% the uses), where they were read from inside other proofs; the data identity q^2 A_1 = A_1 + A_w
% and the iteration to nu^* are written out; (b) implies (a) says how the family of
% prop:ladder-knot-kernel(2) meets (a)-(c); (2) reads the knot kernel from the statement of
% prop:grid-pyramid-not-isotropic(1).
% CHANGED (length take (b), 2026-11-16): wording tightened, route and steps unchanged.
% CHANGED (the theory paper, 2026-11-16): the closing sentence on checking and cited facts shortened; no step changed.
% CHANGED (the theory paper's blueprint round, 2026-11-16, review of Section 6, flags PROSE,
% POINTER and NOTATION): the infinite divisibility of the laws of a carrying family is cited from
% lem:steps-infinitely-divisible, where the null array and Khintchine's theorem were written out
% here; the bare references to thm:cited-A39 became the ledger macro for A39, and Sato's Thm. 25.3
% with its page the ledger macro for A49; the triplet is (A^S, nu, gamma) and the data
% (A, nu) of lem:family-structure, where they were (A, nu, beta) and (A', nu); the weight of the
% couplings is ell, as a is the Laplace ratio; a family carries "the ladder", not "the grid
% pyramid", in (a) implies (b) and (b) implies (a), w being on no lattice there. No step changed.
% CHANGED (review response r1, batch 2b, 2026-11-23, findings B2 and B7): the data of a symmetric
% infinitely divisible law are lem:symmetric-levy-khintchine, and the symmetry and data of the knot
% kernel in (a) implies (b) are lem:knot-data, both cited as the Lean proof calls exists_symData
% and CarriesLatticeLadder.knot_data. They were written out here. No step changed.
\begin{proof}
\leanok
\uses{lem:steps-infinitely-divisible, lem:symmetric-levy-khintchine, lem:knot-data,
lem:lattice-infinitely-divisible, lem:synthesis, prop:grid-pyramid-not-isotropic,
prop:ladder-knot-kernel, prop:no-bounded-step, thm:cited-A39, thm:cited-A41}
\emph{Data.} A symmetric infinitely divisible law has unique \Levy{} data $(A, \nu)$, so a unique
\Levy{} measure, and the triplet $(2A, \nu, 0)$, by Lemma~\ref{lem:symmetric-levy-khintchine}.
Data add under convolution, the exponents adding, and $D_c\mu$ has the data $(c^2A, D_c\nu)$ by the
substitution $y = cx$.

% CHANGED (2026-11-15, a strengthening, the author's decision): the last sentence assumes the
% lattice, the criterion does not.
(1) \emph{The last sentence.} Let $w$ be carried by $\Gamma$ and infinitely divisible, with data
$(A_w, \nu_w)$. By
Lemma~\ref{lem:lattice-infinitely-divisible}, $A_w = 0$ and $\nu_w$ is finite and carried by
$\Gamma \setminus \{0\}$. So $\widehat w = e^{-\psi_w}$ with $\psi_w(\omega) := \int(1 -
\cos\langle\omega, x\rangle)\,\nu_w(dx)$, the transform of the compound Poisson law
$e^{-\nu_w(\RR^d)}\sum_{k \ge 0}\nu_w^{*k}/k!$ of the symmetric measure $\nu_w$, and by the uniqueness
of the transform $w$ is that law. It is not a point mass, a symmetric $\delta_x$ being $\delta_0$.
By Proposition~\ref{prop:no-bounded-step}, a law of bounded support that is not a point mass is not
infinitely divisible, so the support of $w$ is unbounded: $w(\{|x| > R\}) > 0$ for every $R$.

\emph{(a) implies (b).} Let $(\mu_{s,t})$ carry the ladder. Every $\mu_{s,t}$ is infinitely
divisible, by Lemma~\ref{lem:steps-infinitely-divisible}, which applies since $d \ge 1$, $w$ being
a law other than $\delta_0$. Its (I1) is condition (a) of Definition~\ref{def:exact-grid-pyramid},
convolution being commutative, and its (I2) follows from the weak continuity (b) there, as
$x \mapsto e^{i\langle\omega, x\rangle}$ is bounded and continuous. So $w = \mu_{1,q}$ is
infinitely divisible, and it has data $(A_w, \nu_w)$. By Lemma~\ref{lem:knot-data}, $\mu_1$ has the
data $(A_w/(q^2 - 1), \nu^*)$, with $\nu^* := \sum_{m \ge 1}D_{q^{-m}}\nu_w$. So $\nu^*$ is a
\Levy{} measure, and by Tonelli
\[
\int\min\bigl(1, |x|^2\bigr)\,\nu^*(dx) = \int\Bigl[\sum_{m \ge 1}\min\bigl(1,
q^{-2m}|v|^2\bigr)\Bigr]\,\nu_w(dv) < \infty .
\]
For $|v| > 1$ the $\lfloor\log_q|v|\rfloor$ terms of the bracket with $q^m \le |v|$ equal $1$, so
the bracket is at least $\log_q|v| - 1$. Hence
$\log_+|v| \le \log q\,\bigl(\text{bracket} + \mathbf 1_{\{|v| > 1\}}\bigr)$ for every $v$, and
$\nu_w(\{|v| > 1\}) < \infty$ gives $\int\log_+|x|\,\nu_w(dx) < \infty$. This is
Proposition~\ref{prop:ladder-knot-kernel}(1)'s last sentence, for a family assumed only to carry
the ladder.

\emph{(b) implies (a).} The mask $w$ is a symmetric infinitely divisible law other than $\delta_0$
with $\int\log_+|x|\,\nu_w(dx) < \infty$. By Proposition~\ref{prop:ladder-knot-kernel}(2) some
symmetric $q$-covariant family with $S_qt = qt$ has $\mu_{1,q} = w$. Its cascade (I1) is (a) of
Definition~\ref{def:exact-grid-pyramid}, convolution being commutative; the continuity (I2) of its
transforms gives the weak continuity (b) by \Levy's continuity theorem; its covariance at $q$ with
$S_qt = qt$ is (c); and $\mu_{1,q} = w$ is (d). So the ladder is exact.

\emph{(b) is equivalent to (c).} This is the equivalence of the two logarithmic moments in
Proposition~\ref{prop:ladder-knot-kernel}(2), whose proof reads \ledger{A49} with
$g(x) = \log(|x| \vee e)$.

(2) Let the conditions hold, and let $(A_w, \nu_w)$ be the data of $w$, so that $A_w = 0$ by (1).
Let $(\mu_{s,t})$ carry the ladder. As $w$ is carried by $\Gamma$, the ladder is an exact grid
pyramid on the lattice $\Gamma$, so
by Proposition~\ref{prop:grid-pyramid-not-isotropic}(1) the knot kernel is (\ref{eq:knot-product}),
with no Gaussian part and the \Levy{} measure $\nu_{q^0} = \sum_{m \ge 1}D_{q^{-m}}\nu_w = \nu^*$.

Families carrying the ladder are built by Lemma~\ref{lem:synthesis} from data on one octave. For a
weight $\ell$, continuous and strictly increasing on $[1, q]$ with $\ell(1) = 0$ and $\ell(q) = 1$,
take the data $(0, \nu_u)$, $\nu_u := \nu^* + \ell(u)\,\nu_w$, $u \in [1, q]$. Their exponents
$\Psi^*(\xi) + \ell(u)\psi_w(\xi)$, with $\Psi^*$ the exponent of $(0, \nu^*)$, are continuous in $u$;
they increase strictly, as $\psi_w \not\equiv 0$ for $w \ne \delta_0$; and they glue, as
$D_q\nu^* = \nu_w + \nu^*$. The lemma extends them to every $t = q^nu$, $u \in [1, q)$, by
$\nu_t := D_{q^n}\nu_u$, and gives laws $\mu_{s,t}$ with the transforms $e^{-(\Psi_t - \Psi_s)}$,
$\Psi_t$ the exponent of $(0, \nu_t)$ and $\Psi_0 = 0$. They form a symmetric $q$-covariant family
with $S_qt = qt$, and $\mu_{1,q}$ has the transform $e^{-(\Psi_q - \Psi_1)} = e^{-\psi_w}$, so it is
$w$, and the family meets (a)--(d) as in (b) implies (a): it carries the ladder.

\emph{Not unique.} The weights $\log_qu$ and $(\log_qu)^2$ give two such families. At
$t = q^{1/2}$ their laws $\mu_{1,t}$ have the exponents $\tfrac12\psi_w$ and $\tfrac14\psi_w$, which
differ where $\psi_w \ne 0$.

\emph{The uniform coupling} is the family of the weight $\log_q$. Its kernel $\mu_t = \mu_{0,t}$,
$t > 0$, has the transform $e^{-\Psi_t}$, so it has the data $(0, \nu_t)$: at $t = q^nu$ with
$n = \lfloor\log_qt\rfloor$ and $u = q^{-n}t \in [1, q)$, the \Levy{} measure
$D_{q^n}\bigl(\nu^* + (\log_qu)\,\nu_w\bigr)$, the measure $\nu_t$ of the statement.

Every step is machine-checked, and the cited facts are \ledger{A41}, \ledger{A39} and those of the
statements used.
\end{proof}

% draft: the sentence after Proposition 5.5, line 817, the section's one remark on the line
% (Remark 5.8(3)'s logarithmic moment folded in; blind review, flag 18).
\emph{On the line.} Proposition~\ref{prop:exact-masks} is
Proposition~\ref{prop:grid-pyramids-infinitely-divisible}(1)--(2), with the same logarithmic
moment.

The resolvent masks have an explicit \Levy{} measure. Each is the law of a lattice random walk
stopped at an independent geometric time, and its logarithm is a series in the convolution powers
of the step.

% draft: Proposition 5.6 (section 5.3), lines 829-847.
% CHANGED: clause (3)'s "no kernel" reads "no kernel mu_t with t > 0" (mu_0 = delta_0).
% CHANGED: clause (4) is stated on Gamma = hZ^d (h Z^2 for the nine-point resolvent), where the
% instances live; Propositions 5.9 and 5.10 index their kernels by Z^d, read here on hZ^d (the
% referee of sections 5-7, item 1, finding 3). The neighbourhood N is renamed \mathcal{N}
% (BLUEPRINT-INPUTS), t^2/h^2 is written nu_h, and t is said to be the range of kappa_t (blind
% review, flag 3).
% CHANGED: in (4) the range of the resolvent is written tau, since t is the scale of the pyramid in
% (3), and the number t^2/h^2 is not given the letter nu_h beside the Levy measure nu_w: p is
% written in tau and h directly, the same value as Propositions 5.9 and 5.10 give (second blind
% review, flag 4).
\begin{proposition}[the resolvent masks are exact]
\label{prop:resolvent-masks-exact}
\lean{Pyramids.resolvent_masks_exact_1, Pyramids.resolvent_masks_exact_2, Pyramids.resolvent_masks_exact_3, Pyramids.resolvent_masks_exact_4_five, Pyramids.resolvent_masks_exact_4_nine, Pyramids.resolvent_masks_exact_4_general}
\leanok
% CHANGED (wave 3h, 2026-11-15): the uses follow the checked proof, which reads the data of a
% symmetric law (lem:family-structure) and exactness from the moment (prop:ladder-knot-kernel(2)).
% CHANGED (fourth blind pass, 2026-11-15): thm:cited-A39 joins the uses and lem:family-structure
% leaves them; the proof reads the Levy measure of w from A39 itself.
\uses{def:exact-grid-pyramid, prop:grid-pyramid-not-isotropic, prop:exact-masks,
prop:ladder-knot-kernel, prop:discrete-helmholtz, prop:nine-point-family, thm:cited-A39}
% CHANGED (2026-11-15, a strengthening, the author's decision)
% The setting is split clause by clause: clause (1) holds for every symmetric sigma other than
% delta_0 on Gamma; the finite support is assumed in (2) and in (3), which reads the ladder of (2).
% The old wording: "Let $\sigma \ne \delta_0$ be a symmetric probability law on $\Gamma$ with finite
% support, let $0 < p < 1$, and let", with (2) "For every $q > 1$ ..." and (3) "If $d \ge 2$ and ...".
Let $\sigma \ne \delta_0$ be a symmetric probability law on $\Gamma$, let $0 < p < 1$, and let
\begin{equation}\label{eq:resolvent-mask}
w := (1 - p)\sum_{k \ge 0}p^k\sigma^{*k}, \qquad \widehat{w} = \frac{1 - p}{1 - p\,\widehat{\sigma}} .
\end{equation}
\begin{enumerate}
\item $w$ is infinitely divisible. It is the compound Poisson law whose \Levy{} measure $\nu_w$ is
the restriction of $\sum_{k \ge 1}\tfrac{p^k}{k}\sigma^{*k}$ to $\Gamma \setminus \{0\}$, of total
mass at most $\log\frac{1}{1-p}$.
% CHANGED (2026-11-23, review round 3 batch C, item 5, s6-7.S6): a strengthening. The clause
% assumes a logarithmic moment of sigma, where it assumed finite support. The old wording: "Let
% $\sigma$ have finite support. For every $q > 1$ ...".
\item Let $\int\log_+|x|\,\sigma(dx) < \infty$, as when $\sigma$ has finite support. For every
$q > 1$ and every integer $\gamma \ge 1$ the grid pyramid of $(w^{*\gamma}, q)$ is exact, with the
knot kernel
$\widehat{\mu_1} = \prod_{m \ge 1}\widehat{w}(q^{-m}\,\cdot)^\gamma$.
% CHANGED (2026-11-15, a strengthening)
% The finite support of sigma is dropped from (3), which its proof does not read: the ladder of
% (w^{*gamma}, q) is named in the clause instead of through (2). The old wording: "Let $\sigma$ have
% finite support. If $d \ge 2$ and $\operatorname{supp}\sigma$ spans $\RR^d$, no kernel $\mu_t$,
% $t > 0$, of a family carrying that ladder is rotation invariant." The setting above is unchanged.
% CHANGED (2026-11-16, a strengthening): the hypothesis R Gamma = Gamma is dropped from the last
% sentence; neither the proof nor resolvent_masks_exact_3 reads it. The old wording named every
% orthogonal R with R Gamma = Gamma and R_* sigma = sigma; the clause now requires R_* sigma = sigma only.
\item If $d \ge 2$ and $\operatorname{supp}\sigma$ spans $\RR^d$, then for every $q > 1$ and every
integer $\gamma \ge 1$ no kernel $\mu_t$, $t > 0$, of a family carrying the ladder of
$(w^{*\gamma}, q)$ is rotation invariant. At every knot the symmetry group of the kernel is
finite, and it contains every $R \in \Orth(d)$ with $R_*\sigma = \sigma$.
% CHANGED (2026-11-16, notation): the range of the resolvent is written t, the letter of
% prop:discrete-helmholtz, def:grid-family and prop:creation-grid-masks(1), which chapters 6-8 read;
% in this clause t is that range, not a scale of the pyramid of (3). It was written tau.
\item \emph{(Instances.)} Let $\Gamma = h\ZZ^d$, and let $t > 0$, in this clause the range of a
resolvent, not a scale of the pyramid. The $(2d+1)$-point resolvent $\kappa_t$ of
Proposition~\ref{prop:discrete-helmholtz} is (\ref{eq:resolvent-mask}) with $\sigma$ the simple
random walk step and $p = 2dt^2/(h^2 + 2dt^2)$, and then $\nu_w$ charges every nonzero point
of $h\ZZ^d$. On $h\ZZ^2$ the nine-point resolvent $\kappa^9_t$ of
Proposition~\ref{prop:nine-point-family} is (\ref{eq:resolvent-mask}) with its step $\sigma_9$ and
$p = 10t^2/(3h^2 + 10t^2)$. More generally, a mask with
$1/\widehat{w}(\omega) = 1 + \sum_{y \in \mathcal{N}}c_y\bigl(1 - \cos\langle\omega, y\rangle\bigr)$,
for a finite symmetric $\mathcal{N} \subset \Gamma \setminus \{0\}$ and $c_y = c_{-y} \ge 0$ not
all $0$, is (\ref{eq:resolvent-mask}) with $C := \sum_yc_y$, $p := C/(1 + C)$ and
$\sigma := (2C)^{-1}\sum_yc_y(\delta_y + \delta_{-y})$.
\end{enumerate}
% CHANGED (wave 3h, 2026-11-15): machine-checked, so statement and proof are marked proved; the
% status says what each used node carries on the checked route.
% CHANGED (2026-11-15, a strengthening): the Lean of clause (1) drops the finite support of sigma,
% and that of clause (3) drops the finite support, q > 1 (part of carrying the ladder) and R Gamma =
% Gamma; their proofs do not read them. The print keeps the node's setting.
% CHANGED (fourth blind pass, 2026-11-15): regraded [A]. The proof of (1) now reads the Levy measure
% of w from the uniqueness of the triplet, A39(ii), which it had attributed to lem:family-structure
% and prop:grid-pyramid-not-isotropic, whose statements do not say it. The status began "\statusT
% ... (Cites nothing beyond the nodes it uses: the symmetric data of a law give its Levy measure and
% its infinite divisibility through Lemma family-structure and Proposition
% grid-pyramid-not-isotropic, ...".
% CHANGED (fidelity review, 2026-11-15, R26/R50/R69): the status line lists the declarations' hypotheses beside the node's (LINKAGE rule 6).
\statusA\quad\ledger{A39}\quad\emph{(\textbf{Assignment.} \ledger{A39} carries, in (1), that the
\Levy{} measure of $w$ is $\nu_w$: clause (i)'s formula with the triplet $(0, \nu_w, 0)$ is checked
here, and clause (ii), the uniqueness of the triplet, makes $\nu_w$ the \Levy{} measure. Held as
[T], with the proof below: the transform, the series of $\nu_w$, the compound Poisson form and the
infinite divisibility of $w$, and the mask property of the powers. Exactness from the logarithmic
moment is Proposition~\ref{prop:ladder-knot-kernel}(2), and clause (3) is
Proposition~\ref{prop:grid-pyramid-not-isotropic}(2), each spending its own entries. Clause (4) is
machine-checked with no axiom.
% CHANGED (2026-11-16, a strengthening): the sentence that kept R Gamma = Gamma in clause (3) goes
% with the hypothesis.
Clause (3) does not require $R\Gamma = \Gamma$ of a symmetry $R$, which the proof does not read (a
strengthening, 2026-11-16); it needs only $R_*\sigma = \sigma$. Clause
(1) is stated without the finite support of $\sigma$, which its proof does not read (the author's
decision of 2026-11-15, a strengthening).
% CHANGED (2026-11-15, a strengthening): clause (3) no longer prints the finite support.
Clause (3) is stated without it too, as its Lean does not read it: the ladder is named in the
clause, and a family carrying it exists by (2) when $\sigma$ has a logarithmic moment.
% CHANGED (2026-11-23, review round 3 batch C, item 5): clause (2) assumes the logarithmic moment.
Clause (2) assumes only a logarithmic moment of $\sigma$, where it assumed finite support (review
round 3, s6-7.S6, a strengthening of 2026-11-23); the moment of $\nu_w$ that exactness reads
follows from it. \textbf{Not claimed}: that the logarithmic moment of $\sigma$ is necessary.
\textbf{Hypotheses} (the node's, each with what its declarations require): $\Gamma$ a full-rank lattice:
same in (1)--(3), a \texttt{Submodule} over $\ZZ$ with \texttt{DiscreteTopology} and
\texttt{IsZLattice}; $\sigma$ a symmetric probability law on $\Gamma$, $\sigma \ne \delta_0$: same
(\texttt{IsLatticeMask}); $0 < p < 1$: same. (2): $\int\log_+|x|\,\sigma(dx) < \infty$: same, as
a lower integral; $q > 1$: same;
$\gamma \ge 1$: same. (3): $d \ge 2$: same; $\operatorname{supp}\sigma$ spans $\RR^d$: same, as the
span of the atoms of $\sigma$; $q > 1$: same, inside \texttt{CarriesLatticeLadder}; a family
carrying the ladder of $(w^{*\gamma}, q)$, $\gamma \ge 1$: same; $R_*\sigma = \sigma$: same.
(4): $\Gamma = h\ZZ^d$, $h > 0$, $t > 0$: same, with $d \ge 1$ required by
\texttt{resolvent\_masks\_exact\_4\_five} and $d = 2$ by \texttt{resolvent\_masks\_exact\_4\_nine};
in its last sentence, $w$ a mask: weaker, \texttt{resolvent\_masks\_exact\_4\_general} requires only a
probability law on $\RR^d$; $\mathcal{N} \subset \Gamma \setminus \{0\}$ finite and symmetric:
weaker, a finite symmetric subset of $\RR^d \setminus \{0\}$; $c_y = c_{-y} \ge 0$, not all $0$:
same.)}
\end{proposition}

% CHANGED (wave 3h, 2026-11-15)
% The proof follows the machine-checked route. (1): the logarithm series is two real series, sigma
% hat being real; the integral against nu_w is a series of integrals, not Tonelli. (2): w^{*gamma}
% is not delta_0 because its transform is below 1 somewhere; the moment is bounded by log_+ r <= r;
% exactness is prop:ladder-knot-kernel(2) through the data (0, gamma nu_w). (3): the atoms of
% w^{*gamma} are read through (mu * nu){x + y} >= mu{x} nu{y}; the symmetry is checked on the
% partial products. (4): the identities of laws are checked on transforms, with no convolution of
% laws on hZ^d; the lattice path through the same atom inequality. Statement unchanged.
% CHANGED (fourth blind pass, 2026-11-15): in (1), the compound Poisson form gives the infinite
% divisibility and A39(ii) the Levy measure, in place of two misattributed nodes; the mask property
% of the powers w^{*gamma}, which (2) and (3) read, moves into (1), where the finite support is not
% assumed; in (4) the generic step is eta, not tau, which is the range, and its geometric series is
% not attributed to prop:nine-point-family.
% CHANGED (length take (b), 2026-11-16): wording tightened, route and steps unchanged.
% CHANGED (the theory paper, 2026-11-16): the closing sentence on checking and cited facts shortened; no step changed.
% CHANGED (2026-11-16): A39 is cited as \ledger{A39}, and the range in (4) is written t, as in the
% statement; no step changed.
\begin{proof}
\leanok
\uses{prop:grid-pyramid-not-isotropic, prop:exact-masks, prop:ladder-knot-kernel,
prop:discrete-helmholtz, prop:nine-point-family, thm:cited-A39}
Write $\widehat{\mu}$ for the transform of a law $\mu$, $\psi_D(\omega) := \langle\omega,
A\omega\rangle + \int(1 - \cos\langle\omega, x\rangle)\,\nu(dx)$ for symmetric data $D = (A, \nu)$,
and say that $\mu$ has the data $D$ when $\widehat{\mu} = e^{-\psi_D}$.

(1) Each $\sigma^{*k}$ is a probability law, and $\sum_k(1-p)p^k = 1$, so $w$ is a probability law.
Its transform is the series of the transforms of the terms, $\widehat{\sigma^{*k}} =
\widehat{\sigma}^k$, and $|p\,\widehat{\sigma}| \le p < 1$, so
$\widehat{w} = (1-p)\sum_k(p\,\widehat{\sigma})^k = (1-p)/(1 - p\,\widehat{\sigma})$.

The law $\sigma$ is symmetric, so $\widehat{\sigma}(\omega) = z$ is real, with $|z| \le 1$, and so
is $\widehat{\sigma^{*k}}(\omega) = z^k$; as each $\sigma^{*k}$ is symmetric,
$\int(1 - \cos\langle\omega, x\rangle)\,\sigma^{*k}(dx) = 1 - z^k$. The measure
$\sum_{k \ge 1}\tfrac{p^k}{k}\sigma^{*k}$ has total mass $\sum_{k \ge 1}p^k/k = \log\frac{1}{1-p}$,
and so its restriction $\nu_w$ to $\RR^d \setminus \{0\}$ has mass at most $\log\frac{1}{1-p}$. As
$1 - \cos\langle\omega, \cdot\rangle$ is bounded and vanishes at the origin, its integral against
$\nu_w$ is the series of its integrals against the terms, and the real series
$\sum_{k \ge 1}x^k/k = -\log(1 - x)$ at $x = p$ and $x = pz$ give
\[
\int(1 - \cos\langle\omega, x\rangle)\,\nu_w(dx) = \sum_{k \ge 1}\frac{p^k}{k}\bigl(1 - z^k\bigr)
= \log(1 - pz) - \log(1 - p) .
\]
So $\widehat{w}(\omega) = (1-p)/(1-pz) = e^{-\psi(\omega)}$ with $\psi := \psi_{(0, \nu_w)}$. The
measure $\nu_w$ is symmetric, finite, and does not charge $0$, so $w$ has the symmetric data
$(0, \nu_w)$. The compound Poisson law $e^{-\nu_w(\RR^d)}\sum_{k \ge 0}\nu_w^{*k}/k!$ of $\nu_w$ has
the transform $\exp\int(e^{i\langle\omega, x\rangle} - 1)\,\nu_w(dx) = e^{-\psi(\omega)}$, the sine
part integrating to $0$, so it is $w$, by the uniqueness of the transform. It is infinitely
divisible, the $n$-th convolution power of the compound Poisson law of $\nu_w/n$ for every $n$. In
the formula of \ledger{A39}(i) the triplet $(0, \nu_w, 0)$ gives
$\exp\int(e^{i\langle\omega, x\rangle} - 1 - i\langle\omega, x\rangle\mathbf 1_{\{|x| \le 1\}})\,
\nu_w(dx) = e^{-\psi(\omega)}$, the compensator integrating an odd function against $\nu_w$. So by
the uniqueness of the triplet, \ledger{A39}(ii), $\nu_w$ is the \Levy{} measure of $w$. Each
$\sigma^{*k}$ is carried by $\Gamma$, closed under addition and countable, so $\nu_w$ is carried by
$\Gamma \setminus \{0\}$.

\emph{The powers.} Let $\gamma \ge 1$ be an integer. The law $w^{*\gamma}$ has the transform
$\widehat{w}^\gamma = e^{-\gamma\psi}$, so it has the data $(0, \gamma\nu_w)$. It is a mask on
$\Gamma$: a probability law, symmetric and carried by $\Gamma$ as a convolution power of such a
law, and not $\delta_0$. For the last, as $\sigma \ne \delta_0$, some $\omega$ has
$\widehat{\sigma}(\omega) = z < 1$, so $0 < (1-p)/(1-pz) < 1$ and
$\widehat{w^{*\gamma}}(\omega) = \bigl((1-p)/(1-pz)\bigr)^\gamma < 1$.

% CHANGED (2026-11-23, review round 3 batch C, item 5): the moment of nu_w from a logarithmic
% moment L of sigma, through log_+|x + y| <= log_+|x| + log_+|y| + log 2. The old proof bounded
% sigma^{*k} by the ball of radius kR, sigma carried by a finite set.
(2) By (1), $w^{*\gamma}$ is a mask with the data $(0, \gamma\nu_w)$. Let
$L := \int\log_+|x|\,\sigma(dx) < \infty$. For real $a, b \ge 0$,
$\log_+(a + b) \le \log_+(2\max(a, b)) \le \log 2 + \log_+a + \log_+b$, so
$\log_+|x + y| \le \log_+|x| + \log_+|y| + \log 2$ for $x, y \in \RR^d$. Integrating it against
$\sigma^{*(k-1)} \otimes \sigma$, whose image under $(x, y) \mapsto x + y$ is $\sigma^{*k}$, and
inducting on $k$ gives $\int\log_+|x|\,\sigma^{*k}(dx) \le k(L + \log 2)$. So
\[
\int\log_+|x|\,\nu_w(dx) \le \sum_{k \ge 1}\frac{p^k}{k}\,k(L + \log 2)
= (L + \log 2)\sum_{k \ge 1}p^k < \infty .
\]
So Proposition~\ref{prop:ladder-knot-kernel}(2), in the form that
Proposition~\ref{prop:exact-masks}(1) reads, makes the grid pyramid of $(w^{*\gamma}, q)$ exact,
with the moment of $\gamma\nu_w$. For a family carrying the ladder, the product formula
(\ref{eq:knot-product}) of Proposition~\ref{prop:grid-pyramid-not-isotropic}(1) gives
$\widehat{\mu_1}$ as the limit of the partial products of
$\widehat{w^{*\gamma}}(q^{-m}\cdot) = \widehat{w}(q^{-m}\cdot)^\gamma$.

(3) The law $w^{*\gamma}$ is a mask on $\Gamma$ by (1). For laws $\mu$, $\nu$ and points $x$, $y$,
$(\mu * \nu)\{x + y\} \ge \mu\{x\}\,\nu\{y\}$. Applied to $w^{*\gamma} = w^{*(\gamma - 1)} * w$ at
$0 + x$, with $w^{*(\gamma - 1)}\{0\} \ge w\{0\}^{\gamma - 1} \ge (1-p)^{\gamma - 1}$ and
$w\{x\} \ge (1-p)p\,\sigma\{x\}$ (the terms $k = 0$ and $k = 1$ of (\ref{eq:resolvent-mask})), it
gives $w^{*\gamma}\{x\} > 0$ at every atom $x$ of $\sigma$. So the atoms of $w^{*\gamma}$ span
$\RR^d$, and Proposition~\ref{prop:grid-pyramid-not-isotropic}(2) applies to every family carrying
the ladder: no kernel $\mu_t$, $t > 0$, is invariant under $\mathrm{SO}(d)$, and the symmetry group of every
knot kernel is finite.

Let $R \in \Orth(d)$ with $R_*\sigma = \sigma$, that is $\widehat{\sigma}(R^{\tp}\omega) =
\widehat{\sigma}(\omega)$ for every $\omega$. Then $\widehat{w^{*\gamma}}$, a function of
$\widehat{\sigma}$, is invariant in the same way, and so is every partial product of
(\ref{eq:knot-product}), since $R^{\tp}$ commutes with the dilations. The limit
$\widehat{\mu_1}$ is invariant, so $R_*\mu_1 = \mu_1$. By (c) of
Definition~\ref{def:exact-grid-pyramid} the kernel at the knot $q^n$ is $D_{q^n}\mu_1$, whose
transform is $\widehat{\mu_1}(q^n\cdot)$, invariant too.

(4) Read a kernel $\kappa$ on $\ZZ^d$ as the law $\sum_n\kappa(n)\delta_{hn}$ on $h\ZZ^d$; its
transform is the Fourier series $\sum_n\kappa(n)e^{i\langle n, \vartheta\rangle}$ at
$\vartheta = h\omega$. A probability kernel whose Fourier series is real, and with
$\kappa(0) \ne 1$, is a mask on $h\ZZ^d$: symmetric because its transform is real, and not
$\delta_0$. The simple random walk step $\sigma$ and the nine-point step $\sigma_9$ are such
kernels, with $\sigma(0) = \sigma_9(0) = 0$ and real Fourier series
(Propositions~\ref{prop:discrete-helmholtz} and~\ref{prop:nine-point-family}), carried by the
$2d$, resp.\ $8$, points of their stencils.

For a probability step $\eta$ on $\ZZ^d$ with a real Fourier series
$\varphi = \widehat{\eta}(\vartheta)$ and $a > 0$, the kernel $(1 - p)\sum_kp^k\eta^{*k}$,
$p = a/(1+a)$, has the Fourier series $(1-p)/(1 - p\varphi)$, by the geometric series as in (1),
and
\[
\frac{1 - p}{1 - p\varphi} = \frac{1/(1+a)}{(1 + a(1-\varphi))/(1+a)} = \bigl(1 + a(1 - \varphi)\bigr)^{-1} .
\]
Read on $h\ZZ^d$, it is the resolvent mask (\ref{eq:resolvent-mask}) of the law of $\eta$ with the
same $p$, the two laws having the same transform. The resolvent $\kappa_t$ of
Proposition~\ref{prop:discrete-helmholtz} is this kernel with
$\eta = \sigma$ and $a = 2dt^2/h^2$, and $\kappa^9_t$ of Proposition~\ref{prop:nine-point-family}
is this kernel with $\eta = \sigma_9$ and $a = \tfrac{10}{3}t^2/h^2$, as those propositions define
them; then $a/(1+a) = 2dt^2/(h^2 + 2dt^2)$, resp.\ $10t^2/(3h^2 + 10t^2)$.

For the simple random walk step and $n \in \ZZ^d \setminus \{0\}$, let $N := \sum_i|n_i| \ge 1$.
By induction on $N$, $\sigma^{*N}\{hn\} > 0$: if $n_i \ne 0$, let $e := \pm e_i$ with the sign of
$n_i$; then $\sum_i|(n - e)_i| = N - 1$, and $\sigma^{*N}\{hn\} \ge \sigma^{*(N-1)}\{h(n - e)\}\,
\sigma\{he\}$ with $\sigma\{he\} = 1/(2d)$. The point $hn$ is not $0$, so the term $k = N$ of
$\nu_w$ charges it, with the weight $\tfrac{p^N}{N}\sigma^{*N}\{hn\} > 0$.

For the last sentence, $C > 0$ as the $c_y$ are nonnegative and not all $0$. The law
$\sigma = (2C)^{-1}\sum_yc_y(\delta_y + \delta_{-y})$ has total mass $1$, is carried by the
symmetric $\mathcal{N} \not\ni 0$, so it is not $\delta_0$, and has the real transform
$\widehat{\sigma} = C^{-1}\sum_yc_y\cos\langle y, \cdot\rangle$, so it is symmetric. With
$S(\omega) := \sum_yc_y\cos\langle y, \omega\rangle \le C$,
\[
1 + \sum_yc_y\bigl(1 - \cos\langle\omega, y\rangle\bigr) = 1 + C - S(\omega) ,\qquad
\frac{1 - p}{1 - p\,\widehat{\sigma}(\omega)} = \frac{1/(1+C)}{(1 + C - S(\omega))/(1 + C)}
= \frac{1}{1 + C - S(\omega)} ,
\]
with $p = C/(1+C)$. So $\widehat{w} = (1 - p)/(1 - p\,\widehat{\sigma})$, and $w$ is
(\ref{eq:resolvent-mask}) by the uniqueness of the transform.

Every step is machine-checked, and the cited facts are \ledger{A39} and those of the statements
used.
\end{proof}


% CHANGED (fidelity review, 2026-11-15, R38): two glosses softened to what is declared. The masks of
% clause (4) were "exactly" the kernels with the principle on N, an iff whose converse no
% declaration states; the paragraph now gives the direction that Proposition
% grid-implicit-euler-principle yields and says that the converse is not stated. The inverse
% (delta_0 - p sigma)/(1 - p) is read as an identity of transforms, which resolvent_masks_exact_1
% carries.
Clause (4) is what the computations need. For an integer $q \ge 2$, a pyramid that applies the
five-point or nine-point resolvent on each grid and decimates computes the ladder of an exact grid
pyramid (Proposition~\ref{prop:line-and-arrays}(2)), and by clause (3) that pyramid is not
isotropic. On $\ZZ^2$ a mask of clause (4)'s last sentence has a summable inverse,
$(1 + C)\delta_0 - \sum_yc_y\delta_y$, of the form $\delta_0 - L$ with nonnegative weights on
$\mathcal{N}$, so by Proposition~\ref{prop:grid-implicit-euler-principle} it has the implicit-Euler
principle on the symmetric neighbourhood $\mathcal{N}$, and by clause (2) an exact grid pyramid.
Whether every real symmetric kernel of mass one with a summable inverse and the principle on
$\mathcal{N}$ is such a mask is not stated here. The frame mask
$\lambda^{\mathrm L}_a \otimes \lambda^{\mathrm L}_a$ is exact too
(Proposition~\ref{prop:grid-pyramid-not-isotropic}(3)) but is not of the form
(\ref{eq:resolvent-mask}). By (\ref{eq:resolvent-mask}),
$1/\widehat w = (1 - p\,\widehat\sigma)/(1 - p)$, the transform of the signed kernel
$(\delta_0 - p\sigma)/(1-p)$, whose weights off the origin are nonpositive: the inverse of a resolvent
mask in the sense of transforms. The frame mask's inverse
charges the diagonal neighbours with a positive weight
(Proposition~\ref{prop:grid-implicit-euler-principle-frame}).

Which masks does the creation of extrema on a grid, Proposition~\ref{prop:creation-grid}, cover?
Its clauses (1)--(2) need only a nonnegative transform. Its clause (3), which turns an arbitrary
array into a record, needs infinite divisibility. A record is the array that the pyramid algorithm
holds after some number of levels.

% CHANGED: new definition node (new label def:level-records). The draft and the first version of
% this chapter defined the level-k record in prose; two statements below use it, and
% prop:creation-grid(3) defines it inline (second blind review, flag 2).
\begin{definition}[the records of an input]
\label{def:level-records}
\lean{Pyramids.decimRecord}\leanok
\uses{prop:creation-grid}
% CHANGED (2026-11-15, a strengthening)
% The records are defined for every integer q >= 1, as Pyramids.decimRecord is (it takes any
% natural q), so that prop:line-and-arrays(2) can be printed without q >= 2. The old wording:
% "$q \ge 2$ an integer".
% CHANGED (2026-11-23, review round 3, s6-7.S19): wording.
Let $\Gamma \subset \RR^d$ be a full-rank lattice, $q \ge 1$ an integer, $q = 1$ being the computation without decimation, and $w$ a probability law on
$\Gamma$. The \emph{records} of a bounded function $f$ on $\Gamma$ are the arrays $f_k$ on
$q^k\Gamma$, $k \ge 0$, given by
\[
f_0 := f, \qquad f_{j+1} := \bigl(D_{q^j}w * f_j\bigr)\big|_{q^{j+1}\Gamma}, \quad j \ge 0 .
\]
The array $f_k$ is the \emph{level-$k$ record} of $f$. The convolution at a point of
$q^{j+1}\Gamma \subseteq q^j\Gamma$ reads $f_j$ only on $q^j\Gamma$, since $D_{q^j}w$ is carried by
$q^j\Gamma$, and every record is bounded by $\sup|f|$.
\statusT\quad\emph{(A definition. On $h\ZZ^d$ it is the record of
Proposition~\ref{prop:creation-grid}(3). \texttt{decimRecord} is in index coordinates, through a
basis of $\Gamma$, with $f_{j+1}(m) = (w * f_j)(qm)$; it takes a real array $w$ and any natural
$q$, and the summability of $w$ and the boundedness of $f$ are hypotheses of the statements that
use it.)}
% CHANGED (fidelity review, 2026-11-15, R13): the status line says how decimRecord reads the
% definition.
\end{definition}

On $h\ZZ^d$ two statements below rest on one reduction. If every bounded array at level $k$ is a
record, the creation of Proposition~\ref{prop:creation-grid}(1)--(2) at the spacing $q^kh$ happens
inside a pyramid.

% CHANGED: new lemma (new label lem:records-creation), the reduction that the proofs of
% prop:creation-grid-masks(3) and prop:binomial-records(2) shared word for word, stated once and
% cited twice (second blind review, flag 14). Neither statement changes.
\begin{lemma}[a mask whose records exhaust the bounded arrays creates maxima]
\label{lem:records-creation}
% CHANGED (wave 3b, 2026-11-15): the declaration is machine-checked, so statement and proof are
% marked proved.
\lean{Pyramids.records_creation}
\leanok
\uses{def:level-records, prop:creation-grid}
% CHANGED (2026-11-15, a strengthening)
% The hypothesis on the transform is dropped: Proposition creation-grid(1)-(2) no longer reads it.
% The old wording: "let $w \ne \delta_0$ be a symmetric probability law on $h\ZZ^d$ with
% $\widehat{w} \ge 0$. Let".
Let $d \ge 2$, $h > 0$, let $q \ge 2$ and $k \ge 0$ be integers, and let $w \ne \delta_0$ be a
symmetric probability law on $h\ZZ^d$. Let $\mathcal{N}$ be a
neighbourhood as in Proposition~\ref{prop:creation-grid}. Suppose that every bounded array on
$q^kh\ZZ^d$ is the level-$k$ record of a bounded input on $h\ZZ^d$. Then there is a bounded input on
$h\ZZ^d$ whose level-$k$ record has exactly one $q^k\mathcal{N}$-maximum, while the smoothing of that
record by $D_{q^k}w$ has at least two $q^k\mathcal{N}$-maxima, and the level-$(k+1)$ record has at
least two $q^{k+1}\mathcal{N}$-maxima.
% CHANGED (2026-11-15, a strengthening)
% The declaration requires less than the print, which keeps its setting. The status line was a bare
% \statusT before.
% CHANGED (fidelity review, 2026-11-15, R26/R50/R69): the status line lists the declarations' hypotheses beside the node's (LINKAGE rule 6).
\statusT\quad\emph{(Cites nothing beyond the nodes it uses: it reads
Proposition~\ref{prop:creation-grid}(1)--(2), which cite nothing. Machine-checked. The declaration
requires less than the print, which keeps its setting: it holds for every integer $q \ge 1$, and
$\mathcal{N}$ need only contain the $2d$ axis neighbours.
\textbf{Hypotheses} (the node's, each with what its declarations require): $d \ge 2$: same; $h > 0$:
absent, the levels are read in their own indices $\ZZ^d$; $q \ge 2$ an integer: weaker, $q \ge 1$;
$k \ge 0$ an integer: same; $w$ a symmetric probability law on $h\ZZ^d$, $w \ne \delta_0$: same, as
an array on $\ZZ^d$ that is nonnegative, summable, of sum $1$ (\texttt{IsProbKernel}); $\mathcal{N}$
finite, symmetric, in $h\ZZ^d \setminus \{0\}$, containing the $2d$ axis neighbours: weaker, any set
of indices that contains them; every bounded array at level $k$ a level-$k$ record of a bounded
input: same.)}
\end{lemma}

% CHANGED (wave 3b, 2026-11-15): the proof follows the machine-checked route, which reads the
% level-k grid in its own indices and bounds the array by its sum. The old proof checked the
% transform of the level-k mask, which Proposition creation-grid no longer requires, and said "a summable
% array is bounded" without the bound.
% CHANGED (fifth blind pass, 2026-11-15, flag 1): the paragraph read the level-k grid in its indices
% and called a q^kN-maximum an N-maximum there (it is an h^{-1}N-maximum), and did not check the
% hypotheses of creation-grid for q^kN. It now checks them at the spacing q^kh directly.
\begin{proof}
\leanok
\uses{def:level-records, prop:creation-grid}
The mask of the level-$k$ step, $D_{q^k}w$, is a symmetric probability law on $q^kh\ZZ^d$ other
than $\delta_0$, since $D_{q^k}$ is a bijection of $h\ZZ^d$ onto $q^kh\ZZ^d$ that commutes with
$x \mapsto -x$. The set $q^k\mathcal{N}$ is a finite symmetric subset of $q^kh\ZZ^d \setminus \{0\}$
that contains the $2d$ axis neighbours $\pm q^khe_i$, because $\mathcal{N}$ is such a subset of
$h\ZZ^d \setminus \{0\}$ and contains $\pm he_i$. So the hypotheses of
Proposition~\ref{prop:creation-grid} hold at the spacing $q^kh$, with the mask $D_{q^k}w$ and the
neighbourhood $q^k\mathcal{N}$, and its clauses (1)--(2) give a summable $f_k > 0$ on $q^kh\ZZ^d$
with exactly one $q^k\mathcal{N}$-maximum,
such that $D_{q^k}w * f_k$ has at least two $q^k\mathcal{N}$-maxima and its restriction to
$q^{k+1}h\ZZ^d$ has at least two $q^{k+1}\mathcal{N}$-maxima. Being positive and summable, $f_k$ is
bounded by its sum: $f_k(x) \le \sum_yf_k(y)$ for every $x$. By the hypothesis, $f_k$ is the
level-$k$ record of a bounded input $f$ on $h\ZZ^d$. By Definition~\ref{def:level-records}, the
level-$(k+1)$ record of $f$ is $D_{q^k}w * f_k$ restricted to $q^{k+1}h\ZZ^d$, which has the two
maxima.
\end{proof}

% draft: the paragraph after Proposition 4.7 (section 4.2), line 646. Moved here by adr/0004,
% decision 3; new label prop:creation-grid-masks.
% CHANGED: the paragraph's claims are stated as clauses, with the masks written out; "the tensor
% powers of the binomial mask" is read as its d-fold tensor power on hZ^d (Remark 5.8(1) says
% "its tensor powers on hZ^d"); "the count fails to be monotone from rung to rung, before and after
% decimation" is clause (3); the binomial records are prop:binomial-records, after the node.
\begin{proposition}[the masks of the creation on a grid]
\label{prop:creation-grid-masks}
\lean{Pyramids.zdLaw_tensorIdx_eq_tensorPow, Pyramids.creation_grid_masks_1, Pyramids.creation_grid_masks_2, Pyramids.creation_grid_masks_3_resolvent, Pyramids.creation_grid_masks_3_frame}
% CHANGED (wave 3b, 2026-11-15): every declaration is machine-checked, so statement and proof are
% marked proved. The checked proof of (2) reads prop:no-bounded-step for the binomial tensor mask,
% in place of the last sentence of prop:exact-masks(1).
\leanok
% CHANGED (2026-11-23, review round 3 batch C, item 6, s6-7.S8): prop:exact-masks and
% prop:sampled-laplace-exact leave the uses. The frame mask's infinite divisibility is read from its
% data, computed in the proof of prop:grid-pyramid-not-isotropic(3), not from its exact grid pyramid
% through the criterion of prop:exact-masks(1); its transform is no longer read in (1).
\uses{prop:creation-grid, def:level-records, lem:records-creation, prop:discrete-helmholtz,
prop:resolvent-masks-exact, prop:grid-pyramid-not-isotropic, prop:no-bounded-step}
Let $d \ge 2$ and $h > 0$. Let $\beta := \tfrac{1}{16}(\delta_{-2h} + 4\delta_{-h} + 6\delta_0 +
4\delta_h + \delta_{2h})$ be the binomial mask on $h\ZZ$.
% CHANGED (2026-11-23, review round 3 batch C, item 6, s6-7.S8): "with a nonnegative transform" is
% dropped from (1); nothing reads it, as prop:creation-grid does not require it. The old wording:
% "... different from $\delta_0$, with a nonnegative transform: the ...".
\begin{enumerate}
\item Each of the following is a symmetric probability law on $h\ZZ^d$, different from
$\delta_0$: the $(2d+1)$-point resolvent $\kappa_t$ of
Proposition~\ref{prop:discrete-helmholtz}, $t > 0$; the frame mask
$(\lambda^{\mathrm L}_a)^{\otimes d}$ of Proposition~\ref{prop:grid-pyramid-not-isotropic}(3),
$0 < a < 1$; and the tensor power $\beta^{\otimes d}$. So
Proposition~\ref{prop:creation-grid}(1)--(2) applies to each of them.
\item The first two are infinitely divisible. The third is not.
\item Let $w$ be one of the first two, $q \ge 2$ an integer, $k \ge 0$ an integer, and
$\mathcal{N}$ a neighbourhood as in Proposition~\ref{prop:creation-grid}. There is a bounded input
on $h\ZZ^d$ whose level-$k$ record (Definition~\ref{def:level-records}) has exactly one $q^k\mathcal{N}$-maximum, while the smoothing of
that record by $D_{q^k}w$ has at least two $q^k\mathcal{N}$-maxima, and the level-$(k+1)$ record
has at least two $q^{k+1}\mathcal{N}$-maxima.
\end{enumerate}
% CHANGED (2026-11-15, a strengthening)
% The declarations require less than the print, which keeps its setting. The status line was a bare
% \statusT before.
% CHANGED (fidelity review, 2026-11-15, R26/R50/R69): the status line lists the declarations' hypotheses beside the node's (LINKAGE rule 6).
% CHANGED (2026-11-23, review round 3 batch C, item 6): the status says what (2) reads now.
\statusT\quad\emph{(Cites nothing beyond the nodes it uses. Machine-checked. The declarations require
less than the print, which keeps its setting: (1) and (2) hold for every $d \ge 1$, and (3) for
every integer $q \ge 1$ and every $\mathcal{N}$ that contains the $2d$ axis neighbours. In (2) the
frame mask is infinitely divisible as a compound Poisson law, from the data computed in the proof of
Proposition~\ref{prop:grid-pyramid-not-isotropic}(3); since 2026-11-23 (review round 3, s6-7.S8)
it is no longer read from its exact grid pyramid through Proposition~\ref{prop:exact-masks}(1), so
\ledger{A49} leaves the cited facts of (2) and (3). Clause (1) no longer states that the three
transforms are nonnegative, which nothing read.
\textbf{Hypotheses} (the node's, each with what its declarations require): $d \ge 2$: weaker in
\texttt{creation\_grid\_masks\_1} and \texttt{creation\_grid\_masks\_2}, which require $d \ge 1$, same in
the two declarations of (3); $h > 0$: same, absent from \texttt{creation\_grid\_masks\_3\_frame},
which reads the frame mask in the indices $\ZZ^d$; $t > 0$, $0 < a < 1$: same; a one-dimensional
array nonnegative and summable: absent from the node, required by
\texttt{zdLaw\_tensorIdx\_eq\_tensorPow}, which identifies the two spellings of a tensor power, and
met by $\beta$ and $\lambda^{\mathrm L}_a$. (3): $q \ge 2$ an integer: weaker, $q \ge 1$; $k \ge 0$:
same; $\mathcal{N}$ as in Proposition~\ref{prop:creation-grid}: weaker, any set of indices that
contains the $2d$ axis neighbours.)}
\end{proposition}

% CHANGED (wave 3b, 2026-11-15): the proof follows the machine-checked route. It adds the step that
% identifies a tensor power on hZ^d with the array of products (the sum and the transform of the
% tensors are read through it), and the point (pi, ..., pi) where the transform of kappa_t is
% below 1. In (2) the frame mask is infinitely divisible by the criterion of prop:exact-masks(1),
% where the old proof cited prop:grid-pyramid-not-isotropic(1), and the binomial tensor mask is not,
% by prop:no-bounded-step, where the old proof cited the last sentence of prop:exact-masks(1).
% In (3) the old proof also recalled the nonnegative transform, which Proposition creation-grid no
% longer requires.
% CHANGED (fifth blind pass, 2026-11-15, flags 2-5): the kernel of the tensor step is c (k is the
% level); kappa_t is read on hZ^d; the transform of the Laplace mask is cited from
% eq:sampled-laplace-transfer, not from a display in a proof; the frame mask sentence is reworded.
\begin{proof}
\leanok
\uses{prop:creation-grid, lem:records-creation, prop:discrete-helmholtz,
prop:resolvent-masks-exact, prop:grid-pyramid-not-isotropic, prop:no-bounded-step}
\emph{Tensor powers.} Let $c \ge 0$ be summable on $\ZZ$, $\lambda$ the law on $h\ZZ$ with the
weight $c(n)$ at $hn$, and $K(n) := \prod_ic(n_i)$ on $\ZZ^d$. The law on $h\ZZ^d$ with the
weight $K(n)$ at $hn$ is $\lambda^{\otimes d}$. Both are finite measures on $\RR^d$, and a product
measure is determined by its values on boxes $s_1 \times \dots \times s_d$. On a box the first is
$\sum_{n \in \ZZ^d}\prod_ic(n_i)\mathbf{1}[hn_i \in s_i] = \prod_i\sum_{m \in \ZZ}
c(m)\mathbf{1}[hm \in s_i]$, a finite product of sums of nonnegative terms, which is the value of
the second. When $\lambda$ is a probability, so is $\lambda^{\otimes d}$, so $\sum_nK(n) = 1$, and
the transform of $K$ is that of $\lambda^{\otimes d}$, the product $\prod_i\widehat{\lambda}(\omega_i)$.
If $c$ is symmetric, so is $K$, and $K(0) = c(0)^d$.

% CHANGED (2026-11-23, review round 3 batch C, item 6): (1) no longer shows the transforms
% nonnegative, and the transform of beta moves to (2), which reads it; in (2) the frame mask is the
% compound Poisson law of its Levy measure nu_w, from its data in the proof of
% prop:grid-pyramid-not-isotropic(3), in place of its exact grid pyramid and the criterion of
% prop:exact-masks(1).
(1) By Proposition~\ref{prop:discrete-helmholtz}(1), $\kappa_t$, a law on $\ZZ^d$ read on $h\ZZ^d$
through $n \mapsto hn$ (so that its variable is $\vartheta = h\omega$), is a probability law with the
transform
\[
\widehat{\kappa_t}(\omega) = \Bigl(1 + \frac{t^2}{h^2}\sum_{i=1}^d(2 - 2\cos h\omega_i)\Bigr)^{-1},
\]
which that proposition writes through its frequency symbol $D_h$. It is symmetric, as a mixture of
the powers of the symmetric step $\sigma$. It is not $\delta_0$: at $h\omega = (\pi, \dots, \pi)$
its transform is $(1 + 4dt^2/h^2)^{-1} < 1$, while that of $\delta_0$ is $1$. The law
$\lambda^{\mathrm L}_a$ is symmetric and has mass $\tfrac{1-a}{1+a}(1 + \tfrac{2a}{1-a}) = 1$. By
the step on tensor powers, its tensor power is a symmetric probability, and its weight at the origin
is $(\tfrac{1-a}{1+a})^d < 1$. The mask $\beta$ is symmetric with mass $16/16 = 1$, so
$\beta^{\otimes d}$ is a symmetric probability, and its weight at the origin is $(6/16)^d < 1$.

(2) $\kappa_t$ is the mask (\ref{eq:resolvent-mask}) for the simple random walk step, by
Proposition~\ref{prop:resolvent-masks-exact}(4), so it is infinitely divisible by
Proposition~\ref{prop:resolvent-masks-exact}(1). The frame mask $w := (\lambda^{\mathrm L}_a)^{\otimes d}$, read through the step on tensor powers as
the array of products, has the transform
$\widehat{w}(\omega) = \exp\bigl(-\int(1 - \cos\langle\omega, x\rangle)\,\nu_w(dx)\bigr)$ with a
finite symmetric measure $\nu_w$ without atom at $0$, computed in the proof of
Proposition~\ref{prop:grid-pyramid-not-isotropic}(3) (``The mask in $\RR^d$''). The compound
Poisson law $e^{-\nu_w(\RR^d)}\sum_{k \ge 0}\nu_w^{*k}/k!$ has the transform
$\exp\int(e^{i\langle\omega, x\rangle} - 1)\,\nu_w(dx) = \widehat{w}(\omega)$, the sine part
integrating to $0$, so it is $w$, by the uniqueness of the transform. It is infinitely divisible,
the $n$-th convolution power of the compound Poisson law of $\nu_w/n$ for every $n$. The law
$\beta^{\otimes d}$ is carried by the ball of radius $2h\sqrt d$, each factor being carried by
$[-2h, 2h]$. It is no point mass: by
$\cos^4\theta = \tfrac18(3 + 4\cos 2\theta + \cos 4\theta)$,
\[
\widehat{\beta}(\omega) = \tfrac{1}{16}\bigl(6 + 8\cos h\omega + 2\cos 2h\omega\bigr) =
\cos^4(h\omega/2),
\]
and the transform $\prod_i\widehat{\beta}(\omega_i)$ of $\beta^{\otimes d}$ vanishes at
$\omega = (\pi/h)e_1$, where that of a point mass has modulus $1$. By
Proposition~\ref{prop:no-bounded-step} it is not infinitely divisible.

(3) By (1), $w$ is a symmetric probability law on $h\ZZ^d$ other than $\delta_0$. By (2) it is
infinitely divisible, and $q$ is an integer, so Proposition~\ref{prop:creation-grid}(3) makes every
bounded array on $q^kh\ZZ^d$ the level-$k$ record of a bounded input on $h\ZZ^d$.
Lemma~\ref{lem:records-creation} gives the input.
\end{proof}

So on the grids of the five-point resolvent and of the frame mask the count of local maxima fails
to be monotone from rung to rung, before and after decimation. The binomial tensor mask is not
infinitely divisible and its transform vanishes, so Proposition~\ref{prop:creation-grid}(3), whose
proof inverts the mask, does not apply to it. Decimation helps here: it is enough to invert the
mask's subsample on the coarse grid, and that subsample has no zero.

% CHANGED: the binomial records question closed, the author's approval of 2026-11-12. The draft
% marked it open (the paragraph after Proposition 4.7, line 646, and section 8, problem 1); new
% label prop:binomial-records, proved by inverting the subsampled mask.
\begin{proposition}[the records of the binomial masks]
\label{prop:binomial-records}
\lean{Pyramids.binomial_records_1, Pyramids.binomial_records_2}
% CHANGED (wave 3b, 2026-11-15): both declarations are machine-checked, so statement and proof are
% marked proved.
\leanok
\uses{prop:creation-grid, prop:creation-grid-masks, def:level-records, lem:records-creation}
Let $d \ge 2$, $h > 0$, let $q \ge 2$ be an integer, let $\beta$ be the binomial mask of
Proposition~\ref{prop:creation-grid-masks} and $w := \beta^{\otimes d}$ on $h\ZZ^d$.
\begin{enumerate}
\item For every integer $k \ge 0$, every bounded array on $q^kh\ZZ^d$ is the level-$k$ record
(Definition~\ref{def:level-records}) of a bounded input on $h\ZZ^d$.
\item Let $k \ge 0$ be an integer and $\mathcal{N}$ a neighbourhood as in
Proposition~\ref{prop:creation-grid}. There is a bounded input on $h\ZZ^d$ whose level-$k$ record
has exactly one $q^k\mathcal{N}$-maximum, while the smoothing of that record by $D_{q^k}w$ has at
least two $q^k\mathcal{N}$-maxima, and the level-$(k+1)$ record has at least two
$q^{k+1}\mathcal{N}$-maxima.
\end{enumerate}
% CHANGED (2026-11-15, a strengthening)
% The declarations require less than the print, which keeps its setting. The old status line ended at
% "by the inversion of the subsampled mask.)".
% CHANGED (fidelity review, 2026-11-15, R26/R50/R69): the status line lists the declarations' hypotheses beside the node's (LINKAGE rule 6).
\statusT\quad\emph{(Cites nothing beyond the nodes it uses. Clause (1) replaces
Proposition~\ref{prop:creation-grid}(3), which needs infinite divisibility, by the inversion of the
subsampled mask. Machine-checked. The declarations require less than the print, which keeps its
setting: clause (1) holds for every $d \ge 0$, and in clause (2) $\mathcal{N}$ need only contain
the $2d$ axis neighbours.
\textbf{Hypotheses} (the node's, each with what its declarations require): $d \ge 2$: weaker in
\texttt{binomial\_records\_1}, which requires no bound, same in \texttt{binomial\_records\_2}; $h > 0$:
absent, the levels are read in their own indices $\ZZ^d$; $q \ge 2$ an integer: same; $k \ge 0$:
same; in (1), the array bounded: same. (2): $\mathcal{N}$ as in
Proposition~\ref{prop:creation-grid}: weaker, any set of indices that contains the $2d$ axis
neighbours.)}
\end{proposition}

% CHANGED (wave 3b, 2026-11-15): the proof follows the machine-checked route. The inverse of the
% subsampled mask is explicit, 2 sqrt(2) rho^{|n|} with rho = 2 sqrt(2) - 3 at q = 2 and
% (16/6) delta_0 at q >= 3, in place of the Neumann series of epsilon = (delta_{-1} + delta_1)/6;
% S * S^{-1} = delta_0 is a finite sum over the box {-1,0,1}^d that factors into d line identities,
% in place of "the convolution of product kernels being taken factor by factor"; the step names
% the associativity of a summable, a summable and a bounded kernel. In (2) the old proof recalled
% the nonnegative transform, which Lemma records-creation no longer requires.
\begin{proof}
\leanok
\uses{prop:creation-grid-masks, def:level-records, lem:records-creation}
(1) \emph{The subsampled mask.} Let $s$ be the kernel on $\ZZ$ that reads $\beta$ on $qh\ZZ$:
$s(m) := \beta(\{qmh\})$. The mask $\beta$ charges only $0, \pm h, \pm 2h$, so $s$ is carried by
$\{-1, 0, 1\}$: at $q = 2$ it is $s = \tfrac{1}{16}(\delta_{-1} + 6\delta_0 + \delta_1)$, and at
$q \ge 3$ it is $s = \tfrac{6}{16}\delta_0$. Its inverse on $\ZZ$ is explicit. At $q \ge 3$ put
$s^{-1} := \tfrac{16}{6}\delta_0$. At $q = 2$ let $\rho := 2\sqrt2 - 3$, the root of
$\rho^2 + 6\rho + 1 = 0$ with $|\rho| < 1$, and put $s^{-1}(n) := 2\sqrt2\,\rho^{|n|}$, summable on
$\ZZ$. Then $\sum_j s(j)s^{-1}(n - j) = \tfrac{2\sqrt2}{16}\bigl(\rho^{|n-1|} + 6\rho^{|n|} +
\rho^{|n+1|}\bigr)$, which is $\tfrac{2\sqrt2}{16}\rho^{|n|-1}(1 + 6\rho + \rho^2) = 0$ at $n \ne 0$
and $\tfrac{2\sqrt2}{16}(6 + 2\rho) = \tfrac{2\sqrt2 \cdot 4\sqrt2}{16} = 1$ at $n = 0$. So
$s * s^{-1} = \delta_0$ in both cases.

Put $S := s^{\otimes d}$ and $S^{-1} := (s^{-1})^{\otimes d}$ on $\ZZ^d$. $S$ is carried by the box
$\{-1, 0, 1\}^d$, and $S^{-1}$ is summable, $|S^{-1}|$ being the tensor power of the summable
kernel $|s^{-1}|$. So $(S * S^{-1})(n) = \sum_{m \in \{-1,0,1\}^d}\prod_i s(m_i)s^{-1}(n_i - m_i) =
\prod_i\sum_{j \in \{-1,0,1\}}s(j)s^{-1}(n_i - j) = \prod_i\delta_0(n_i)$, a finite sum of products
being the product of the sums: $S * S^{-1} = \delta_0$. Since $w(\{qhm\}) = \prod_i\beta(\{qhm_i\})
= S(m)$ for $m \in \ZZ^d$, $S$ is the subsample of $w$ on $qh\ZZ^d$.

\emph{One step.} Let $j \ge 0$ and let $g$ be a bounded array on $q^{j+1}h\ZZ^d$. Write
$G(m) := g(q^{j+1}hm)$, $m \in \ZZ^d$, and $V := S^{-1} * G$, bounded by
$\sum|S^{-1}|\cdot\sup|G|$. Then $S * V = (S * S^{-1}) * G = G$, the convolution of two summable
kernels and a bounded one being associative. Define $u$ on $q^jh\ZZ^d$ by $u(q^{j+1}hm) := V(m)$ and
$u := 0$ off $q^{j+1}h\ZZ^d$; it is bounded. The mask of step $j$, $D_{q^j}w$, has the weight
$w(\{hn\})$ at $q^jhn$, and $u$ vanishes off $q^{j+1}h\ZZ^d$, so the sum defining the smoothing at
$q^{j+1}hm$ runs over $q^{j+1}h\ZZ^d$:
\begin{align*}
(D_{q^j}w * u)(q^{j+1}hm) &= \sum_{m' \in \ZZ^d}w\bigl(\{qh(m - m')\}\bigr)\,V(m') = (S * V)(m)
= g(q^{j+1}hm).
\end{align*}
So $g$ is the array that step $j$ keeps from the bounded $u$.

\emph{All levels.} By induction on $k$. The level-$0$ record of $f$ is $f$ itself. Let $k \ge 1$
and $g$ bounded on $q^kh\ZZ^d$. The step $j = k-1$ gives a bounded $u$ on $q^{k-1}h\ZZ^d$ whose
smoothing by $D_{q^{k-1}}w$, kept on $q^kh\ZZ^d$, is $g$. By induction $u$ is the level-$(k-1)$
record of a bounded input $f$ on $h\ZZ^d$. Then $g$ is the level-$k$ record of $f$, by the
recursion of Definition~\ref{def:level-records}. The argument reads no bound on $d$.

(2) By Proposition~\ref{prop:creation-grid-masks}(1), $w$ is a symmetric probability law on
$h\ZZ^d$ other than $\delta_0$, its weight at $0$ being $(6/16)^d < 1$. By (1) every bounded array
on $q^kh\ZZ^d$ is the level-$k$ record of a bounded input, and Lemma~\ref{lem:records-creation}
gives the input.
\end{proof}

So the count of local maxima fails to be monotone from rung to rung, before and after decimation,
for the binomial tensor masks as well, at every integer $q \ge 2$. Nothing is claimed on records at
a non-integer $q$. Proposition~\ref{prop:grid-implicit-euler-principle} bounds a value at an
$\mathcal{N}$-maximum and counts nothing, so it does not conflict with this.

What remains is to say what the pyramid algorithm returns, and to identify the exact grid pyramids
on the line with those that are exact in the sense of Definition~\ref{def:grid-pyramid-line}, a
definition on operators.

% draft: Proposition 5.7 (section 5.3), lines 849-858.
% CHANGED: in (1), "admissible in the sense of def:grid-pyramid-line" reads "exact in the sense of
% def:grid-pyramid-line", the word that definition now uses, and the clause says that it is what
% makes the two definitions agree at d = 1 (second blind review, flag 1; first review of chapter 2).
% CHANGED: in (2), the array returned is named the level-n record of def:level-records (flag 2).
\begin{proposition}[the line, and the arrays]
\label{prop:line-and-arrays}
\lean{Pyramids.line_and_arrays_1, Pyramids.line_and_arrays_2}
\leanok
% CHANGED (wave 3b, 2026-11-15): the uses follow the checked proof, which does not read A30.
\uses{def:exact-grid-pyramid, def:grid-pyramid-line, def:exact-senses, def:level-records}
\begin{enumerate}
\item \emph{(The line.)} In $d = 1$ a grid pyramid is exact in the sense of
Definition~\ref{def:exact-grid-pyramid}, on laws, if and only if it is exact in the sense of
Definition~\ref{def:grid-pyramid-line}, on operators. So the two definitions name the same
pyramids at $d = 1$, and this equivalence is what identifies them.
% CHANGED (2026-11-15, a strengthening)
% The record formula is stated for every integer q >= 1, which line_and_arrays_2 proves (it takes
% any natural q); exactness, defined for q > 1, keeps q >= 2. The old wording: "Let $q \ge 2$ be an
% integer, ..." and "If the grid pyramid of $(w, q)$ is exact, this is".
% CHANGED (2026-11-23, review round 3, s6-7.S19): wording.
\item \emph{(The arrays.)} Let $q \ge 1$ be an integer ($q = 1$ is the computation without decimation), $w$ a probability law on $\Gamma$ and $f$
bounded on $\Gamma$. The algorithm that, for $k = 0, \dots, n-1$, smooths by $D_{q^k}w$ on
$q^k\Gamma$ and keeps the values on $q^{k+1}\Gamma$ returns the level-$n$ record of $f$
(Definition~\ref{def:level-records}), which is
$\bigl(D_{q^{n-1}}w * \dots * D_qw * w * f\bigr)$ restricted to $q^n\Gamma$. If $q \ge 2$ and the
grid pyramid of $(w, q)$ is exact, this is the ladder increment $\mu_{1,q^n}$ applied to $f$ and read on the coarse
grid, so the computation is exact in sense (i) for that increment, in every dimension.
\end{enumerate}
% CHANGED (2026-11-15, a strengthening): proved from fewer cited facts. The status line was
% [A] on ledger entry A30, A30 (iso v0.1, Lem. 3.3(1)-(2), p. 11) spent in clause (1) in both
% directions; the checked proof reaches every cited fact through proved nodes and reads no A30.
% The Lean of clause (2) also does not take q >= 2, which its proof does not read.
% CHANGED (fidelity review, 2026-11-15, R26/R50/R69): the status line lists the declarations' hypotheses beside the node's (LINKAGE rule 6).
\statusT\quad\emph{(No ledger entry: every cited fact is carried by the nodes used. In clause (1),
the passage from operators to laws is the reading of the operator families on the line through
their kernels (Lemma~\ref{lem:line-axioms-kernels}, with Lemma~\ref{lem:family-dilations}(3) for
the gauge); the passage from laws to operators is the agreement of the two criteria,
Proposition~\ref{prop:exact-masks}(1) on the lattice $h\ZZ$ and
Proposition~\ref{prop:grid-pyramids-infinitely-divisible}(2), each of which carries its own cited
facts through the nodes it uses (on the checked route A39, A41, A42, A48 and A49). Clause (2) cites nothing. Proved
here: the lift to $\RR^1$ and back, the record formula and the ladder product. The Lean of
clause (2) reads no bound on the integer $q$, and the print states the record formula for every
$q \ge 1$ (2026-11-15, a strengthening); its second sentence does not read the exactness
of the pyramid, only a family carrying the ladder. \textbf{Not claimed}: anything about the arrays
when the grid pyramid is not exact, beyond the formula.
\textbf{Hypotheses} (the node's, each with what its declarations require): (1): $d = 1$: same, the line
typed as $\RR$, not as $\RR^1$; a full-rank lattice of the line: same, as $h\ZZ$ with $h > 0$; $w$ a
symmetric probability law on it, $w \ne \delta_0$: same (\texttt{IsGridMask}); $q > 1$: same. (2):
$q \ge 1$ an integer: weaker, any natural number; $\Gamma$ a full-rank lattice: same, a
\texttt{Submodule} over $\ZZ$ with \texttt{DiscreteTopology} and \texttt{IsZLattice}, with a basis
$b$ as a binder, which every such lattice has; $w$ a probability law on $\Gamma$: same; $f$ bounded:
same, as an array in the indices of $b$; in the second sentence, $q \ge 2$: absent, implied by the
exactness required; the pyramid exact: same, with every family carrying the ladder.)}
\end{proposition}

\begin{proof}
\leanok
% CHANGED (wave 3b, 2026-11-15)
% The proof of (1) follows the checked route: both directions through proved nodes, with no
% reading of iso v0.1, Lem. 3.3 (A30). The old route read each operator as the convolution with a
% unique law and turned convergence of transforms into L1-continuity, by that lemma.
% CHANGED (fourth blind pass, 2026-11-15): the covariance of the kernels is
% lem:line-axioms-kernels(3), not lem:family-dilations(3), which leaves the proof's citations
% (it stays in the uses for the gauge of the first paragraph); (b) from (I2) by Levy's continuity
% theorem; the converse applies prop:exact-masks(1) at d = 1 directly, the line being R^1, in place
% of the lift through x -> x e_1, which the print does not need.
\uses{prop:exact-masks, prop:grid-pyramids-infinitely-divisible, lem:line-axioms-kernels,
lem:family-dilations, def:q-covariant-line, prop:equivalent-weights}
(1) In $d = 1$ a full-rank lattice is $h\ZZ$ for some $h > 0$. Definition~\ref{def:grid-pyramid-line}
requires a $q^{\ZZ}$-covariant family of operators (Definition~\ref{def:q-covariant-line}) in the
canonical discrete gauge $S_qt = qt$ of Lemma~\ref{lem:family-dilations}(3), whose kernels have the
steps $\mu_{q^n,q^{n+1}} = T_{q^n}w$. On the line $T_c = D_c$.

\emph{Exact on operators implies exact on laws.} Let $(\Phi_{s,t})$ be a family as in
Definition~\ref{def:grid-pyramid-line}, with kernels $\mu_{s,t}$. Read through its kernels
(Lemma~\ref{lem:line-axioms-kernels}(1)), it is a symmetric family of laws in dimension $1$. Its
cascade (I1) is (a) of Definition~\ref{def:exact-grid-pyramid}, convolution being commutative, and
the continuity (I2) of its transforms gives the weak continuity (b) by \Levy's continuity theorem.
Its covariance in the canonical discrete gauge is $T_q\mu_{s,t} = \mu_{qs,qt}$
(Lemma~\ref{lem:line-axioms-kernels}(3) with $\lambda = q$ and $S = S_q$), which is (c), and the
step $n = 0$ of the ladder is $\mu_{1,q} = T_1w = w$, which is (d). So the kernels carry the grid
pyramid in the sense of Definition~\ref{def:exact-grid-pyramid}.

\emph{Exact on laws implies exact on operators.} Let $(\mu_{s,t})$ carry the grid pyramid in the
sense of Definition~\ref{def:exact-grid-pyramid}. Here $\RR$ is $\RR^d$ at $d = 1$, and $w$ is a
symmetric probability law other than $\delta_0$, carried by the lattice $h\ZZ$. By the criterion
of Proposition~\ref{prop:exact-masks}(1), at $d = 1$, $w$ is infinitely divisible with
$\int\log_+|x|\,w(dx) < \infty$. That is the second criterion of
Proposition~\ref{prop:grid-pyramids-infinitely-divisible}(2), so the grid pyramid of $(w, q)$ is
exact in the sense of Definition~\ref{def:grid-pyramid-line}: a $q^{\ZZ}$-covariant family of
operators in the canonical discrete gauge has the steps of the ladder.

(2) As $q$ is an integer, $q\Gamma \subseteq \Gamma$, so the grids $q^k\Gamma$ are nested. The law
$D_{q^k}w$ is carried by $q^k\Gamma$, so for $x \in q^k\Gamma$
\[
(D_{q^k}w * g)(x) = \sum_{y \in q^k\Gamma}(D_{q^k}w)(\{y\})\,g(x - y)
\]
uses only the values of $g$ on $q^k\Gamma$: decimation discards nothing the later steps read. Let
$g_0 := f$ and let $g_{k+1}$ be $D_{q^k}w * g_k$ restricted to $q^{k+1}\Gamma$, the array the
algorithm holds after step $k$. These are the records of $f$ (Definition~\ref{def:level-records}). Put $F_k := D_{q^{k-1}}w * \dots * w * f$ on $\Gamma$, with
$F_0 := f$. If $g_k$ is $F_k$ restricted to $q^k\Gamma$, then for $x \in q^{k+1}\Gamma$ the display
gives $g_{k+1}(x) = (D_{q^k}w * F_k)(x) = F_{k+1}(x)$. By induction the algorithm returns $F_n$
restricted to $q^n\Gamma$. If the pyramid is exact, (a), (c) and (d) of
Definition~\ref{def:exact-grid-pyramid} give
\[
\mu_{1,q^n} = \mu_{q^{n-1},q^n} * \dots * \mu_{q,q^2} * \mu_{1,q} = D_{q^{n-1}}w * \dots * D_qw * w,
\]
so the returned array is $\mu_{1,q^n} * f$ restricted to $q^n\Gamma$. It equals the lattice law it
stands for applied to $f$, which is sense (i) of Definition~\ref{def:exact-senses}.
\end{proof}

% CHANGED: "which is a continuum kernel (Proposition 5.4(1))" dropped (blind review, flag 8):
% clause (1) gives the Levy measure of mu_{q^n}, not the absence of atoms; the kernel is named by
% its product formula instead.
Clause (2) says what a pyramid computes. The input is read as the record at scale $1$, and level
$n$ is the increment of the ladder from there. It is not the kernel $\mu_{q^n} = D_{q^n}\mu_1$ from
scale zero, the infinite product (\ref{eq:knot-product}) dilated. For the sampled Laplace mask on
the line the distance between the computed arrays and the samples of the \Matern{} levels is
therefore the error of Proposition~\ref{prop:sampled-laplace-exact}(2), and not the aliased mass of
Proposition~\ref{prop:sampled-laplace-exact}(3).

% draft: Remark 5.8 (section 5.3), lines 860-865. Its clause (1) asserts and argues, so its
% mathematics is the node prop:binomial-knot-relation (new label); clauses (2) and (4) stay prose,
% each sentence citing a node; clause (3) is in the section's remark on the line.
The definition requires continuity in scale, not for the knot relation. The relation
$D_q\mu_1 = \mu_1 * w$ alone is satisfied by pyramids that
Proposition~\ref{prop:grid-pyramid-not-isotropic} does not cover.

% draft: Remark 5.8(1), line 862, its first three sentences.
% CHANGED: the binomial mask is written as a measure on hZ, and "not infinitely divisible" is
% asserted of the tensor power on hZ^d as well, through its marginal (the draft cites Proposition
% 2.16, a statement on the line).
\begin{proposition}[the binomial mask satisfies the knot relation]
\label{prop:binomial-knot-relation}
\lean{Pyramids.binomial_knot_relation, Pyramids.binomial_knot_relation_tensor}
\leanok
% CHANGED (wave 3g, 2026-11-15): the uses follow the checked proof, which reads the line from
% prop:equivalent-weights(2) and the infinite divisibility of the mask from lem:family-structure.
% CHANGED (fourth blind pass, 2026-11-15): lem:family-structure leaves the uses, the proof citing
% prop:no-bounded-step instead.
\uses{def:exact-grid-pyramid, prop:equivalent-weights, prop:classical-kernels-no-exact-cascade,
prop:no-bounded-step}
Let $\beta := \tfrac{1}{16}(\delta_{-2h} + 4\delta_{-h} + 6\delta_0 + 4\delta_h + \delta_{2h})$ be the
binomial mask on $h\ZZ$ and $q = 2$. The product (\ref{eq:knot-product}) converges,
\[
\prod_{m \ge 1}\widehat{\beta}\bigl(2^{-m}\omega\bigr) = \prod_{m \ge 1}\cos^4\bigl(2^{-m-1}h\omega\bigr)
= \Bigl(\frac{\sin(h\omega/2)}{h\omega/2}\Bigr)^4 ,
\]
to the transform of the cubic B-spline $\mu_1$, the law of the sum of four independent variables
uniform on $[-h/2, h/2]$. It satisfies $D_2\mu_1 = \mu_1 * \beta$. The same holds for the tensor
power $\beta^{\otimes d}$ on $h\ZZ^d$, with the tensor power $\mu_1^{\otimes d}$. These limits have
bounded support and are not infinitely divisible, and the grid pyramids of $(\beta^{\otimes d}, 2)$,
$d \ge 1$, are not exact.
% CHANGED (wave 3g, 2026-11-15): both declarations are machine-checked, so statement and proof are
% marked proved; the status line names the nodes that carry the cited facts.
% CHANGED (fourth blind pass, 2026-11-15): both cited facts now enter through
% prop:no-bounded-step.
% CHANGED (fidelity review, 2026-11-15, R26/R50/R69): the status line lists the declarations' hypotheses beside the node's (LINKAGE rule 6).
\statusT\quad\emph{(Cites nothing of its own. Khintchine's theorem (A41) and Sato's Cor.~24.4
(A46) enter through Proposition~\ref{prop:no-bounded-step}, which excludes infinite divisibility
for a law of bounded support that is not a point mass and makes no such law a step of a family with
(I1) and (I2); on the line the first is also
Proposition~\ref{prop:classical-kernels-no-exact-cascade}. Every clause is machine-checked.
\textbf{Hypotheses} (the node's, each with what its declarations require): $h > 0$: same; $q = 2$: same;
the line: same, typed as $\RR$ in \texttt{binomial\_knot\_relation}, not as $\RR^1$; $d \ge 1$ for
the tensor power: same, in \texttt{binomial\_knot\_relation\_tensor}.)}
\end{proposition}

% CHANGED (wave 3g, 2026-11-15)
% The proof follows the machine-checked route. The line is prop:equivalent-weights(2); the tensor
% power mu_1^{(x)d} and the mask beta^{(x)d} are excluded from infinite divisibility by
% prop:no-bounded-step directly (bounded support, and a transform with a zero, so no point mass), in
% place of the marginal argument; and an exact grid pyramid is excluded because its mask mu_{1,2} is
% infinitely divisible by the first step of lem:family-structure, in place of
% prop:grid-pyramid-not-isotropic(1). Statement unchanged.
% CHANGED (fourth blind pass, 2026-11-15): the exclusion of a carrying family cites the second
% sentence of prop:no-bounded-step, a statement, in place of a step inside the proof of
% lem:family-structure, which leaves the uses; the sentence on the image of the line, empty in
% print, is dropped; the weak convergence of the partial convolutions is said.
% CHANGED (the theory paper, 2026-11-16): the closing sentence on checking and cited facts shortened; no step changed.
\begin{proof}
\leanok
\uses{prop:equivalent-weights, prop:classical-kernels-no-exact-cascade, prop:no-bounded-step}
\emph{The line.} The five point masses of $\beta$ give $\widehat{\beta}(\omega) = \cos^4(h\omega/2)$.
By Proposition~\ref{prop:equivalent-weights}(2) the partial products
$\prod_{m=1}^M\widehat{\beta}(2^{-m}\omega)$ converge at every $\omega$ to
$\widehat{\mu_1}(\omega) = \operatorname{sinc}(h\omega/2)^4$, the transform of the cubic B-spline
$\mu_1$, where $\operatorname{sinc}\theta := \sin\theta/\theta$, and $D_2\mu_1 = \mu_1 * \beta$. The
partial convolutions $D_{2^{-1}}\beta * \dots * D_{2^{-M}}\beta$ have the partial products as
transforms, so by \Levy's continuity theorem they converge weakly to $\mu_1$.

\emph{The tensor powers.} The transform of a tensor power is the product over the axes of the
transforms of its factor. So the partial product of $\beta^{\otimes d}$ at $\omega$ is
$\prod_{i=1}^d$ of the partial products of $\beta$ at $\omega_i$, and as $M \to \infty$ this finite
product of convergent sequences tends to $\prod_i\widehat{\mu_1}(\omega_i)$, the transform of
$\mu_1^{\otimes d}$; the partial convolutions have the partial products as transforms, so they
converge weakly to $\mu_1^{\otimes d}$, again by \Levy's continuity theorem. The relation
$D_2\mu_1^{\otimes d} = \mu_1^{\otimes d} * \beta^{\otimes d}$ is the relation of the line in each
factor of the transform.

\emph{Bounded support, and not infinitely divisible.} The law $\mu_1$ is carried by $[-2h, 2h]$, a
sum of four variables with values in $[-h/2, h/2]$, and so is $\beta$. If a law $\nu$ on $\RR$ is
carried by $[-r, r]$, then $\nu^{\otimes d}$ is carried by the ball of radius $\sqrt d\,r$: a point
with $|x_i| \le r$ for every $i$ has $|x|^2 = \sum_ix_i^2 \le dr^2$, and the set where some
$|x_i| > r$ is a finite union of sets of product measure $0$. A tensor power $\nu^{\otimes d}$,
$d \ge 1$, is no point mass when $\widehat\nu(t) = 0$ for some $t$: its transform vanishes at
$te_1$, and the transform of a point mass has modulus $1$. Here $\widehat{\mu_1}(2\pi/h) =
\operatorname{sinc}(\pi)^4 = 0$ and $\widehat{\beta}(\pi/h) = \cos^4(\pi/2) = 0$. So by
Proposition~\ref{prop:no-bounded-step} neither $\mu_1^{\otimes d}$ nor $\beta^{\otimes d}$ is
infinitely divisible; on the line this is Proposition~\ref{prop:classical-kernels-no-exact-cascade}
for $\mu_1$ and for $\beta$.

\emph{Not exact.} Let a family $(\mu_{s,t})$ carry the grid pyramid of $(\beta^{\otimes d}, 2)$.
By (a) it has (I1), and by (b) it has (I2), weak convergence giving the convergence of the
transforms at every $\omega$. Then $\mu_{1,2} = \beta^{\otimes d}$ is a step of a family of
probability measures on $\RR^d$ with (I1) and (I2), and a law of bounded support that is not a
point mass. The second sentence of Proposition~\ref{prop:no-bounded-step} excludes this. So no family
carries the pyramid, for every $d \ge 1$.

Every step is machine-checked, and the cited facts, \ledger{A41} and \ledger{A46}, enter only
through Proposition~\ref{prop:no-bounded-step}.
\end{proof}

% CHANGED (review round 2, batch C, 2026-11-23): the last sentence, "Whether some positive finite
% mask has a rotation-invariant refinable limit, a question about refinable functions and not about
% hemigroups, is not addressed here.", is replaced by prop:no-isotropic-refinable-finite-mask, which
% answers it (open question (4) of the theory paper's section 7.2).
Here the carrier step of Proposition~\ref{prop:grid-pyramid-not-isotropic} has nothing to act on:
condition (b) is what makes the mask infinitely divisible.
% CHANGED (2026-11-23, review round 3 batch C, item 7): "For a mask of finite support" reads "For a
% mask with an exponential moment, finite support included", the reach of the widened node.
For a mask with an exponential moment, finite support included, the knot
relation alone excludes a rotation-invariant law in $d \ge 2$, by the analyticity of the mask's
transform.

% CHANGED (review round 2, batch C, 2026-11-23): new node. The mask is any probability law carried
% by a finite set, symmetric or not, on a lattice or not, and the law is any solution of the knot
% relation, not only the limit of the product (eq:knot-product).
% CHANGED (2026-11-23, review round 3 batch C, item 7, s1-3.S13 and s6-7.S7): a strengthening. The
% statement was the finite case alone: "Let $d \ge 2$ and $q > 1$, let $\mu$ be a probability law
% on $\RR^d$ with $D_R\mu = \mu$ for every $R \in \mathrm{SO}(d)$, and let $w$ be a probability law
% on $\RR^d$ carried by a finite set. If $D_q\mu = \mu * w$, then $w = \delta_0$ and
% $\mu = \delta_0$. In particular, for a mask $w$ of finite support in the sense of
% Definition~\ref{def:exact-grid-pyramid}, no rotation-invariant probability law satisfies the knot
% relation $D_q\mu = \mu * w$." Clause (1) takes a countable carrier and an exponential moment,
% clause (2) finitely many complex weights. The decay cannot be dropped: a band-limited mask on
% hZ^d of infinite support has a rotation-invariant refinable law (review round 3, s6-7 2.2(a)).
% The title said "of a finite mask".
\begin{proposition}[no rotation-invariant refinable law of a mask with an exponential moment]
\label{prop:no-isotropic-refinable-finite-mask}
\lean{Pyramids.no_isotropic_refinable_mask_exp, Pyramids.no_isotropic_refinable_signed_finite_mask, Pyramids.no_isotropic_refinable_finite_mask}
\leanok
\uses{def:exact-grid-pyramid}
Let $d \ge 2$ and $q > 1$, and let $\mu$ be a probability law on $\RR^d$ with $D_R\mu = \mu$ for
every $R \in \mathrm{SO}(d)$.
\begin{enumerate}
\item Let $w$ be a probability law on $\RR^d$ carried by a countable set, with an exponential
moment: $\int e^{c|x|}\,w(dx) < \infty$ for some $c > 0$. If $D_q\mu = \mu * w$, then
$w = \delta_0$ and $\mu = \delta_0$.
\item Let $w = \sum_{x \in S}a_x\delta_x$ for a finite set $S \subset \RR^d$ and complex weights
$a_x$, a signed or complex measure. If $\widehat\mu(q\omega) = \widehat\mu(\omega)\,
\sum_{x \in S}a_xe^{i\langle x, \omega\rangle}$ for every $\omega$, as when $D_q\mu = \mu * w$,
then $w = \delta_0$ and $\mu = \delta_0$.
\end{enumerate}
In particular, for a mask $w$ in the sense of Definition~\ref{def:exact-grid-pyramid} with an
exponential moment, finite support included, no rotation-invariant probability law satisfies the knot
relation $D_q\mu = \mu * w$.
% CHANGED (2026-11-23, review round 3 batch C, item 7): the status names the three declarations and
% their readings.
% CHANGED (2026-11-23, review round 3 batch C, item 14): the band-limited example of "Not claimed"
% is now the node prop:band-limited-refinable-mask, which the status line names.
\statusT\quad\emph{(Cites nothing. Machine-checked, in Lean core; the identity theorem for analytic
functions, the analyticity of the complex moment generating function in its strip, Dedekind's
lemma on the linear independence of characters and the uniqueness of the transform are Mathlib's.
Clause (1) is \texttt{no\_isotropic\_refinable\_mask\_exp}, clause (2)
\texttt{no\_isotropic\_refinable\_signed\_finite\_mask}, and the finite case of (1), the statement of
the second review round, \texttt{no\_isotropic\_refinable\_finite\_mask}, now a corollary.
\textbf{Not claimed}: the statement without a decay condition, which is false (a band-limited mask
on $h\ZZ^d$ of infinite support has a rotation-invariant refinable law,
Proposition~\ref{prop:band-limited-refinable-mask}); a signed
mask of infinite support. \textbf{Hypotheses} (the node's, each with what its declarations
require): $d \ge 2$: same; $q > 1$: same; $\mu$ a probability law with $D_R\mu = \mu$ for
$R \in \mathrm{SO}(d)$: same, \texttt{IsSORotationInvariant}. (1): $w$ a probability law carried by
a countable set: same, $w$ of mass $0$ off a countable set $C$; the exponential moment: same, the
integrability of $e^{c|x|}$ against $w$; the knot relation: same, the image of $\mu$ under
$x \mapsto qx$ equal to $\mu * w$. (2): the finite set and the weights: same, a finitely supported
function $a$ from $\RR^d$ to $\CC$; the relation on transforms: same, which $D_q\mu = \mu * w$
implies, the transform of a convolution being the product of the transforms; the conclusion
$w = \delta_0$: $a$ is the indicator of $\{0\}$. The symmetry of $w$, a lattice, the mass of $w$
and $w \ne \delta_0$ are not required: the knot relation at $\omega = 0$ gives mass $1$.)}
\end{proposition}

% CHANGED (2026-11-23, review round 3 batch C, item 7): the proof covers both clauses. The transform
% is used along lines through 0 only, where it is analytic in one variable; the atoms are read from
% the countable carrier in (1) and from the uniqueness of the weights of an exponential sum in (2).
\begin{proof}
\leanok
Write $\psi$ for the transform of $w$: $\psi := \widehat{w}$ in (1), and
$\psi(\omega) := \sum_{x \in S}a_xe^{i\langle x, \omega\rangle}$ in (2). In (1) the knot relation
gives $\widehat\mu(q\omega) = \widehat\mu(\omega)\psi(\omega)$, the transform of a convolution being
the product of the transforms; in (2) this is the hypothesis. For $R \in \mathrm{SO}(d)$ the
invariance gives $\widehat\mu(R^{\tp}\omega) = \widehat\mu(\omega)$ for every $\omega$.

\emph{Analytic along lines.} Fix $v \in \RR^d$. In (2), $s \mapsto \psi(sv)$ is a finite sum of
exponentials of linear functions of $s$, so it is real-analytic on $\RR$. In (1), let
$c' := c/(|v| + 1)$. For real $|\tau| < c'$, $e^{\tau\langle x, v\rangle} \le
e^{|\tau|\,|v|\,|x|} \le e^{c|x|}$, so $e^{\tau\langle x, v\rangle}$ is integrable against $w$. The
function $F(z) := \int e^{z\langle x, v\rangle}\,w(dx)$ is therefore defined and holomorphic on the
strip $|\operatorname{Re}z| < c'$, which contains the imaginary axis, and $\psi(sv) = F(is)$. So
$s \mapsto \psi(sv)$ is real-analytic on $\RR$.

\emph{The transform of the mask is rotation invariant.} Fix $R \in \mathrm{SO}(d)$. The transform
$\widehat\mu$ is continuous with $\widehat\mu(0) = 1$, so it has no zero on a neighbourhood $U$ of
$0$. For $\omega \in U$,
\[
\psi(R^{\tp}\omega)\,\widehat\mu(\omega) = \psi(R^{\tp}\omega)\,\widehat\mu(R^{\tp}\omega)
= \widehat\mu(qR^{\tp}\omega) = \widehat\mu(q\omega) = \psi(\omega)\,\widehat\mu(\omega),
\]
so $\psi(R^{\tp}\omega) = \psi(\omega)$ on $U$. Now fix $\omega$. The functions
$s \mapsto \psi(sR^{\tp}\omega)$ and $s \mapsto \psi(s\omega)$ are real-analytic on $\RR$, and they
agree for $s$ near $0$, where $s\omega \in U$. By the identity theorem they agree on $\RR$, and at
$s = 1$, $\psi(R^{\tp}\omega) = \psi(\omega)$.

\emph{The mask is $\delta_0$.} In (1), $\psi(R^{\tp}\cdot)$ is the transform of $D_Rw$, so
$D_Rw = w$ by the uniqueness of the transform, and as $R$ is injective, $w(\{Rx\}) = w(\{x\})$ for
every $x$. In (2), $\psi(R^{\tp}\omega) = \sum_{x \in S}a_xe^{i\langle Rx, \omega\rangle}$ is the
exponential sum with the weight $a_x$ at $Rx$. The characters $\omega \mapsto
e^{i\langle y, \omega\rangle}$, $y \in \RR^d$, are distinct homomorphisms from $\RR^d$ to the
multiplicative group of $\CC$: for $y \ne y'$ and $u := y - y'$, at $\omega = \pi u/|u|^2$ their
quotient is $e^{i\pi} = -1$. By Dedekind's lemma they are linearly independent over $\CC$, so the
weights of an exponential sum are unique, and $a_{Rx} = a_x$ for every $x$, with $a_y := 0$ for
$y \notin S$. In both clauses let $x \ne 0$ carry a nonzero weight, $w(\{x\}) > 0$ in (1) and
$a_x \ne 0$ in (2), and let $|y| = |x|$. Some $R \in \mathrm{SO}(d)$ has $Rx = y$, and the sphere
$\{|y| = |x|\}$ is uncountable, both as in the proof of
Proposition~\ref{prop:grid-pyramid-not-isotropic}(2). So every point of that sphere carries the
same nonzero weight, and lies in the countable carrier in (1), in the finite set $S$ in (2), which
is impossible. So $w$ has no weight off $0$. In (1), $w(\RR^d \setminus \{0\})$ is at most the mass
off the countable carrier $C$ plus $\sum_{x \in C,\,x \ne 0}w(\{x\})$, which is $0$, and
$w = \delta_0$. In (2), $\psi \equiv a_0$, and the relation at $\omega = 0$ gives
$\widehat\mu(0) = \widehat\mu(0)\,a_0$, so $a_0 = 1$ and $w = \delta_0$.

\emph{The law is $\delta_0$.} Now $\psi \equiv 1$, so $\widehat\mu(q\omega) =
\widehat\mu(\omega)$, and by induction $\widehat\mu(\omega) = \widehat\mu(q^{-n}\omega)$ for every
$n \ge 1$. As $n \to \infty$, $q^{-n}\omega \to 0$, and by continuity $\widehat\mu(\omega) =
\widehat\mu(0) = 1$. So $\widehat\mu \equiv 1$, the transform of $\delta_0$, and $\mu = \delta_0$.

The last sentence of the statement is clause (1): a mask is a probability law carried by a lattice,
a countable set, and a law carried by a finite set has every exponential moment. Every step is
machine-checked, with no cited fact.
\end{proof}

% CHANGED (2026-11-23, review round 3 batch C, item 14, decision D2): new prose and a new node.
Without a decay condition the knot relation alone allows isotropy, and so does exactness without
the continuity (b). A mask whose transform vanishes outside small balls about the points of the
dual grid has a rotation-invariant refinable law. Such a mask is not infinitely divisible, so no
family continuous in scale carries it, but a family that jumps at the knots does.

% CHANGED (2026-11-23, review round 3 batch C, item 14, decision D2 of
% records/pyramids/PLAN-review-response-theory-round3.md): new node, the band-limited example of
% review round 3 (s6-7 appendix 2.2(a), checked step by step by its verifier), with the verifier's
% family without (b). Proved in print only and not ready, as prop:expanding-grid-pyramid-not-isotropic
% is; its statement is typed in the skeleton. The density p is a hypothesis with its existence a
% clause, where the review fixed p := c |inverse transform of psi|^2: the proof reads only that p is
% a rotation-invariant Schwartz probability density with transform vanishing off the ball B_r.
\begin{proposition}[a band-limited mask with a rotation-invariant refinable law]
\label{prop:band-limited-refinable-mask}
\lean{Pyramids.band_limited_density_exists, Pyramids.band_limited_refinable_mask_1,
Pyramids.band_limited_refinable_mask_2, Pyramids.band_limited_refinable_mask_3}
\notready
\uses{def:exact-grid-pyramid, lem:steps-infinitely-divisible, prop:grid-pyramid-not-isotropic,
prop:no-isotropic-refinable-finite-mask, thm:cited-A40, thm:cited-A41}
Let $d \ge 1$, $q > 1$, $h > 0$ and $0 < r \le 2\pi/\bigl((q + 1)h\bigr)$. Let $p$ be a
rotation-invariant probability density on $\RR^d$ of the Schwartz class whose transform
$\varphi := \widehat p$ vanishes for $|\omega| \ge r$; such densities exist for every $r > 0$. Put
\[
  w := \sum_{n \in \ZZ^d}h^d\,p(hn)\,\delta_{hn} .
\]
\begin{enumerate}
\item \emph{(The mask.)} $w$ is a symmetric probability law on $h\ZZ^d$ with $w \ne \delta_0$, and
$\widehat w(\omega) = \sum_{k \in \ZZ^d}\varphi(\omega - 2\pi k/h)$. On $|\omega| < 2\pi/h - r$
the transform $\widehat w$ equals $\varphi$, so it vanishes for $r \le |\omega| < 2\pi/h - r$. The
mask has no exponential moment, so its support is infinite, and it is not infinitely divisible.
\item \emph{(A rotation-invariant refinable law.)} Let $Y_1, Y_2, \dots$ be independent with
density $p$. The series $\sum_{m \ge 1}q^{-m}Y_m$ converges almost surely, and its law $\mu$ is
rotation invariant, has a $C^\infty$ density and satisfies $D_q\mu = \mu * w$. Its transform is
$\prod_{m \ge 1}\varphi(q^{-m}\omega)$, the limit of the partial products
$\prod_{m=1}^M\widehat w(q^{-m}\omega)$ of (\ref{eq:knot-product}).
\item \emph{(The pyramid.)} The grid pyramid of $(w, q)$ on $h\ZZ^d$ is not exact. Put
$N(t) := \max\{n \in \ZZ : q^n \le t\}$ for $t > 0$, $\mu_{0,0} := \delta_0$,
$\mu_{0,t} := D_{q^{N(t)}}\mu$ for $t > 0$, and for $0 < s \le t$ let $\mu_{s,t}$ be the
convolution of the laws $D_{q^{n-1}}w$ over the integers $n$ with $s < q^n \le t$, and $\delta_0$
if there is none. This family has (a), (c) and (d) of Definition~\ref{def:exact-grid-pyramid}, and
every kernel $\mu_t$, $t > 0$,
is rotation invariant. It does not have (b): $\mu_t = D_{q^{n-1}}\mu$ for
$q^{n-1} \le t < q^n$, and $\mu_{q^n} = D_{q^n}\mu \ne D_{q^{n-1}}\mu$, at every knot $q^n$.
\end{enumerate}
\statusA\quad\ledger{A40}\quad\ledger{A41}\quad\emph{(\textbf{Assignment.} \ledger{A40} carries
that the transform of an infinitely divisible law has no zero, so that $w$ is not infinitely
divisible, in (1). \ledger{A41} carries, through Lemma~\ref{lem:steps-infinitely-divisible}, that a
step of a family with (a) and (b) is infinitely divisible, so that the pyramid is not exact, in (3).
Proved here, in print only: everything else, with textbook facts of analysis and probability that
are not ledger entries: the Poisson summation formula on $\ZZ^d$ for a Schwartz function,
Plancherel's identity, the convolution theorem, dominated convergence and the uniqueness of the
transform. \textbf{Not claimed}: a mask with an exponential moment, which
Proposition~\ref{prop:no-isotropic-refinable-finite-mask} excludes; a family continuous in scale
through this mask, which (3) excludes. \textbf{Left}, and why the node is not marked proved: the
proof is not formalized, by the author's decision D2 of the third review round's plan, since its
Lean needs the Poisson summation formula on $\ZZ^d$, which Mathlib states on $\RR$ only; its
statement is typed in the skeleton. Priced at about 1200 to 2000 lines: Poisson summation on
$\ZZ^d$ for Schwartz functions (about 500 to 800), the existence of the density through the inverse
transform of a radial bump (about 300 to 500), the random series, its density and the knot relation
(about 300 to 500), and the family of (3) (about 150). \textbf{Hypotheses} (the node's, each with
what its declarations require): $d \ge 1$, $q > 1$, $h > 0$, $0 < r \le 2\pi/((q+1)h)$: same, and
$r > 0$ alone for the existence; $p$: same (\texttt{IsBandLimitedDensity}), a Schwartz function with
$p \ge 0$, $\int p = 1$, $p(x) = p(y)$ for $|x| = |y|$, and the transform of the law $p(x)\,dx$ zero
for $|\omega| \ge r$; $w$: same (\texttt{sampledMask}). Conclusions: (1) as printed, the vanishing on
the shell left as the consequence of the equality with $\varphi$ that it is; (2) with $\mu$ typed by
its transform, the limit of the partial products of $\varphi$, which determines it, in place of
the law of the series; (3) with the family typed as one that exists with (a), (c), (d), rotation-invariant
kernels, kernels constant on $[q^{n-1}, q^n)$ and a jump at every knot, in place of the formula.)}
\end{proposition}

\begin{proof}
\uses{lem:steps-infinitely-divisible, prop:no-isotropic-refinable-finite-mask, thm:cited-A40,
thm:cited-A41}
The transform is $\widehat\rho(\omega) = \int e^{i\langle\omega, x\rangle}\rho(dx)$, also for a
function $\rho$.

\emph{Such densities exist.} Let $\psi \ne 0$ be a radial $C^\infty$ function on $\RR^d$ supported
in the closed ball of radius $r/2$, and let
$g(x) := (2\pi)^{-d}\int\psi(\omega)e^{-i\langle x, \omega\rangle}\,d\omega$, so that
$\widehat g = \psi$. As $\psi$ is real, even and radial, $g$ is real and radial, and it is a
Schwartz function. Then $g^2$ is a Schwartz function with transform $(2\pi)^{-d}\psi * \psi$, and
$\int g^2 = (2\pi)^{-d}\int\psi^2 > 0$ by Plancherel's identity. So $p := g^2/\int g^2$ is a
rotation-invariant probability density of the Schwartz class with
$\widehat p = \psi * \psi/\int\psi^2$. A point $\eta$ contributes to
$(\psi * \psi)(\omega) = \int\psi(\eta)\psi(\omega - \eta)\,d\eta$ only if $|\eta| \le r/2$ and
$|\omega - \eta| \le r/2$. For $|\omega| > r$ there is no such $\eta$, and for $|\omega| = r$ only
$\eta = \omega/2$, a null set. So $\widehat p$ vanishes for $|\omega| \ge r$.

(1) \emph{The transform.} The weights $h^dp(hn)$ are nonnegative and summable, $p$ being a Schwartz
function, so $w$ is a finite measure on $h\ZZ^d$. The Poisson summation formula
$\sum_{n \in \ZZ^d}f(n) = \sum_{k \in \ZZ^d}\int f(x)e^{-2\pi i\langle k, x\rangle}\,dx$, applied to
the Schwartz function $f(x) := h^dp(hx)e^{i\langle\omega, hx\rangle}$, with the substitution
$y = hx$, gives
\[
  \widehat w(\omega) = \sum_{n \in \ZZ^d}h^dp(hn)e^{i\langle\omega, hn\rangle}
  = \sum_{k \in \ZZ^d}\varphi(\omega - 2\pi k/h) .
\]
For $k \ne 0$, $|2\pi k/h| \ge 2\pi/h > r$, so at $\omega = 0$ the sum is $\varphi(0) = 1$, and $w$
is a probability law. It is symmetric, as $p$ is even. If $|\omega| < 2\pi/h - r$ and $k \ne 0$,
then $|\omega - 2\pi k/h| \ge 2\pi/h - |\omega| > r$, so $\widehat w(\omega) = \varphi(\omega)$,
which is $0$ for $|\omega| \ge r$. The shell $r \le |\omega| < 2\pi/h - r$ is not empty, as
$2r \le 4\pi/((q+1)h) < 2\pi/h$.

\emph{What the zeros exclude.} The transform of $\delta_0$ has no zero, so $w \ne \delta_0$, and an
infinitely divisible law has a transform without zero (\ledger{A40}), so $w$ is not infinitely
divisible. If $\int e^{c|x|}\,w(dx) < \infty$ for some $c > 0$, then $s \mapsto \widehat w(se_1)$
is real-analytic on $\RR$, as in the proof of
Proposition~\ref{prop:no-isotropic-refinable-finite-mask} (``Analytic along lines''). It vanishes on
the interval $r < s < 2\pi/h - r$, so by the identity theorem it vanishes everywhere, against
$\widehat w(0) = 1$. So $w$ has no exponential moment. A law carried by a finite set has every
exponential moment, so the support of $w$ is infinite.

(2) \emph{The law.} As $\int|x|\,p(x)\,dx < \infty$, $E\sum_mq^{-m}|Y_m| < \infty$, so the series
converges absolutely almost surely. By independence and dominated convergence,
$\widehat\mu(\omega) = \lim_M\prod_{m=1}^M\varphi(q^{-m}\omega)$. Each factor is radial, so
$\widehat\mu(R^{\tp}\omega) = \widehat\mu(\omega)$ for every rotation $R$, and $\mu$ is rotation
invariant. With $\mu'$ the law of $\sum_{m \ge 2}q^{-m}Y_m$, independent of $q^{-1}Y_1$, the law
$\mu$ has the density $x \mapsto \int q^dp\bigl(q(x - y)\bigr)\,\mu'(dy)$, which is $C^\infty$ by
differentiation under the integral, the derivatives of $p$ being bounded.

\emph{The knot relation.} The product gives $\widehat\mu(q\omega) = \varphi(\omega)\widehat\mu(\omega)$.
If $|\omega| < 2\pi/h - r$, then $\widehat w(\omega) = \varphi(\omega)$ by (1), and
$\widehat\mu(q\omega) = \widehat w(\omega)\widehat\mu(\omega)$. If $|\omega| \ge 2\pi/h - r$, then
$|\omega| \ge qr$, by the bound on $r$. Let $m \ge 1$ be the least integer with
$q^{-m}|\omega| < qr$; then $q^{-m}|\omega| \ge r$, as $q^{-m+1}|\omega| \ge qr$. So
$\varphi(q^{-m}\omega) = 0$ and $\widehat\mu(\omega) = 0$, and both sides of
$\widehat\mu(q\omega) = \widehat w(\omega)\widehat\mu(\omega)$ vanish, the left one being
$\varphi(\omega)\widehat\mu(\omega)$. The transforms of $D_q\mu$ and $\mu * w$ agree, so
$D_q\mu = \mu * w$.

\emph{The partial products of the mask.} If $|\omega|/q < 2\pi/h - r$, every $q^{-m}\omega$,
$m \ge 1$, lies in the ball where $\widehat w = \varphi$, and the partial products are those of
$\widehat\mu(\omega)$. Otherwise $|\omega| \ge q(2\pi/h - r) \ge qr$, and the $m \ge 1$ found above
has $q^{-m}|\omega| \in [r, qr)$, inside $[r, 2\pi/h - r)$ as $qr \le 2\pi/h - r$. There
$\widehat w(q^{-m}\omega) = \varphi(q^{-m}\omega) = 0$, so for $M \ge m$ both partial products
vanish, and $\widehat\mu(\omega) = 0$. In both cases the partial products of the mask converge to
$\widehat\mu(\omega)$.

(3) \emph{Not exact.} If a family with (a)--(d) carried the pyramid, its step $w = \mu_{1,q}$ would
be infinitely divisible, by Lemma~\ref{lem:steps-infinitely-divisible}, which (1) excludes.

\emph{The family.} For $0 < s \le t$ only finitely many integers $n$ have $s < q^n \le t$, so
$\mu_{s,t}$ is a finite convolution of probability laws, $\delta_0$ when there is none; in
particular $\mu_{t,t} = \delta_0$. For $0 < s \le t \le u$ the integers with $s < q^n \le t$ and
those with $t < q^n \le u$ make up those with $s < q^n \le u$, so
$\mu_{s,t} * \mu_{t,u} = \mu_{s,u}$. The knot relation, dilated by $D_{q^{n-1}}$, is
$D_{q^n}\mu = D_{q^{n-1}}\mu * D_{q^{n-1}}w$, and $t < q^n \le u$ exactly when
$N(t) < n \le N(u)$; so for $0 < t \le u$,
\[
  \mu_{0,u} = D_{q^{N(u)}}\mu
  = D_{q^{N(t)}}\mu * D_{q^{N(t)}}w * D_{q^{N(t)+1}}w * \dots * D_{q^{N(u)-1}}w
  = \mu_{0,t} * \mu_{t,u} .
\]
With $\mu_{0,0} = \delta_0$ the cases $t = 0$ hold too, and this is (a).
For (c): $s < q^n \le t$ exactly when $qs < q^{n+1} \le qt$, and
$D_qD_{q^{n-1}}w = D_{q^n}w$, so $\mu_{qs,qt} = D_q\mu_{s,t}$ for $s > 0$; and $N(qt) = N(t) + 1$
gives it for $s = 0 < t$, while $D_q\delta_0 = \delta_0$. For (d): the one integer with
$1 < q^n \le q$ is $n = 1$, so $\mu_{1,q} = D_{q^0}w = w$. Each kernel $\mu_t = D_{q^{N(t)}}\mu$,
$t > 0$, is a dilate of a rotation-invariant law, and a dilation commutes with every rotation, so
it is rotation invariant. Finally $N(t) = n - 1$ for $q^{n-1} \le t < q^n$. If
$D_{q^n}\mu = D_{q^{n-1}}\mu$, then $D_q\mu = \mu$, and by induction $\mu = D_{q^{-j}}\mu$ for every
$j \ge 1$, which tends weakly to $\delta_0$ as $j \to \infty$; but $\mu$ has a density. So the
kernels jump at every knot.

The cited facts are \ledger{A40}, in (1), and \ledger{A41}, through
Lemma~\ref{lem:steps-infinitely-divisible}, in (3).
\end{proof}

% draft: Remark 5.8(2), line 863.
The obstruction belongs to the grid. On the continuum, the isotropic \Matern{} pyramid
$\prod_{i \ge 1}(1 + c^2q^{-2i}t^2|\omega|^2)^{-1}$ in $\RR^d$ is $q^{\ZZ}$-covariant, indeed fully
covariant, and isotropic (Theorem~\ref{thm:matern-pyramid}). Its octave steps are
Bessel-potential laws, and its kernels are infinite convolutions of them. All are infinitely
divisible in every dimension. Its increments are infinitely
divisible too, and as functions of $u = |\omega|^2$ their exponents are Thorin--Bernstein functions
exactly when $t/s \in q^{\NN}$ (Proposition~\ref{prop:matern-increments}). No lattice, the
hexagonal included, carries it exactly. When the mask's support spans $\RR^d$, a lattice allows a
finite symmetry group, and for the frame \Matern{} pyramid that group is the lattice's own
(Proposition~\ref{prop:grid-pyramid-not-isotropic}(3)).

% draft: Remark 5.8(4), line 865.
% CHANGED: "non-separable" now carries its reason (blind review, flag 27): the Levy measure of a
% product law is carried by the axes, and this one is not.
The discrete Helmholtz mask, the five-point resolvent $\kappa_t$ on $h\ZZ^2$, is infinitely
divisible, and its grid pyramid is exact for every $q > 1$
(Proposition~\ref{prop:resolvent-masks-exact}). Its \Levy{} measure charges every lattice point off
the origin (Proposition~\ref{prop:resolvent-masks-exact}(4)), so its atoms leave the axes. The
mask is therefore not a product law. The factors of such a product would be its coordinate
marginals, the images of an infinitely divisible law under linear maps, and so infinitely
divisible. The \Levy{} measure of the product would then be carried by the axes, as in the proof of
Proposition~\ref{prop:grid-pyramid-not-isotropic}(3). This is a statement about the
ladder of the mask, in the sense of Definition~\ref{def:exact-grid-pyramid}. It is not the
subsampling of the discrete Helmholtz family on $h\ZZ^d$, a different object, which by
Proposition~\ref{prop:helmholtz-subsampling}(3) is not exact in $d \ge 2$, at any index, in the
weaker sense of being a multiple of a coarse member.

\section{The discrete Helmholtz family and its stencils}
\label{sec:discrete-helmholtz}

On a grid the Helmholtz solve keeps, in every dimension, the one property the cascade needs: its
kernel is positive. The inverse of the Helmholtz operator is a geometric mixture of the steps of a
random walk. That gives a family which is an exact cascade on one grid, with increments that are
mixtures of the identity and the kernel, as on the line. The section then compares three stencils
on the $3 \times 3$ neighbourhood of the plane: the five-point stencil, the nine-point stencil, and
the full Mehrstellen scheme, which also corrects the right-hand side. The two others buy rotation
covariance, and each gives up a property that the five-point stencil has.

In this section $D_h$, $D_5$, $D_9$ and $D^*$ are symbols and not the dilations $D_c$ of measures.
$D_h$ is a function of the phase $\vartheta$, and $D_5$, $D_9$, $D^*$ are functions of the frequency
$\omega$, with $\vartheta = h\omega$. Kernels on the grid $h\ZZ^d$ are indexed by $\ZZ^d$, and the
transform of a summable kernel $\lambda$ on $\ZZ^d$ is $\widehat\lambda(\vartheta) := \sum_n \lambda(n)
e^{i\langle n, \vartheta\rangle}$. Write $g(x) := 2 - 2\cos x$, and $\rho := h/t$ for a kernel of range
$t$, the ratio of Section~\ref{sec:price-measured} and Chapter~\ref{ch:computation}.

% draft: the proof of Proposition 5.9(3) (section 5.4), line 889, and its continuation in the proof of Proposition 5.10(3), line 915.
% CHANGED (second blind review 2026-11-12, flag 2; session decision): the Taylor bounds of the cosine, stated until now only inside the proof of prop:discrete-helmholtz(3) and continued inside the proof of prop:nine-point-family-rotation, as a lemma node of their own, so that every node that spends them cites a statement. Its proof is the chain of integrations moved from those two proofs.
\begin{lemma}[the Taylor bounds of the cosine]
\label{lem:cosine-bounds}
\lean{Pyramids.cosine_bounds}\leanok
For every real $x$, with $g(x) = 2 - 2\cos x$,
\[
  0 \le x^2 - g(x) \le \frac{x^4}{12}, \qquad
  0 \le g(x) - x^2 + \frac{x^4}{12} \le \frac{x^6}{360},
\]
\[
  -\frac{x^8}{20160} \le g(x) - x^2 + \frac{x^4}{12} - \frac{x^6}{360} \le 0 .
\]
% CHANGED (fidelity review, 2026-11-15, R26/R50/R69): the status line lists the declarations' hypotheses beside the node's (LINKAGE rule 6).
\statusT\quad\emph{(Cites nothing: eight integrations of $\cos \le 1$ from $0$. The bound of order
$6$ is the one inside Proposition~\ref{prop:discrete-helmholtz}(3), and the bound of order $8$ is
used by Proposition~\ref{prop:nine-point-family-rotation}. \textbf{Hypotheses} (the node's, each
with what its declarations require): $x$ real: same; \texttt{cosine\_bounds} requires nothing else.)}
\end{lemma}

% CHANGED (wave 1a, 2026-11-13): the first paragraph names the checked route, each step as a function vanishing at 0 with a nonnegative derivative, in place of the integration of an inequality; the chain of bounds is unchanged.
\begin{proof}
\leanok
Let $x \ge 0$. A differentiable function that vanishes at $0$ and has a nonnegative derivative on
$[0, x]$ is nonnegative at $x$. Write each bound below as the larger side minus the smaller side
being nonnegative. Each of these differences vanishes at $0$, and, since the derivative of $\sin$ is
$\cos$ and that of $\cos$ is $-\sin$, its derivative is the previous difference, the first one's
being $1 - \cos$. A bound for $\cos$ therefore gives one for $\sin$ of the same direction, and a
bound for $\sin$ one for $\cos$ of the opposite direction. Eight steps give in turn
\[
  \sin x \le x, \qquad \cos x \ge 1 - \frac{x^2}{2}, \qquad \sin x \ge x - \frac{x^3}{6}, \qquad
  \cos x \le 1 - \frac{x^2}{2} + \frac{x^4}{24},
\]
\[
  \sin x \le x - \frac{x^3}{6} + \frac{x^5}{120}, \qquad
  \cos x \ge 1 - \frac{x^2}{2} + \frac{x^4}{24} - \frac{x^6}{720},
\]
\[
  \sin x \ge x - \frac{x^3}{6} + \frac{x^5}{120} - \frac{x^7}{5040}, \qquad
  \cos x \le 1 - \frac{x^2}{2} + \frac{x^4}{24} - \frac{x^6}{720} + \frac{x^8}{40320} .
\]
Both sides of each bound for $\cos$ are even in $x$, so the four bounds for $\cos$ hold for every
real $x$. Multiply them by $-2$ and add $2$. The bounds $\cos x \ge 1 - x^2/2$ and
$\cos x \le 1 - x^2/2 + x^4/24$ give $x^2 - x^4/12 \le g(x) \le x^2$, which is the first bound. The
second of these and $\cos x \ge 1 - x^2/2 + x^4/24 - x^6/720$ give
$x^2 - x^4/12 \le g(x) \le x^2 - x^4/12 + x^6/360$, which is the second. The last two bounds for
$\cos$ give
\[
  x^2 - \frac{x^4}{12} + \frac{x^6}{360} - \frac{x^8}{20160} \le g(x)
  \le x^2 - \frac{x^4}{12} + \frac{x^6}{360} ,
\]
which is the third.
\end{proof}

% draft: Proposition 5.9 (section 5.4), lines 873-889 of draft/spatial-hemigroup-pyramids.md.
% Moved from the affine blueprint (06-paper-vi.tex), statement and proof machine-checked there and
% here (Formalization/Pyramids/DiscreteHelmholtz.lean).
% CHANGED: renamed c := t^2/h^2 to nu_h, and the family lambda_t, lambda_{s,t}, lambda_{r,t} to Lambda_t, Lambda_{s,t}, Lambda_{r,t} (BLUEPRINT-INPUTS).
% CHANGED: clause (4) names the two kernels it equates in the new symbols, as BLUEPRINT-INPUTS requires (the lambda_t row). The sentence "So at d = 1 and gamma = 1, Lambda_{t_h} = lambda_t" is that equation in the family's symbol; it follows from the clause at once and is kept, as the row requires it (second blind review, flag 10, answered without change).
% CHANGED (blind review 2026-11-12, flags 11, 12, 25): clause (2) writes the affine node's "r < s < t" as "0 < r < s < t", its reading, since Lambda_{r,s} is defined for 0 < r only (the end points s = 0, r = 0, which discrete_helmholtz_cascade also covers, are named in the annotation, not stated); "invariant under translations" read as commutation of the convolution with translations, as discrete_helmholtz_invariance states it; clause (4) states inline what prop:recursion-helmholtz(2) gives, which stands in chapter 6.
% CHANGED: widened, the author's approval of 2026-11-12: clause (2) states the end points, the increments for 0 <= s < t, Lambda_{0,t} = Lambda_t, and the composition for 0 <= r < s < t, which is exactly what Pyramids.discrete_helmholtz_cascade proves (second blind review, flag 1).
% CHANGED (fidelity review, 2026-11-15, R37): the tag adds
% discrete_helmholtz_increment_translation, which states that each increment Lambda_{s,t},
% 0 <= s < t, commutes with translations; discrete_helmholtz_invariance stated it for the members
% only. The statement is unchanged.
\begin{proposition}[the discrete Helmholtz family on $h\ZZ^d$]
\label{prop:discrete-helmholtz}
\uses{prop:recursion-helmholtz}
\lean{Pyramids.discrete_helmholtz_kernel, Pyramids.discrete_helmholtz_cascade,
Pyramids.discrete_helmholtz_invariance, Pyramids.discrete_helmholtz_symbol,
Pyramids.discrete_helmholtz_line, Pyramids.discrete_helmholtz_increment_translation}\leanok
Let $d \ge 1$, $h > 0$, and let $\Delta_h$ be the $(2d+1)$-point Laplacian on $h\ZZ^d$,
$\Delta_hu(n) := h^{-2}\sum_{i=1}^d\bigl(u(n+e_i) - 2u(n) + u(n-e_i)\bigr)$, with symbol $-D_h$,
$D_h(\vartheta) := h^{-2}\sum_i(2 - 2\cos\vartheta_i)$. Let $\sigma$ be the simple random walk step,
mass $1/2d$ at each $\pm e_i$. For $t > 0$ put $\nu_h := t^2/h^2$,
$p := 2d\nu_h/(1+2d\nu_h) \in (0,1)$, and
\[
  \kappa_t := (1-p)\sum_{k \ge 0}p^k\,\sigma^{*k} .
\]
\begin{enumerate}
% CHANGED (2026-11-23, review round 3 batch C, item 8, s6-7.R9): a strengthening. Uniqueness among
% bounded solutions, where it was stated among summable ones ("the unique summable solution"); a
% summable kernel is bounded, so the old form follows.
\item $\kappa_t$ is a probability measure on $\ZZ^d$ with
$\widehat{\kappa_t}(\vartheta) = (1 + t^2D_h(\vartheta))^{-1}$, and it is the unique bounded
solution of $(1 - t^2\Delta_h)\kappa_t = \delta_0$.
% CHANGED (2026-11-14, a strengthening): gamma >= 1 widened to gamma >= 0; neither declaration of
% the clause read it.
\item For an integer $\gamma \ge 0$ the \emph{discrete Helmholtz family}
$\Lambda_t := \kappa_t^{*\gamma}$, with transfer $(1 + t^2D_h)^{-\gamma}$, is an exact cascade: for
$0 \le s < t$ the increment $\Lambda_{s,t}$ with transfer $((1 + s^2D_h)/(1 + t^2D_h))^{\gamma}$ is the
$\gamma$-fold convolution power of the probability measure
$\tfrac{s^2}{t^2}\delta_0 + (1 - \tfrac{s^2}{t^2})\kappa_t$. Its end point is
$\Lambda_{0,t} = \Lambda_t$, and $\Lambda_{r,t} = \Lambda_{s,t} * \Lambda_{r,s}$ for
$0 \le r < s < t$. Convolution with each
member and each increment commutes with the translations of $\ZZ^d$, and the kernels are invariant
under the signed permutations of the axes.
\item \emph{(Rotation covariance fails at order $h^2$.)} For $\omega \in \RR^d$,
\[
  \Bigl|D_h(h\omega) - |\omega|^2 + \frac{h^2}{12}\sum_i\omega_i^4\Bigr|
  \le \frac{h^4}{360}\sum_i\omega_i^6 ,
\]
and for $d \ge 2$ the correction $\sum_i\omega_i^4$ is not a function of $|\omega|$.
\item At $d = 1$ the family is the sampled Laplace family. Let $a := e^{-h/t}$,
$\lambda_t(n) := \tfrac{1-a}{1+a}a^{|n|}$ the Laplace kernel sampled on $h\ZZ$ and renormalized,
and $t_h := h/(2\sinh(h/2t))$ the renormalized range, so that $\lambda_t$ is the unique summable
solution of $(1 - t_h^2\Delta_h)u = \delta_0$ (Proposition~\ref{prop:recursion-helmholtz}(2)). Then
$\kappa_{t_h} = \lambda_t$. So at $d = 1$ and $\gamma = 1$, $\Lambda_{t_h} = \lambda_t$.
\end{enumerate}
% CHANGED (fidelity review, 2026-11-15, R26/R50/R69): the status line lists the declarations' hypotheses beside the node's (LINKAGE rule 6).
\statusT\quad\emph{(Cites nothing: a geometric series, a maximum principle for bounded solutions,
and a Taylor bound. \textbf{The Lean}, with each declaration's hypotheses against the node's.
% CHANGED (2026-11-23, review round 3 batch C, item 8): the Lean states bounded uniqueness.
Clause (1) is \texttt{discrete\_helmholtz\_kernel}, under $d \ge 1$, $h > 0$, $t > 0$, as in the
node. It also proves $0 < p < 1$, and it states uniqueness among bounded solutions, the form of the
node since 2026-11-23 (review round 3, s6-7.R9, a strengthening; it was stated among summable
ones). Clause (2) is
% CHANGED (2026-11-14, a strengthening): the unused hypotheses dropped from the Lean.
\texttt{discrete\_helmholtz\_cascade}, under $d \ge 1$, $h > 0$, as in the node, for every integer
$\gamma \ge 0$ (its unused $\gamma \ge 1$ was dropped on 2026-11-14). It
states exactly the clause: the transfer of $\Lambda_t$ for $t > 0$; for $0 \le s < t$ that the
mixture and its $\gamma$-fold power are probability kernels, with the transfer stated;
$\Lambda_{0,t} = \Lambda_t$ for $t > 0$; and the composition for $0 \le r < s < t$. The increment is
the power of the mixture by its definition, \texttt{helmIncrement}. The translation and permutation sentence is
\texttt{discrete\_helmholtz\_invariance}. It requires nothing of $d$, $h$ and $\gamma$ (the three
hypotheses it carried unused were dropped on 2026-11-14), and it reads translations on summable
inputs. Clause (3) is \texttt{discrete\_helmholtz\_symbol}, under
$h > 0$ alone, so for every $d$. It says ``not a function of $|\omega|$'' as ``no $f$ has
$\sum_i\omega_i^4 = f(\sum_i\omega_i^2)$'', which is the same claim. Clause (4) is
\texttt{discrete\_helmholtz\_line}, under $h > 0$, $t > 0$, as in the node. Each rests on Lean core
alone. \textbf{Hypotheses} (the node's, each with what its declarations require): $d \ge 1$: same in
\texttt{discrete\_helmholtz\_kernel} and \texttt{discrete\_helmholtz\_cascade}, absent from
\texttt{discrete\_helmholtz\_invariance}, \texttt{discrete\_helmholtz\_increment\_translation} and
\texttt{discrete\_helmholtz\_symbol}, fixed at $1$ in \texttt{discrete\_helmholtz\_line}; $h > 0$:
same, absent from the two translation and invariance declarations; (1) $t > 0$: same; (2)
$\gamma \ge 0$ an integer: same, a natural number; $t > 0$ for the members: same; $0 \le s < t$
for the increments: same in \texttt{discrete\_helmholtz\_cascade} and
\texttt{discrete\_helmholtz\_increment\_translation}, stronger in
\texttt{discrete\_helmholtz\_invariance}, which requires $0 < s$ for the invariance of an increment
under the signed permutations, harmless, as the increment at $s = 0$ is the member
($\Lambda_{0,t} = \Lambda_t$, \texttt{discrete\_helmholtz\_cascade}), whose invariance it states;
$0 \le r < s < t$: same; the inputs of a translation: stronger for the members, which
\texttt{discrete\_helmholtz\_invariance} reads on summable inputs only, harmless, as
\texttt{discrete\_helmholtz\_increment\_translation} requires nothing of the input and covers the member
as the increment at $s = 0$; (3) $\omega \in \RR^d$: same; $d \ge 2$ for the second sentence: same;
(4) $d = 1$, $t > 0$: same.)}
\end{proposition}

% CHANGED (second blind review 2026-11-12, flag 2): the proof of clause (3) cites lem:cosine-bounds, where the chain of integrations now stands; the statement and the \lean tag are unchanged.
% CHANGED (fidelity review, 2026-11-15, R66): the edge to lem:cosine-bounds stays, as the bound is
% that lemma's; a sentence says that the declaration re-derives it in place.
\begin{proof}
\uses{prop:recursion-helmholtz, lem:cosine-bounds}
\leanok
\emph{Clause (1).} Convolution of summable kernels on $\ZZ^d$ is commutative, associative and
linear in each factor. It preserves probability kernels, and the Fourier series of a convolution is
the product of the Fourier series (Fubini for absolutely summable families). The step $\sigma$ is a
probability kernel with $\widehat\sigma(\vartheta) = d^{-1}\sum_i\cos\vartheta_i =: \phi(\vartheta)$, so
$\sigma^{*k}$ is a probability kernel with transform $\phi^k$. As $\sum_k\sum_np^k\sigma^{*k}(n) =
\sum_kp^k = (1-p)^{-1}$, $\kappa_t$ is a convex combination of the $\sigma^{*k}$ and a probability
kernel, and $\widehat{\kappa_t} = (1-p)\sum_kp^k\phi^k = (1-p)/(1-p\phi)$, using $|p\phi| < 1$. Now
$1 + t^2D_h = 1 + 2d\nu_h - 2d\nu_h\,\phi = (1 + 2d\nu_h)(1 - p\phi)$ and
$1 - p = (1+2d\nu_h)^{-1}$, which gives the transfer. Next
$(\kappa_t * \sigma)(n) = (1-p)\sum_kp^k\sigma^{*(k+1)}(n) = (\kappa_t(n) - (1-p)\delta_0(n))/p$,
and $\Delta_hu = h^{-2}(2d\,u * \sigma - 2d\,u)$, since
$(u * \sigma)(n) = (2d)^{-1}\sum_i(u(n-e_i) + u(n+e_i))$. With $2d\nu_h/p = 1 + 2d\nu_h$ this gives
$(1 - t^2\Delta_h)\kappa_t = (1 + 2d\nu_h)(1-p)\delta_0 = \delta_0$.
% CHANGED (2026-11-23, review round 3 batch C, item 8): the uniqueness is concluded for a bounded
% solution u, the difference with kappa_t being bounded as kappa_t is summable; the sentence "A
% summable kernel is bounded." closed the old proof.
For uniqueness let $u$ be a bounded solution. As $\kappa_t$ is summable, it is bounded, and so is
$e := u - \kappa_t$, with $(1 - t^2\Delta_h)e = 0$. Then
$(1 + 2d\nu_h)e(n) = \nu_h\sum_i(e(n+e_i) + e(n-e_i))$, so $|e(n)| \le p\,M$ with $M := \sup|e|$.
Hence $M \le pM$ and $M = 0$, that is $u = \kappa_t$.

\emph{Clause (2).} The transfers of $\Lambda_t = \kappa_t^{*\gamma}$ and of the increment follow
from the product rule for transforms, with $\tfrac{s^2}{t^2} + (1 - \tfrac{s^2}{t^2})(1 +
t^2D_h)^{-1} = (1 + s^2D_h)/(1 + t^2D_h)$. The mixture is a probability kernel since $0 \le s < t$,
and so is its power. At $s = 0$ the mixture is $\kappa_t$, so $\Lambda_{0,t} = \kappa_t^{*\gamma} =
\Lambda_t$. For $0 < s < t$ the resolvent identity
$\kappa_s * \kappa_t = (t^2\kappa_t - s^2\kappa_s)/(t^2 - s^2)$ holds. Indeed $\Delta_h$ passes
through the convolution, so both sides are bounded solutions of $(1 - t^2\Delta_h)u = \kappa_s$, the
right side because $s^2\Delta_h\kappa_s = \kappa_s - \delta_0$ and $t^2\Delta_h\kappa_t = \kappa_t -
\delta_0$. They agree by the uniqueness of clause (1). For $0 \le r < s < t$ the scale $s$ is
positive, so the identity applies to $\kappa_s * \kappa_t$, and expanding the convolution of the two
mixtures with it gives $\mu_{s,t} * \mu_{r,s} = \mu_{r,t}$, where
$\mu_{s,t} := \tfrac{s^2}{t^2}\delta_0 + (1 - \tfrac{s^2}{t^2})\kappa_t$. As convolution is
commutative and associative, $\mu_{s,t}^{*\gamma} * \mu_{r,s}^{*\gamma} = (\mu_{s,t} *
\mu_{r,s})^{*\gamma} = \mu_{r,t}^{*\gamma}$. Convolution with a kernel commutes with translations.
A signed permutation $\mathfrak{s}$ of the axes is an additive bijection of $\ZZ^d$ that permutes
$\{\pm e_i\}$, so $\sigma\circ\mathfrak{s} = \sigma$. Convolution of $\mathfrak{s}$-invariant kernels
is $\mathfrak{s}$-invariant (reindex the sum by $\mathfrak{s}$), hence so are $\sigma^{*k}$,
$\kappa_t$, the mixtures and their powers.

\emph{Clause (3).} By Lemma~\ref{lem:cosine-bounds},
$0 \le 2 - 2\cos x - x^2 + x^4/12 \le x^6/360$ for every real $x$. Apply it at $x = h\omega_i$, divide
by $h^2$ and sum. For $d \ge 2$, $\omega = e_1$ and $\omega = (e_1 + e_2)/\sqrt2$ have $|\omega| = 1$
and $\sum_i\omega_i^4$ equal to $1$ and $\tfrac12$.

\emph{Clause (4).} At $d = 1$, by clause (1) at the range $t_h > 0$, the sequence $\kappa_{t_h}$ on
$\ZZ$ is a summable solution of $(1 - t_h^2\Delta_h)u = \delta_0$, which by
Proposition~\ref{prop:recursion-helmholtz}(2) has the unique summable solution $\lambda_t$.

Every step is machine-checked, with no cited fact.
\end{proof}

% draft: the paragraph after Proposition 5.9, lines 891-895: prose (the naming, the M-matrix gloss,
% and the computation of a stage).
For $\gamma > d/2$ the continuum kernel with transform $(1 + t^2|\omega|^2)^{-\gamma}$ is the
\Matern{} covariance of smoothness $\gamma - d/2$ [\sscite{whittle1954stationary}, p.~448;
\sscite{lindgren2011explicit}, (2), p.~426], and $\Lambda_t$ is its grid counterpart. For $\gamma \le d/2$,
the plane's index $1$ among them, the kernel is unbounded at the origin, so it is no covariance
function. Clause (1) is the M-matrix property made concrete. A finite matrix with positive
diagonal, nonpositive off-diagonal entries and strict diagonal dominance has an entrywise
nonnegative inverse. The finite sections of $1 - t^2\Delta_h$ are such matrices, and clause (1)
writes the inverse on all of $\ZZ^d$ as a positive series.

The proposition also says how a stage is computed. Applying $\kappa_t$ to an array $f$ is one sparse
solve, $(1 - t^2\Delta_h)v = f$, by one of the methods of Section~\ref{sec:computation-solvers}:
transform division by $1 + t^2D_h$, a sparse factorization, or multigrid. The series that defines
$\kappa_t$ proves positivity and is not the method. Read as a method, the series is the Jacobi
iteration
$v \mapsto (1-p)f + p\,\sigma * v$, which contracts by $p$ per sweep, and $p$ is close to $1$ when
$t \gg h$. The increment of clause (2) from $s$ to $t$ is $\gamma$ repetitions of
$f \mapsto \tfrac{s^2}{t^2}f + (1 - \tfrac{s^2}{t^2})\kappa_t * f$, each one solve and one weighted
average. This is the line's stage with the forward--backward pair replaced by the solve. Under
transform division the $\gamma$ repetitions are one multiplication by
$\bigl((1 + s^2D_h)/(1 + t^2D_h)\bigr)^\gamma$. A level of the pyramid is simpler: $\gamma$ solves at
the octave's range, with no average, followed by keeping every $q$-th sample along each axis
(Section~\ref{sec:computation-decimation}). The stencils below change only the matrix of the solve, and under transform division only
its symbol. The full Mehrstellen scheme also smooths $f$ by a five-point average before the solve.

Clause (3) is the price of the five-point stencil. Its $h^2$ term $-\tfrac{h^2}{12}\sum_i\omega_i^4$
depends on direction. In the plane two other stencils on the $3 \times 3$ neighbourhood change that
term. The nine-point stencil gives up the principle on the axis neighbours $\mathcal{N}_4$
(Proposition~\ref{prop:grid-implicit-euler-principle-nine-axis}) and keeps it on the eight
neighbours. The full Mehrstellen scheme gives up the principle on every neighbourhood
(Proposition~\ref{prop:full-mehrstellen-no-principle}). The weights
$(-20, 4, 1)/6h^2$ of the nine-point stencil below are those of the Hermitian, or Mehrstellen,
formula of [\sscite{collatz1960numerical}, Ch.~V \S 2.5, p.~386], where they appear as the leading part
of a relation with correction terms.

% draft: the definitions before Proposition 5.10 (section 5.4), lines 897-900: D_5, Delta_9, D_9 and the first form of (5.5).
% CHANGED: a definition node of its own, so that the nodes on the nine-point stencil can \uses it; the draft defines these inline.
\begin{definition}[the five-point and nine-point symbols on $h\ZZ^2$]
\label{def:nine-point-laplacian}
\lean{Pyramids.gcos, Pyramids.fivePointSymbol, Pyramids.ninePointLap, Pyramids.ninePointSymbol, Pyramids.ninePointPhase}\leanok
\uses{prop:discrete-helmholtz}
Let $h > 0$. For $\omega \in \RR^2$ write $x := h\omega_1$, $y := h\omega_2$. The \emph{five-point
symbol} is $D_5(\omega) := D_h(h\omega) = h^{-2}\bigl(g(x) + g(y)\bigr)$, the symbol of
Proposition~\ref{prop:discrete-helmholtz} at $d = 2$. The \emph{nine-point Laplacian} on $h\ZZ^2$ is
\[
  \Delta_9u(n) := \frac{1}{6h^2}\Bigl(4\sum_{\pm e_i}u(n \pm e_i)
  + \sum_{\pm e_1 \pm e_2}u(n \pm e_1 \pm e_2) - 20u(n)\Bigr),
\]
the sums running over the four axis neighbours and the four diagonal neighbours. Its symbol is
$-D_9$, where
\[
  h^2D_9(\omega) := \tfrac{10}{3} - \tfrac43(\cos x + \cos y) - \tfrac23\cos x\cos y ,
\]
that is, $\Delta_9$ maps $n \mapsto e^{ih\langle n, \omega\rangle}$ to $-D_9(\omega)$ times itself.
Write $\Phi_9(\theta) := h^2D_9(\theta/h)$ for the same function of the phase $\theta = h\omega$; it
does not depend on $h$.
\statusT\quad\emph{(A definition. The symbol is read off the weights: the axis neighbours contribute
$\tfrac{4}{6}\cdot 2(\cos x + \cos y)$ and the diagonal ones $\tfrac16\cdot 4\cos x\cos y$.)}
\end{definition}

% draft: equation (5.5), line 900, the second and third forms.
% CHANGED: the identity (5.5) as a lemma of its own, with its proof; the draft states it without one.
\begin{lemma}[the nine-point symbol]
\label{lem:nine-point-symbol}
\lean{Pyramids.nine_point_symbol}\leanok
\uses{def:nine-point-laplacian}
With $x = h\omega_1$, $y = h\omega_2$, the symbol $h^2D_9$ has the three forms
\[
  h^2D_9 = \tfrac{10}{3} - \tfrac43(\cos x + \cos y) - \tfrac23\cos x\cos y ,
\]
\[
  h^2D_9 = \tfrac23\bigl(g(x) + g(y)\bigr) + \tfrac16\bigl(g(x+y) + g(x-y)\bigr) , \qquad
  h^2D_9 = h^2D_5 - \tfrac16g(x)g(y) .
\]
% CHANGED (fidelity review, 2026-11-15, R26/R50/R69): the status line lists the declarations' hypotheses beside the node's (LINKAGE rule 6).
\statusT\quad\emph{(Cites nothing: the addition formula for the cosine. \textbf{Hypotheses} (the
node's, each with what its declarations require): $h > 0$, from Definition~\ref{def:nine-point-laplacian}:
same in \texttt{nine\_point\_symbol}; $\omega \in \RR^2$: same.)}
\end{lemma}

\begin{proof}
\uses{def:nine-point-laplacian}
\leanok
By the addition formula, $\cos(x+y) + \cos(x-y) = 2\cos x\cos y$, so
$g(x+y) + g(x-y) = 4 - 4\cos x\cos y$. Hence
$\tfrac23(g(x) + g(y)) + \tfrac16(g(x+y) + g(x-y)) = \tfrac83 - \tfrac43(\cos x + \cos y) + \tfrac23
- \tfrac23\cos x\cos y$, which is the first form. For the third,
$\tfrac16g(x)g(y) = \tfrac16(4 - 4\cos x - 4\cos y + 4\cos x\cos y)$, and
$h^2D_5 - \tfrac16g(x)g(y) = 4 - \tfrac23 - (2 - \tfrac23)(\cos x + \cos y) - \tfrac23\cos x\cos y$,
which is the first form again.
\end{proof}

% draft: Proposition 5.10(1)-(2) (section 5.4), lines 902-906, proof line 913.
% CHANGED: split from Proposition 5.10: clauses (1)-(2) without the random-walk representation and the Levy measure (BLUEPRINT-INPUTS), which are prop:nine-point-family-levy.
% CHANGED: renamed nu := t^2/h^2 to nu_h and the family lambda_t to Lambda_t (BLUEPRINT-INPUTS).
% CHANGED: the transfer is read at the phase h omega, the readings the draft leaves to the reader.
% CHANGED (blind review 2026-11-12, flags 11, 25, 26, 32): the composition required for 0 < r < s < t, as the draft reads; translation invariance read as commutation; uniqueness among bounded solutions said to include the summable form of prop:discrete-helmholtz; the nine-point kernel, family and step written kappa^9_t, Lambda^9_t, sigma_9, apart from the five-point ones.
% CHANGED: widened, the author's approval of 2026-11-12: clause (2) states the end points, the increments for 0 <= s < t, Lambda^9_{0,t} = Lambda^9_t, and the composition for 0 <= r < s < t, which the proof already covered (second blind review, flag 1).
\begin{proposition}[the nine-point family on $h\ZZ^2$]
\label{prop:nine-point-family}
\lean{Pyramids.nine_point_family_1, Pyramids.nine_point_family_2, Pyramids.nine_point_family_2_invariance}\leanok
\uses{def:nine-point-laplacian, lem:nine-point-symbol}
Let $h > 0$ and let $\sigma_9$ be the step on $\ZZ^2$ with mass $1/5$ at each $\pm e_i$ and $1/20$ at
each $\pm e_1 \pm e_2$. For $t > 0$ put $\nu_h := t^2/h^2$,
$p := \tfrac{10}{3}\nu_h/(1 + \tfrac{10}{3}\nu_h) \in (0,1)$ and
$\kappa^9_t := (1-p)\sum_{k \ge 0}p^k\sigma_9^{*k}$.
\begin{enumerate}
\item $\kappa^9_t$ is a probability measure on $\ZZ^2$ with
$\widehat{\kappa^9_t}(h\omega) = (1 + t^2D_9(\omega))^{-1}$, and the unique bounded solution of
$(1 - t^2\Delta_9)\kappa^9_t = \delta_0$, in particular the unique summable one.
% CHANGED (2026-11-14, a strengthening): gamma >= 1 widened to gamma >= 0; neither declaration of
% the clause read it. The invariance declaration does not require h > 0 either; the node keeps h > 0,
% which clause (1) and the cascade read.
\item For an integer $\gamma \ge 0$ the family $\Lambda^9_t := (\kappa^9_t)^{*\gamma}$ is an exact
cascade: for $0 \le s < t$ the increment $\Lambda^9_{s,t}$ with transfer
$((1 + s^2D_9)/(1 + t^2D_9))^\gamma$ is the $\gamma$-fold power of the probability measure
$\tfrac{s^2}{t^2}\delta_0 + (1 - \tfrac{s^2}{t^2})\kappa^9_t$. Its end point is
$\Lambda^9_{0,t} = \Lambda^9_t$, and $\Lambda^9_{r,t} = \Lambda^9_{s,t} * \Lambda^9_{r,s}$ for
$0 \le r < s < t$. Convolution with each member
and each increment commutes with the translations of $\ZZ^2$, and the kernels are invariant under
the signed permutations.
\end{enumerate}
% CHANGED (fidelity review, 2026-11-15, R26/R50/R69): the status line lists the declarations' hypotheses beside the node's (LINKAGE rule 6).
\statusT\quad\emph{(Cites nothing: a geometric series, a maximum principle for bounded solutions,
and the resolvent identity. \textbf{Not claimed}: the representation as a walk at a geometric time
and the \Levy{} measure, which are Proposition~\ref{prop:nine-point-family-levy}.
\textbf{Hypotheses} (the node's, each with what its declarations require): $h > 0$: same in
\texttt{nine\_point\_family\_1} and \texttt{nine\_point\_family\_2}, absent from
\texttt{nine\_point\_family\_2\_invariance} (weaker); (1) $t > 0$: same; the competitors in the
uniqueness bounded, or summable: same, two conjuncts; (2) $\gamma \ge 0$ an integer: same, a
natural number; $t > 0$ for the members: same; $0 \le s < t$ for the increments: same;
$0 \le r < s < t$: same; the inputs of a translation: same, \texttt{nine\_point\_family\_2\_invariance}
requires nothing of them.)}
\end{proposition}

% CHANGED (wave 1e, 2026-11-13): the resolvent identity of clause (2) is checked through the factorization of the operator and the associativity of convolution, in place of "Delta_9 passes through convolution"; the rest of the proof is unchanged.
\begin{proof}
\uses{def:nine-point-laplacian, lem:nine-point-symbol}
\leanok
\emph{The step.} $\sigma_9$ is a probability kernel, invariant under the signed permutations, in
particular symmetric. At the phase $\theta = (x, y)$ its four axis points contribute
$\tfrac15\cdot2(\cos x + \cos y)$ and its four diagonal points $\tfrac1{20}\cdot4\cos x\cos y$, so
$\widehat\sigma_9(\theta) = \tfrac25(\cos x + \cos y) + \tfrac15\cos x\cos y$. Put $R := 10/3$, so that
$p = R\nu_h/(1 + R\nu_h)$ and $1 - p = (1 + R\nu_h)^{-1}$. By the first form of
Lemma~\ref{lem:nine-point-symbol},
\[
  1 + t^2D_9(\omega) = 1 + \nu_h\bigl(R - \tfrac43(\cos x + \cos y) - \tfrac23\cos x\cos y\bigr) ,
\]
and the right side is $(1 + R\nu_h)\bigl(1 - p\,\widehat\sigma_9(h\omega)\bigr)$. The same
factorization holds for the operator. For a bounded $u$,
\[
  R\nu_h(\sigma_9 * u)(n) = \frac{\nu_h}{6}\Bigl(4\sum_{\pm e_i}u(n \pm e_i)
  + \sum_{\pm e_1 \pm e_2}u(n \pm e_1 \pm e_2)\Bigr) ,
\]
and $20\nu_h/6 = R\nu_h$, so $(1 - t^2\Delta_9)u = (1 + R\nu_h)u - R\nu_h\,\sigma_9 * u$.

\emph{Clause (1).} As $|p\widehat\sigma_9| \le p < 1$, the series converges, and $\kappa^9_t$ is a convex
combination of the probability kernels $\sigma_9^{*k}$ with the weights $(1-p)p^k$, which sum to $1$.
So $\kappa^9_t$ is a probability kernel, and
$\widehat{\kappa^9_t} = (1-p)\sum_kp^k\widehat\sigma_9^k = (1-p)/(1 - p\widehat\sigma_9)$, which is
$(1 + t^2D_9)^{-1}$ by the factorization. Next
$\kappa^9_t * \sigma_9 = (1-p)\sum_kp^k\sigma_9^{*(k+1)} = (\kappa^9_t - (1-p)\delta_0)/p$. Hence, with
$R\nu_h/p = 1 + R\nu_h$,
\[
  (1 - t^2\Delta_9)\kappa^9_t = (1 + R\nu_h)\kappa^9_t - (1 + R\nu_h)\bigl(\kappa^9_t - (1-p)\delta_0\bigr)
  = (1 + R\nu_h)(1-p)\,\delta_0 = \delta_0 .
\]
For uniqueness, let $e$ be bounded with $(1 - t^2\Delta_9)e = 0$. Then
$(1 + R\nu_h)e = R\nu_h\,\sigma_9 * e$, so $|e(n)| \le p\sup|e|$ at every $n$, since $\sigma_9$ is a
probability kernel. Hence $\sup|e| \le p\sup|e|$, and $e = 0$. The difference of two bounded
solutions solves the homogeneous equation, so the bounded solution is unique.

\emph{Clause (2).} For $0 < s < t$ the resolvent identity
$\kappa^9_s * \kappa^9_t = (t^2\kappa^9_t - s^2\kappa^9_s)/(t^2 - s^2)$ holds. Indeed, by the
factorization at the time $t$, $(1 - t^2\Delta_9)u = (1 + Rt^2/h^2)u - (Rt^2/h^2)\,u * \sigma_9$ for
every $u$. Convolution of summable kernels is associative and linear in each argument, so
$(1 - t^2\Delta_9)(\kappa^9_s * \kappa^9_t) = \kappa^9_s * \bigl((1 - t^2\Delta_9)\kappa^9_t\bigr) =
\kappa^9_s * \delta_0 = \kappa^9_s$. On the other
side, $s^2\Delta_9\kappa^9_s = \kappa^9_s - \delta_0$ gives
$(1 - t^2\Delta_9)(s^2\kappa^9_s) = s^2\kappa^9_s - t^2(\kappa^9_s - \delta_0)$, so
$(1 - t^2\Delta_9)(t^2\kappa^9_t - s^2\kappa^9_s) = t^2\delta_0 - s^2\kappa^9_s + t^2\kappa^9_s - t^2\delta_0 =
(t^2 - s^2)\kappa^9_s$. Both sides are summable, hence bounded, and they agree by the uniqueness of
clause (1).

For $0 \le s < t$ put $\mu_{s,t} := \tfrac{s^2}{t^2}\delta_0 + (1 - \tfrac{s^2}{t^2})\kappa^9_t$, a
probability kernel, with $\mu_{0,t} = \kappa^9_t$ and transform
$\tfrac{s^2}{t^2} + (1 - \tfrac{s^2}{t^2})(1 + t^2D_9)^{-1} = (1 + s^2D_9)/(1 + t^2D_9)$. For
$0 \le r < s < t$, expanding $\mu_{s,t} * \mu_{r,s}$ and replacing $\kappa^9_t * \kappa^9_s$ by the
resolvent identity gives
\[
  \frac{r^2}{t^2}\delta_0 + \frac{s^2 - r^2}{t^2}\kappa^9_s + \frac{(t^2 - s^2)r^2}{t^2s^2}\kappa^9_t
  + \frac{s^2 - r^2}{t^2s^2}\bigl(t^2\kappa^9_t - s^2\kappa^9_s\bigr)
  = \frac{r^2}{t^2}\delta_0 + \Bigl(1 - \frac{r^2}{t^2}\Bigr)\kappa^9_t = \mu_{r,t} .
\]
For $0 \le s < t$ put $\Lambda^9_{s,t} := \mu_{s,t}^{*\gamma}$, a probability kernel with transfer
$((1 + s^2D_9)/(1 + t^2D_9))^\gamma$. As $\mu_{0,t} = \kappa^9_t$, $\Lambda^9_{0,t} = (\kappa^9_t)^{*\gamma}
= \Lambda^9_t$. As convolution is commutative and associative,
$\Lambda^9_{s,t} * \Lambda^9_{r,s} = (\mu_{s,t} * \mu_{r,s})^{*\gamma} = \Lambda^9_{r,t}$ for
$0 \le r < s < t$. Convolution with a kernel commutes with
translations. A signed permutation $\mathfrak{s}$ of the axes permutes the axis points and the
diagonal points among themselves, so $\sigma_9\circ\mathfrak{s} = \sigma_9$, and a convolution of
$\mathfrak{s}$-invariant kernels is $\mathfrak{s}$-invariant. Hence so are $\sigma_9^{*k}$, $\kappa^9_t$,
the mixtures and their powers.
\end{proof}
